%%
%% memoria.tex
%%
\documentclass[a4paper, 12pt]{book}

\pdfminorversion=4

\pagestyle{headings}

% Paquete para dibujos
\usepackage{tikz}

% Márgenes, hace que los márgenes laterales sean iguales
\usepackage[a4paper,top=3cm,right=3cm,left=3cm,bottom=3cm]{geometry}

% Referencias externas
\usepackage{xr-hyper}

% Paquete para generar enlaces clickables
\usepackage[pdftex, breaklinks=false, colorlinks=true, linkcolor=black, anchorcolor=black, urlcolor=blue, citecolor=red]{hyperref}

% Encoding y tratamiento de fuentes
\usepackage[T1]{fontenc}
\usepackage[spanish]{babel}
\usepackage[utf8]{inputenc}

% Paquete con el entorno tabularx, que define un tipo de celda X que hace word-wrapping
\usepackage{tabularx}

% Paquete con el entorno longtable y longtabu para tablas de varias páginas
\usepackage{longtable}

% Provee la flag H para forzar los floats en el sitio donde los pongas
\usepackage{float}

% Provee el comando \captionof para dar etiquetas a elementos no flotados
\usepackage{caption}

% Paquete que provee el comando \FloatBarrier para poner una barrera para floats
\usepackage{placeins}

% Paquete para controlar el interlineado a nivel de entorno
\usepackage{setspace}

% Paquete para la fuente charter
\usepackage{charter}

% Paquete para la fuente para códigos
\usepackage{inconsolata}

% Paquete para la fuente helvética
\usepackage[scaled=0.92]{helvet}

% Paquete para GNUPlot
\usepackage{gnuplottex}

% Guiones de hyphenado
\usepackage{hyphenat}

% Extra ToC listings
\usepackage{tocbibind}

% Gráficos
\usepackage{graphicx}

% Para poner figuras una al lado de otra
\usepackage[lofdepth,lotdepth]{subfig}

% Paquete para figuras giradas
\usepackage{rotating}

% Borra la indentación de párrafos y añade separación entre ellos
\usepackage{parskip}

% Paquete para quitar la numeración de las páginas en blanco
\usepackage{emptypage}

% Cálculos de márgenes
\usepackage{calc}

% Entornos para tablas y figuras giradas
\usepackage{rotating}

%%%%%%%%%%%%%%%%%%%%%%%%%%%%%%%%%%%%%%%%%%%%%%%%%%%%%%%%%%%%%%%%%%%%%%%%%%%%%%%%%%%%%%%%%%%%%%%%%%%%%%%%%%%%%%
% Paquetes en local (directorio 'paquetes')

% Automatiza el resaltado de sintaxis usando pygments
\usepackage{paquetes/minted}

\usemintedstyle[console]{vs}

\newminted[bashcode]{console}{} % tango vs
\newminted{python}{}


% Soporte de acrónimos
\usepackage[printonlyused]{paquetes/acronym-custom}

% Cambia el nombre del apéndice de bibliografía
\addto\captionsspanish{
\renewcommand\bibname{Bibliografía y referencias}
}

%%%%%%%%%%%%%%%%%%%%%%%%%%%%%%%%%%%%%%%%%%%%%%%%%%%%%%%%%%%%%%%%%%%%%%%%%%%%%%%%%%%%%%%%%%%%%%%%%%%%%%%%%%%%%%

\begin{document}

% Órdenes LaTeX variadas
% http://www.tex.ac.uk/cgi-bin/texfaq2html?label=chngmargonfly

\newenvironment{changemargin}[2]{%
  \begin{list}{}{%
    \setlength{\topsep}{0pt}%
    \setlength{\leftmargin}{#1}%
    \setlength{\rightmargin}{#2}%
    \setlength{\listparindent}{\parindent}%
    \setlength{\itemindent}{\parindent}%
    \setlength{\parsep}{\parskip}%
  }%
  \item[]}{\end{list}}

\newenvironment{nota}{
  \begin{changemargin}{2em}{2em}
    \textbf{\textsc{Nota: }}
}{
  \end{changemargin}
}

\newcommand{\nombreProyecto}{SiteUp: plataforma para la vigilancia de la disponibilidad de servicios de Internet}
\newcommand{\jugador}{\textit{Jugador}}
\newcommand{\sistema}{\textit{Sistema}}


% %% PRODUCTOS
% \newcommand{\nombrepostprocesador}{ACL2:\colonhyp{}Procesador}
% \newcommand{\nombrevisor}{XMLEye}
% \newcommand{\nombreyaxml}{YAXML:\colonhyp{}Reverse}
% \newcommand{\postprocesador}{\texttt{\nombrepostprocesador}\xspace}
% \newcommand{\visor}{\nombrevisor\xspace}
% \newcommand{\yaxml}{\texttt{\nombreyaxml}\xspace}

% \newcommand{\biblioteca}[1]{\index{#1}\texttt{#1}}

% % YAML/YAXML
% \newcommand{\etiqueta}[1]{{\bfseries \ttfamily #1}}

% % ACL2

% \newcommand{\orden}[1]{\texttt{#1}}   % Nombre de orden ACL2
% \newcommand{\fichero}[1]{\texttt{#1}} % Nombre de fichero
% \newcommand{\evento}[1]{\texttt{#1}}  % Nombre de un evento
% \newcommand{\libro}[1]{\textsc{#1}}   % Libro ACL2

% % PERL

% \newcommand{\modulo}[1]{\index{módulo Perl!#1}\index{#1|see{módulo Perl!#1}}\texttt{#1}}
% \newcommand{\funcion}[1]{\textit{#1}}

% % JAVA

% \newcommand{\clase}[1]{\textit{#1}}   % Clase Java
% \newcommand{\metodo}[1]{\texttt{#1}}  % Método (también Perl)
% \newcommand{\paquete}[1]{\texttt{#1}}

% %% MANUALES

% \newcommand{\accesoteclado}[1]{\textsc{#1}} % Acceso de teclado

% %% OTROS

% \newcommand{\patron}[1]{\emph{#1}}

% \newcolumntype{,}{>{$}r<{$}}
% \newcommand{\Index}[1]{#1\emph{\index{#1}}}

% \newcounter{pasoacept}

% \newenvironment{pruebaaceptacion}{
%   \setcounter{pasoacept}{0}
%   \begin{center}
%   \begin{tabular}{| >{\stepcounter{pasoacept}\arabic{pasoacept}. }p{.4\linewidth} | p{.5\linewidth}|}
%     \hline
%     \multicolumn{1}{| c |}{\textbf{Paso seguido}} & \multicolumn{1}{c|}{\textbf{Resultado esperado}} \\
%     \hline
%     \hline
% }{
%   \hline
%   \end{tabular}
%   \end{center}
% }

% \renewcommand{\lstlistlistingname}{Listados}
% \renewcommand{\lstlistingname}{Listado}

% \newcommand{\CPP}
% {\mbox{C\hspace{-.1em}\raise.2ex\hbox{+\hspace{-.1em}+}}\xspace}

% \setlength{\extrarowheight}{4pt}

% \lstset{
%   extendedchars,
%   flexiblecolumns,
%   stringstyle=\ttfamily,
%   showstringspaces=false,
%   frame=tb
% }

% %% PARTE DE DOCBOOK

% \newcommand{\application}[1]{\index{#1}\emph{#1}}
% \newcommand{\cmdsynopsis}[1]{\nohyphens{\texttt{#1}}}
% \newcommand*{\command}[1]{\nohyphens{\textbf{\texttt{#1}}}}
% \newcommand{\constant}[1]{\texttt{#1}}
% \newcommand{\computeroutput}[1]{#1}
% \newcommand*{\email}[1]{\nohyphens{\texttt{#1}}}
% \newcommand*{\envar}[1]{\nohyphens{\texttt{#1}}}
% \newcommand*{\filename}[1]{\texttt{#1}}
% \newcommand*{\guibutton}[2][]{\emph{#2}}
% \newcommand*{\guilabel}[2][]{\emph{#2}}
% \newcommand*{\guimenuitem}[2][]{\emph{#2}}
% \newcommand*{\guimenu}[2][]{\emph{#2}}
% \newcommand*{\keycombo}[1]{\textsc{#1}}
% \newcommand*{\keysym}[1]{\textsc{#1}}
% \newcommand*{\option}[1]{\nohyphens{\texttt{#1}}}
% \newcommand{\prompt}[1]{#1}
% \newcommand{\varname}[1]{\nohyphens{\texttt{#1}}}

% \newcommand{\note}[1]{\vskip 1em
%     \fbox{\parbox{\textwidth}{\textsc{Nota}
%     \vskip 1em #1}} \vskip 1em}


%%% Local Variables:
%%% mode: latex
%%% TeX-master: "memoria"
%%% End:


% Portada
\begin{titlepage}
  \centering
  \includegraphics[width=.3\textwidth]{0_misc/logo_uca}

  \bigskip
  \bigskip
  \bigskip

  \begin{changemargin}{2em}{2em}
    \centering

    {\Huge \textsc{\nohyphens{Escuela Superior de Ingeniería}}}

    \bigskip
    \bigskip
    \bigskip

    {\huge \nohyphens{Ingeniería Informática}}

    \bigskip
    \bigskip
    \bigskip
    \bigskip
    \bigskip
    \bigskip

    \begin{doublespace}
      {\LARGE \nohyphens{\nombreProyecto}}
    \end{doublespace}


    \bigskip
    \bigskip
    \bigskip
    \bigskip

    \bigskip
    \bigskip
    \bigskip
    \bigskip
    \bigskip
    \bigskip
    \bigskip

  \end{changemargin}

  {\Large José Tomás Tocino García \\}

  \bigskip

  {\large Cádiz, \today}

\end{titlepage}


%%% Local Variables: 
%%% mode: latex
%%% TeX-master: "../memoria"
%%% End: 

\cleardoublepage

% Segunda portada
{
  \thispagestyle{empty}
  \centering
  \includegraphics[width=.2\textwidth]{0_misc/logo_uca}

  \bigskip
  \bigskip
  \bigskip
  
  \begin{changemargin}{2em}{2em}

    \begin{center}
      {\Huge \textsc{\nohyphens{Escuela Superior de Ingeniería}}}
      
      \bigskip
      \bigskip
      
      {\huge \nohyphens{Ingeniería Informática}}
      
      \bigskip
      \bigskip
      \bigskip
      \bigskip
      
      \begin{doublespace}
        {\LARGE \nohyphens{\nombreProyecto}}
      \end{doublespace}


      \bigskip
      \bigskip
      \bigskip
      \bigskip
      
    \end{center}
  \end{changemargin}
  \begin{changemargin}{2em}{1em}
  \begin{flushleft}
    \Large

    \textsc{Departamento}: \nohyphens{Ingeniería Informática.} \\
    \textsc{Director del proyecto}: \nohyphens{Iván Ruiz Rube} \\
    \textsc{Autor del Proyecto}: \nohyphens{José Tomás Tocino García}. \\
  \end{flushleft}

  \end{changemargin}  

  \bigskip
  \bigskip
  \bigskip
  
  \begin{flushright}
    \large
    Cádiz, \today

    \bigskip
    \bigskip
    \bigskip
    \bigskip    
    \bigskip
    \bigskip

    Fdo.: José Tomás Tocino García
    
  \end{flushright}

}

%%% Local Variables: 
%%% mode: latex
%%% TeX-master: "../memoria"
%%% End: 



\cleardoublepage

% Texto de licencia (pequeño)
\bigskip
\bigskip

Este documento se halla bajo la licencia \ac{FDL}. Según estipula la
licencia, se muestra aquí el aviso de copyright. Se ha usado la
versión inglesa de la licencia, al ser la única reconocida
oficialmente por la \ac{FSF}.

\begin{quote}
  Copyright \copyright~\the\year~José Tomás Tocino García.

  Permission is granted to copy, distribute and/or modify this document
  under the terms of the GNU Free Documentation License, Version 1.2
  or any later version published by the Free Software Foundation;
  with no Invariant Sections, no Front-Cover Texts, and no Back-Cover Texts.
  A copy of the license is included in the section entitled "GNU
  Free Documentation License".
\end{quote}

%%% Local Variables: 
%%% mode: latex
%%% TeX-master: "../memoria"
%%% End: 

\cleardoublepage

% Agradecimientos
\section*{Agradecimientos}
Quiero agradecer y dedicar la presente memoria y, por extensión, el proyecto
SiteUp al completo:
% \begin{itemize}
% \item A Julian Raschke por crear y mantener Gosu.
% \item A Manuel Palomo Duarte y Rafael Rodríguez Galván, de la OSLUCA, por confiar en mí
%   como becario de la oficina, apoyarme y darme facilidades para desarrollar el
%   proyecto durante la beca.
% \item A Antonio García Domínguez, como co-tutor del proyecto, por sus acertados
%   consejos, su visión para descubrir cosas mejorables y su implicación con el
%   software y el conocimiento libre, gracias a la cual esta memoria fue posible.
% \item A mis compañeros más cercanos de la carrera, que me han apoyado y de los
%   que me he nutrido para desarrollar este proyecto.
% \item A mi familia y mi pareja, que finalmente no tendrán que recibir como
%   excusa el desarrollo del proyecto a la hora de dedicarles el tiempo que se
%   merecen.
% \item A Robert Penner por su magnífico conjunto de ecuaciones para animaciones,
%   sin las que oFlute tendría una interfaz mucho más pobre.
% \item A los que apuestan por el software libre y el conocimiento abierto. Sus
%   herramientas, bibliotecas y documentos han hecho posible que oFlute llegase a
%   buen puerto.
%\end{itemize}


\cleardoublepage

% Resumen
\documentclass[a4paper,12pt]{article}

\usepackage[left=2.5cm,top=2.5cm,right=2.5cm,bottom=2.5cm]{geometry}

\usepackage[T1]{fontenc}
\usepackage[spanish]{babel}
\usepackage[utf8]{inputenc}

\usepackage[pdftex, breaklinks=false, colorlinks=true, linkcolor=black, anchorcolor=black, urlcolor=blue, citecolor=red]{hyperref}

%\usepackage{framed}

% \usepackage{amsfonts}
% \usepackage{amssymb}
% \usepackage{amsthm}
% \usepackage{amsmath}
% \usepackage{array}
% \usepackage{caption}
% \usepackage{color}
% \usepackage[mirror]{crop} %<--------------LOL
% \usepackage{eurosym} %para el simbolo del euro
% \usepackage{fancybox}
% \usepackage[Rejne]{fncychap}% Sonny, Glenn, Lenny, Conny, Rejne, Bjarne
% \usepackage{float}
% \usepackage{fullpage}
% \usepackage{graphicx}
% \usepackage{multirow}
% \usepackage{pifont}
% \usepackage{stmaryrd}
% \usepackage{supertabular}
% \usepackage[dotinlabels]{titletoc}
% \usepackage[all]{xy}
% \usepackage{wasysym}




%\usepackage{tikz}
\usepackage{times}

%\usepackage{xxcolor}

%\usepackage{mathdesign}
%\usepackage{titlesec}


%\usepackage{estiloBase}
%\usepackage{colores}
%\usepackage{bera}
\usepackage{comandos}

% \usepackage{setspace}

% \usepackage{calc}

% \addto\captionsspanish{
% \renewcommand\bibname{Bibliografía y referencias}
% }

\def \titulo{SiteUp: plataforma para la vigilancia de la disponibilidad de servicios de Internet}
\def \autor{Alumno: José Tomás Tocino García\\Tutores: Iván Ruiz Rube}
\def \fecha{Mayo de 2014}

% Directorio de imágenes
%\graphicspath{{../img/}}

\usepackage{parskip}
\usepackage{abstract}

\begin{document}
\portada

\vspace{0.25cm}

\begin{abstract}
\textbf{SiteUp} es un proyecto para la vigilancia de la disponibilidad
  de servicios de Internet. Consta de una plataforma web, en la que los usuarios
  tienen la posibilidad dar de alta chequeos de varios tipos: envío de pings,
  chequeo de puertos, chequeo de registros DNS y comprobaciones mediante
  peticiones HTTP. Estos chequeos son ejecutados por el sistema de forma
  periódica y generan notificaciones cuando detectan fallos. Estas
  notificaciones se envían mediante correo electrónico o a través de SiteUp
  Client, una aplicación Android desarrollada a tal efecto. Los usuarios tienen
  también la posibilidad de revisar la información obtenida de los chequeos a
  través de la web a lo largo del tiempo. \\

  \textbf{Palabras clave:} Internet, Web, Servicio, SaaS, Vigilancia,
  Disponibilidad, Monitorización

\end{abstract}


\vspace{0.5cm}

\begin{center}
{\footnotesize Este documento se halla bajo la licencia FDL de GNU (Free Documentation
  License)\\ \url{http://www.gnu.org/licenses/fdl.html} }   
\end{center}



\tableofcontents

\section{Introducción}

\subsection{Contexto y motivación}

Las tecnologías de la información en general e Internet en particular son ya
parte integral de la sociedad. Casi todos los ámbitos de la vida, desde las
interacciones sociales hasta la búsqueda de empleo, cuentan ya con su reflejo en
las tecnologías de la información. Además, han surgido nuevos modelos
empresariales propios de Internet que han crecido a niveles comparables a los de
los negocios tradicionales. Empresas puramente digitales como Facebook o Twitter
ya cotizan en bolsa y realizan operaciones bursátiles del orden de miles de
millones de dólares~\cite{facebook-acquires-whatsapp}.

Se pone así de manifiesto la importancia de la fiablidad de los servicios e
infraestructuras de los que dependen estos nuevos modelos de negocio. La
disponibilidad debe ser siempre cercana al 100\%, dado que en caso contrario los
potenciales usuarios del servicio se encontrarán con que no pueden acceder a él,
dando lugar incluso a pérdidas económicas. Es el caso de Amazon, que llegó a
perder 4.8 millones de dólares al sufrir un fallo que dejó inaccesible su web
durante 40 minutos~\cite{amazon}.

\textbf{SiteUp} se modela como una herramienta de monitorización de servicios de
Internet accesible a través de la web. Los usuarios tendrán la posibilidad de
crear y gestionar una serie de \textit{chequeos} de diversos tipos sobre los
servicios web que elijan. La aplicación irá recopilando información relativa a
esos chequeos, e informará al usuario en caso de que las verificaciones que se
hayan dado de alta no coincidan con los resultados esperados.

Además, el usuario tendrá la posibilidad de recibir notificaciones de manera
instantánea a través del correo electrónico y de una aplicación para la
plataforma móvil \textbf{Android}. 


\subsection{Objetivos}
Los principales objetivos a alcanzar con \textbf{SiteUp} son los siguientes:

\begin{itemize}

\item Crear un conjunto de herramientas para la monitorización y el chequeo de
  diversos aspectos del estado de un servicio de Internet.
\item Crear una aplicación online, de acceso público, que permita la creación y
  gestión de chequeos de manera sencilla, basada internamente en las
  herramientas mencionadas en el punto anterior.
\item Habilitar a esta aplicación de un sistema de notificaciones mediante correo
  electrónico que alerte a los usuarios de posibles cambios en la disponibilidad
  de los servicios monitorizados.
\item Crear una aplicación móvil para que los usuarios tengan la opción de
  recibir notificaciones instantáneas provenientes de la aplicación web con
  información de sus chequeos.

\item Investigar y conocer los vectores de vigilancia usados habitualmente para
  monitorizar servicios de Internet.
\item Ampliar mis conocimientos sobre desarrollo web en general y las
  tecnologías de back-end en particular.
\item Adquirir soltura en el uso del lenguaje de programación Python en entornos
  web.
\item Obtener una base de conocimientos mínima sobre el desarrollo de
  aplicaciones sobre la plataforma móvil Android.
\item Utilizar un enfoque de análisis, diseño y codificación orientado
  a objetos, de una forma lo más clara y modular posible, para
  permitir ampliaciones y modificaciones sobre la aplicación por
  terceras personas.
\item Hacer uso de herramientas básicas en el desarrollo de software,
  como son los \textbf{sistemas de control de versiones} para llevar
  un control realista del desarrollo del software, así como hacer de
  las veces de sistema de copias de seguridad.

\end{itemize}

% \section{Planificación}
% El proyecto se ha desarrollado siguiendo un calendario basado en fases,
% utilizando un modelo de desarrollo iterativo incremental.

% \subsection{Primera iteración: conocimientos preliminares}
% En esta etapa se adquirieron los fundamentos teóricos para poder afrontar el
% desarrollo con todas las garantías. Se llevaron a cabo labores de documentación
% y aprendizaje autodidacta con las que se asentaron los conocimientos necesarios.

% \subsection{Segunda iteración: analizador básico}
% La segunda iteración se basó en el diseño de un analizador de notas básico, que
% sería el corazón del programa. Del éxito del desarrollo temprano del módulo que
% se encargaría del análisis de sonidos dependería la viabilidad completa del
% proyecto.

% \subsection{Tercera iteración: interfaz gráfica de usuario}
% En esta tercera iteración se propusieron numerosos diseños para la interfaz
% gráfica de usuario y se comenzó el desarrollo de los elementos de la interfaz,
% haciendo énfasis en conseguir un aspecto dinámico y jovial.

% \subsection{Cuarta iteración: motor de lecciones}
% En esta iteración se llevó a cabo el motor de lecciones, que presenta una serie
% de unidades didácticas en formato multimedia, compuestas de imágenes y
% textos. El sistema resultante es muy sencillo de ampliar y utilizar.

% \subsection{Quinta iteración: motor de canciones}
% Durante la quinta iteración se elaboró el sistema de canciones, encargado de
% listar y cargar las diferentes canciones, y puntuar al usuario según cómo
% interprete, mediante la flauta, las canciones que aparece en pantalla.. Es la
% parte de la aplicación con mayor interactividad.

% \subsection{Diagrama de Gantt}
% Se ha diseñado un diagrama de Gantt para reflejar la distribución de las tareas
% a lo largo del tiempo (figura~\ref{fig:gantt} en la página~\pageref{fig:gantt}).

% \section{Descripción general}
% \textbf{oFlute} se modela como una herramienta lúdico-educativa para alumnos que
% comiencen a aprender a usar la flauta dulce, proporcionando un entorno atractivo
% y ameno para el estudiante. Éstos tendrán la posibilidad de comprobar sus
% conocimientos sobre el uso de la flauta de forma totalmente práctica, gracias a
% un motor de análisis del sonido capaz de detectar las notas que emite el jugador
% con la flauta, capturadas por un micrófono, mediante el que la aplicación
% valorará la pericia del estudiante con la flauta.

% Además, los jugadores podrán recorrer una serie de pequeñas lecciones sobre
% música en general, y el uso de la flauta dulce en particular. Estas lecciones
% son totalmente ampliables, dando al usuario la posibilidad de crear las suyas
% propias.

% La aplicación cuenta con varias funcionalidades bien diferenciadas. A
% continuación se detallan cada una de ellas.

% \begin{figure}[htp!]
%   \vspace{1cm}
%   \centering
%   \includegraphics[width=0.8\textwidth]{imagen_menuPrincipal}
%   \caption{Pantalla del menú principal}
% \end{figure}

% \pagebreak

% \begin{figure}[htp!]
%   \centering
%   \includegraphics[width=0.95\textwidth]{imagen_diagrama_gantt}
%   \caption{Diagrama Gantt de iteraciones}
%   \label{fig:gantt}
% \end{figure}

% \pagebreak

% \subsection{Análisis de notas}
% Permite a los usuarios comprobar, de manera individual y pausada, que la
% interpretación de cada una de las notas en la flauta es correcta. Para ello, se
% presenta un pequeño analizador que responderá al sonido emitido por el usuario
% con la flauta mostrando la nota tocada en pantalla.

% \begin{figure}[h!]
%   \centering
%   \includegraphics[width=0.8\textwidth]{imagen_seccionAnalizador}
%   \caption{Pantalla del analizador de notas}
% \end{figure}

% \subsection{Motor de canciones}
% Mediante este sistema, el usuario tendrá la oportunidad de interpretar canciones
% completas a la vez que el computador analiza la eficacia del jugador,
% otorgándole una puntuación en tiempo real. Además, el motor de canciones es
% fácilmente expansible mediante ficheros de definición de canciones.

% \subsection{Motor de lecciones}
% Este sistema ofrece al usuario una serie de lecciones multimedia, con las que
% podrá aprender sobre diferentes aspectos de la música en general y la flauta en
% particular. Las lecciones cuentan con imágenes, texto y animaciones, que harán
% que el aprendizaje sea entretenido y ameno. También es posible añadir nuevas
% lecciones al sistema de forma sencilla.

% \begin{figure}[h!]
%   \centering
%   \includegraphics[width=0.8\textwidth]{imagen_seccionLecciones1}
%   \caption{Pantalla del menú de selección de lecciones}
% \end{figure}

% \subsection{Calibración del micrófono}
% \textbf{oFlute} ofrece la posibilidad de calibrar el micrófono, de forma que el
% sistema se adapte al ruido ambiental y el análisis del sonido capturado sea lo
% más exacto posible

% \section{Implementación}
% Durante el desarrollo del proyecto surgieron numerosas cuestiones y decisiones
% de implementación. A continuación se presenta una selección de las que más
% interés han generado. En la memoria del proyecto se desarrollan en mayor
% extensión.

% \subsection{Implementación del analizador básico}
% La creación del analizador básico de notas era crucial para el buen transcurso
% del resto del proyecto, por lo que fue uno de los objetivos que antes se
% abordaron. El desarrollo se dividió en dos partes. Por un lado, había que
% iniciar la captura de audio, lanzando el subsistema de sonido y empezando a
% capturar datos. Por otro lado, se tenía que hacer el análisis de los datos
% leídos para determinar qué nota se estaba tocando.

% La gestión del subsistema de audio se hizo mediante la \textit{Simple API} de
% PulseAudio, la biblioteca de audio que se empleó en oFlute. Se utilizó para la
% configuración y creación de un flujo de audio de entrada que permitiera captar,
% muestrear y samplear el sonido del micrófono, ofreciéndolo en forma de datos en
% punto flotante en un búffer.

% El siguiente paso fue el análisis del audio capturado. Para ello, utilizamos la
% Transformada Rápida de Fourier mediante su implementación en la biblioteca
% KissFFT. De forma iterativa, se iba analizando el búffer de sonido cada vez que
% éste se llenaba, dando lugar a una aproximación de la frecuencia fundamental
% capturada y, por ende, la nota que se estaba interpretando con la flauta.

% Una vez implementado y probado el sistema, se dió por bueno el analizador y se
% siguió construyendo el resto de la aplicación

% \subsection{Carga y uso de fuentes TrueType en Gosu}
% La biblioteca principal utilizada en la implementación del proyecto es Gosu, una
% librería de desarrollo de videojuegos 2D, multi-plataforma y con numerosas
% características. Desafortunadamente, no era posible utilizar fuentes TrueType
% cuando se utilizaba la biblioteca en GNU/Linux.

% Utilizando los conocimientos de otra popular biblioteca, SDL, se consiguió
% implementar un sistema para cargar y usar fuentes TrueType en Gosu de forma
% eficiente. Esta implementación finalmente acabó formando parte oficial de la
% biblioteca Gosu.

% \subsection{Animaciones dinámicas}
% Una de las decisiones iniciales de diseño fue la de hacer la interfaz gráfica de
% usuario lo más atractiva posible, intentando utilizar gráficos amigables y, en
% la medida de lo posible, animaciones y efectos dinámicos.

% Con esto, se tornaba necesario crear un sistema de animaciones lo más versátil
% posible, de forma que dotar a los elementos de la interfaz de movimiento fuera
% un proceso sencillo. 

% Para satisfacer este objetivo se ideó una serie de clases de animación que
% permiten animar las propiedades de los objetos de manera muy sencilla. Mediante
% el uso de las ecuaciones de animación de Robert Penner, el sistema creado cuenta
% con un gran número de opciones y diferentes formas de movimiento, según el tipo
% de función utilizada para calcular los valores: cuadrático, cúbico, etcétera.

% Gracias a este sistema, la interfaz de usuario de oFlute cuenta con un gran
% dinamismo y vistosidad.

% \subsection{Gestión de estados}
% La gestión de estados es uno de los retos principales a la hora de desarrollar
% un videojuego. Es necesario contar un sistema de gestión de estados robusto, que
% permita pasar de un estado a otro de forma sencilla, sin errores, y manteniendo
% información entre transiciones si fuera necesario.

% En el caso de oFlute, se creó un gestor de estados que permitió modularizar
% fácilmente el programa, haciendo una clara división entre secciones y
% facilitando la transición entre ellas. Además, mediante pruebas exhaustivas y el
% uso de herramientas de depuración, el gestor está implementado de forma que la
% carga y descarga de estados sea limpia y sin errores, evitando fugas de memoria.

% \subsection{Internacionalización}
% Durante el desarrolllo del proyecto surgió la necesidad de preparar el proyecto
% para su traducción. Inicialmente se optó por utilizar un sistema propio de
% traducción, basado en un script en Python y en una pequeña clase que generaba un
% diccionario según el idioma elegido. Sin embargo, esta opción resultó
% ineficiente y, sobre todo, alejada del resto de soluciones estándar a la hora de
% traducción de proyectos. 

% Por ello, se decidió investigar sobre las tecnologías de traducción más
% utilizadas en el panorama del software libre, y finalmente se optó por
% \textbf{GNU Gettext}~\cite{refgettext}. Para afianzar los conocimientos y
% facilitar el aprendizaje a otros desarrolladores, se editó un conciso manual
% sobre internacionalización de proyectos, presente como apéndice en la memoria del proyecto fin de



% \section{Conclusiones y difusión}
% Durante el transcurso del desarrollo de oFlute, y sobre todo al término del
% mismo, se han obtenido unas conclusiones y unos resultados, tanto de forma
% personal como para con la comunidad, que se resumen en esta sección y se tratan
% más en profundidad en el capítulo dedicado a tal efecto en la memoria principal.

% \subsection{Objetivos}
% Al término del desarrollo del proyecto, el proyecto ha completado todos los
% objetivos a cumplir, detallados en el planteamiento inicial. En particular:

% \begin{itemize}
% \item Se llevó a cabo la creación del módulo de análisis de sonido, consiguiendo
%   una efectividad en la detección de las notas muy alta.
% \item Se hizo uso del módulo en una sección de análisis de notas y en el propio
%   juego de interpretación de canciones, con resultados muy satisfactorios y
%   funcionamiento correcto.
% \item Además, el sistema de lecciones planteado se llevó a cabo correctamente,
%   consiguiendo un motor dinámico bastante potente y fácilmente ampliable, con
%   opción a añadir nuevas lecciones al sistema.
% \item Finalmente, todos los objetivos se llevaron a cabo respetando la premisa
%   de mantener una interfaz de usuario amigable y fluida, agradable a la vista y
%   sencilla de usar.
% \end{itemize}

% \subsection{Posibles mejoras}
% Hay un gran número de posibles mejoras y extensiones de la funcionalidad de
% oFlute, en unos casos descubiertas de forma personal y en otros casos aportadas
% por terceras personas. Algunas de las más interesantes son las siguientes:
% \begin{itemize}
% \item Extensión del sistema de lecciones para permitir lecciones de varios
%   pasos, así como integrar las lecciones con el analizador de audio para
%   comprobar los conocimientos in-situ.
% \item Mejora de la jugabilidad del sistema de canciones, añadiendo bonus y
%   puntos extra.
% \item Intentar extender el sistema a otros instrumentos o a la voz humana.
% \item Portar el videojuego a la plataforma Windows.
% \end{itemize}

% \subsection{Conclusiones personales}

% oFlute ha sido el proyecto más longevo al que me he enfrentado hasta ahora. A
% pesar de tener conciencia de la envergadura del mismo desde el principio, el
% tiempo para completarlo ha superado todas mis expectativas, sobre todo en lo que
% a documentación se refiere. 

% Gracias a oFlute he aprendido las técnicas básicas de la programación de audio,
% sobre todo en la parte técnica más que teórica. Es esta parte teórica, sobre
% todo la de análisis, la que me gustaría reforzar en un futuro, ahora que cuento
% con las bases para conseguirlo. Los videojuegos relacionados con el audio
% son un nicho aún poco explorado y que puede dar muchas satisfacciones, sobre
% todo cuando el producto se orienta a un público joven.

% Por otro lado, oFlute también me ha ayudado a aprender a usar bastantes
% tecnologías auxiliares, en algunos casos meras herramientas que han facilitado
% el trabajo y, en otros casos, elementos que resultaron ser de inestimable ayuda
% e imprescindible uso al término del proyecto.

% Una de estas tecnologías es \textbf{Boost}~\cite{boost}, un conjunto de
% bibliotecas para C++ que amplían en gran medida la biblioteca estándar del
% lenguaje. Un amplio número de componentes de Boost formarán parte del nuevo
% estándar C++0x, por lo que me ha servido para ponerme al día en las novedades
% que están por llegar.

% Al tratarse de la principal biblioteca utilizada durante el desarrollo, oFlute
% me ha provisto de un profundo conocimiento de \textbf{Gosu}~\cite{gosu},
% permitiéndome implementar videojuegos con mucha más fluidez y labrándome un
% pequeño \textit{framework} personal que utilizar de base en próximos
% proyectos. 

% Otra de las tecnologías que he aprendido ha sido \textbf{GNU Gettext}, que ha
% servido para internacionalizar el proyecto. Aunque inicialmente se pensó en
% utilizar una solución propia para la internacionalización el proyecto,
% posteriormente se decidió usar esta tecnología, ampliamente conocida y
% utilizada.

% \subsection{Difusión}

% \subsubsection{Conocimiento generado}

% Además de servir como recurso para el correcto transcurso del proyecto, los
% conocimientos adquiridos me permitieron, en numerosos casos, generar
% documentación adicional e impartir talleres relacionados las tecnologías
% previamente mencionadas.

% En cuanto a Boost, se llevó a cabo un taller durante los \textit{Cursos de
%   Verano de la OSLUCA}~\cite{cursosverano}, en el que se explicaron las partes
% más importantes de esta colección de bibliotecas, con numerosos ejemplos de
% muchas de ellas. Toda la documentación es libre~\cite{materialesCursoBoost}.

% Por otro lado, impartí un taller~\cite{tallergosu} sobre la biblioteca Gosu en
% colaboración con la ADVUCA~\cite{advuca}, cuya afluencia superó las 50
% personas. Los materiales pueden descargarse
% libremente~\cite{tallergosumateriales}.

% Por último, tras la aplicación de GNU Gettext en oFlute, escribí una guía
% concisa sobre traducción de proyectos con esta herramienta, que se puede
% encontrar en uno de los apéndices de la memoria, así como de forma online en la
% url \url{http://hdl.handle.net/10498/10772}.


% \subsubsection{Freegemas}

% Uno de los proyectos \textit{``hijos''} de oFlute ha sido
% \textbf{Freegemas}~\cite{freegemas}, un clon libre del popular juego tipo puzzle
% \textit{Bejeweled}. Freegemas es multiplataforma, funciona tanto en GNU/Linux
% como en Windows, y además forma parte oficial de Guadalinex~\cite{guadalinex},
% por lo que es posible encontrarlo en los repositorios oficiales.

% Además, Freegemas sirvió como base para una serie de tres artículos que
% publiqué, junto al director del presente PFC, en la revista Linux
% Magazine~\cite{linuxmagazine} sobre desarrollo de videojuegos en C++. Es posible
% encontrar estos artículos en el archivo de la
% revista~\cite{refarticulo1}~\cite{refarticulo2}~\cite{refarticulo3} bajo una
% licencia Creative Commons.

% \subsubsection{IV Concurso Universitario de Software Libre}

% Por otro lado, gracias a oFlute tuve la oportunidad de participar en el IV
% Concurso Universitario de Software Libre~\cite{cusl}. En el transcurso del
% concurso formé parte de una comunidad muy unida, en la que reinó el apoyo y la
% ayuda entre los concursantes. La final del concurso se celebró en la Escuela
% Superior de Ingeniería de Cádiz, en la que oFlute obtuvo una mención
% especial~\cite{cusl2}.

% Previa a la final nacional tuvo lugar la fase local del concurso, en el que el
% proyecto también recibió un accésit al mejor proyecto de
% innovación~\cite{cusllocal}.

% \bibliographystyle{hispa-annote}
% \bibliography{bibliografia}

\end{document}

\cleardoublepage

% Índices varios
\tableofcontents
\listoffigures

% Cambia el nombre del índice de tablas
\renewcommand{\listtablename}{Índice de tablas}
\renewcommand{\tablename}{Tabla}
\listoftables

%\setlength{\parskip}{1.2ex plus 0.4ex minus 0.1ex}

\chapter{Introducción}
\label{chap:introduccion}
En este capítulo se explicarán en detalle las instrucciones de uso de la
plataforma web.

\section{Introducción a SiteUp}

\textbf{SiteUp} es un servicio web creado para la monitorización de servicios de
Internet, especialmente pensado para dueños de sitios web y servidores que
deseen estar informados en caso de cualquier eventualidad.

El uso de internet en todos los ámbitos de la vida, desde las interacciones
sociales hasta la búsqueda de empleo, han dado lugar a la creación de servicios
accesibles a través de Internet y a nuevos modelos de negocio basados en
aquellos. Resulta por tanto imprescindible contar con sistemas para monitorizar
la disponibilidad de estos servicios y, en caso necesario, actuar de manera que
puedan subsanarse las causas de los problemas detectados.

SiteUp cubre precisamente estas necesidades, permitiendo a los usuarios:

\begin{itemize}
\item Crear chequeos mediante envío de paquetes ICMP (Ping), para revisar que un
  servidor está activo y conectado a la red.
\item Crear chequeos de registros DNS, de forma que el contenido de estos
  registros sea el correcto en todo momento.
\item Crear chequeos de puertos remotos, para vigilar que los servicios
  concretos dentro de un servidor son accesibles.
\item Crear chequeos mediante peticiones HTTP, para controlar que un servicio
  web responde de forma correcta y con el contenido adecuado.
\item Establecer las opciones de periodicidad de estos chequeos, de acuerdo a
  las necesidades particulares de cada caso.
\item Recibir notificaciones tanto a través de correo electrónico como mediante
  una aplicación móvil para el sistema operativo Android.
\end{itemize}

\section{Requisitos previos}

Para el uso correcto de la plataforma web SiteUp, el principal requisito es
contar con un dispositivo con acceso a internet que disponga de un navegador
reciente, como por ejemplo Google Chrome. Cualquier tipo de tablet, teléfono
móvil u ordenador personal podrá acceder a la plataforma web SiteUp, pero para
aprovechar mejor el potencial del sistema se recomienda usar una computadora
personal con las siguientes características:

\begin{itemize}
\item Navegador reciente. Se recomienda el uso de Google Chrome.
\item Pantalla con resolución mínima de 1280 píxeles en su lado longitudinal.
\item Conexión a internet con un caudal de descarga de al menos 512kb/s.
\end{itemize}

Si desea hacer uso de la aplicación móvil, será necesario que el usuario cuente
con un dispositivo con las siguientes características mínimas:

\begin{itemize}
\item Sistema operativo Android, versión mínima 2.3, versión recomendada 4.3.
\item Pantalla de al menos 3 pulgadas de diagonal y resolución mínima de 320x480
  píxeles.
\end{itemize}

\section{Uso de la plataforma web}

\subsection{Acceso}

A fecha de abril de 2014, la dirección de acceso a la plataforma web de SiteUp
es \url{http://siteup.josetomastocino.com}. Simplemente introduzca esa dirección
en su navegador y automáticamente será redirigido a la plataforma web, donde
verá la pantalla de bienvenida.

\subsection{Registro de usuario}

\begin{figure}[hbtp]
  \centering
  \includegraphics[width=0.3\textwidth]{apendice_manual_usuario/general_botones.png}
  \caption{Botones principales para el usuario anónimo}
  \label{fig:botones-1}
\end{figure}

Si es la primera vez que accede al sitio web deberá registrarse. Para ello, debe
pulsar en el botón \textbf{Sign up}, que se ve en la figura \ref{fig:botones-1}.

Tras ello, aparecerá el formulario de registro de usuario, que se ve en la
figura \ref{fig:registro}

\begin{figure}[hbtp]
  \centering
  \includegraphics[width=0.5\textwidth]{apendice_manual_usuario/formulario_registro.png}
  \caption{Formulario de registro de usuario}
  \label{fig:registro}
\end{figure}

En el formulario deberá introducir los siguientes datos:

\begin{itemize}
\item \textbf{Nombre de usuario}. Debe tener 30 caracteres como máximo, y solo
  pueden usarse letras, números y guiones.
\item \textbf{Dirección de correo electrónico}. Asegúrese de que es una
  dirección de correo correcta.
\item \textbf{Contraseña}.
\end{itemize}

Tras introducir los datos, deberá pulsar el botón \textbf{Sign up} del
formulario. Si el registro es correcto, el sistema dará de alta su usuario y
mostrará la pantalla \ref{fig:registro-correcto}. Si no, volverá a aparecer el
formulario con información sobre los errores en los datos.

\begin{figure}[hbtp]
  \centering
  \includegraphics[width=\textwidth]{apendice_manual_usuario/pantalla_register_completo.png}
  \caption{Mensaje indicando el registro correcto}
  \label{fig:registro-correcto}
\end{figure}

\subsection{Inicio de sesión}

Una vez que haya registrado su usuario, deberá pulsar en el botón \textbf{Log
  in}, como aparece en la imagen \ref{fig:botones-1}. Esto le llevará al
formulario de inicio de sesión que aparece en la figura \ref{fig:login}

\begin{figure}[hbtp]
  \centering
  \includegraphics[width=0.5\textwidth]{apendice_manual_usuario/general_formulario_login.png}
  \caption{Formulario de inicio de sesión}
  \label{fig:login}
\end{figure}

Deberá introducir su nombre de usuario y la contraseña que escribió al
registrarse, y posteriormente pulsar en el botón \textbf{Log in} del
formulario. Si los datos son correctos, será llevado al panel principal de la
web, donde podrá gestionar sus chequeos. Si no son correctos, volverá a aparecer
el formulario, con información sobre los errores encontrados.

\subsection{Recuperación de contraseña}

Si se registró como usuario pero no recuerda su contraseña es posible
recuperarla.

Para ello, vaya al formulario de inicio de sesión como ser ha explicado en el
paso anterior. En ese formulario, en la zona inferior, verá un enlace titulado
\textbf{Forgot password?} que deberá pulsar. Después de hacerlo se mostrará el
formulario de recuperación de contraseña, ilustrado en la figura
\ref{fig:formulario-reset}.

\begin{figure}[hbtp]
  \centering
  \includegraphics[width=0.5\textwidth]{apendice_manual_usuario/general_formulario_reset.png}
  \caption{Formulario de recuperación de contraseña}
  \label{fig:formulario-reset}
\end{figure}

En este formulario deberá introducir la dirección de correo electrónico que
utilizó al registrarse, y pulsar el botón \textbf{Send reset code}. Si lo hace
correctamente, aparecerá la pantalla de confirmación que se ve en la figura
\ref{fig:pantalla-reset-enviado}.

\begin{figure}[hbtp]
  \centering
  \includegraphics[width=0.5\textwidth]{apendice_manual_usuario/pantalla_reset_enviado.png}
  \caption{Pantalla de confirmación de envío de código de reinicio}
  \label{fig:pantalla-reset-enviado}
\end{figure}

El sistema habrá enviado un enlace a su dirección de correo electrónico. Deberá
pulsar el enlace que aparece en ese correo. Asegúrese de que recibe
correctamente el correo. Si no lo recibe, revise la carpeta de correo no
deseado. El enlace del correo le llevará a un formulario de reseteo de
contraseña como el que aparece en la figura \ref{fig:formulario-reset-2}.

\begin{figure}[hbtp]
  \centering
  \includegraphics[width=0.5\textwidth]{apendice_manual_usuario/general_formulario_reset_2.png}
  \caption{}
  \label{fig:formulario-reset-2}
\end{figure}

En este formulario deberá escribir la nueva contraseña de su cuenta dos veces,
para confirmar que la escribe correctamente. Tras ello, deberá pulsar el botón
\textbf{Set new password}, lo cual confirmará los cambios. 

Tras este proceso, podrá iniciar sesión como se ha descrito previamente, esta
vez utilizando la nueva contraseña.


\subsection{Pantalla de listado de chequeos}

\begin{figure}[hbtp]
  \centering
  \includegraphics[width=0.8\textwidth]{apendice_manual_usuario/pantalla-dashboard.png}
  \caption{Panel de control}
  \label{fig:dashboard}
\end{figure}

\begin{figure}[hbtp]
  \centering
  \includegraphics[width=0.4\textwidth]{apendice_manual_usuario/general_botones_2.png}
  \caption{Botones principales para el usuario logueado}
  \label{fig:botones-2}
\end{figure}

En la figura \ref{fig:dashboard} se puede ver el aspecto del panel de control de
la plataforma web SiteUp. Desde este panel podrá hacer una gran cantidad de
acciones, desde crear grupos de chequeos hasta editar sus detalles de usuario.

Es posible acceder a esta pantalla desde cualquier sitio de la web si está
logueado, pulsando en el botón \textbf{Dashboard} de la cabecera de la web, tal
y como se ve en la figura \ref{fig:botones-2}.

\subsection{Cerrar sesión}

Si ha iniciado sesión y desea cerrarla, debe dirigirse a los botones
principales, situados en el lado derecho de la cabecera de la web, tal y como
aparecen en la figura \ref{fig:botones-2}. Deberá pulsar el botón \textbf{Log
  out}, tras lo cual se cerrara su sesión y será redirigido a la pantalla
inicial de SiteUp.

\subsection{Edición de datos de usuario}

Si desea editar alguno de los datos de usuario que ingresó cuando hizo el
registro, deberá primero pulsar el botón \textbf{Your profile}, situado en la
cabecera de la web. Es el segundo de los botones que aparecen, tal y como se ve
en la figura \ref{fig:botones-2}.

Una vez hecho esto, aparecerá el formulario de edición de usuario, ilustrado en
la figura \ref{fig:formulario-usuario}, en el que verá sus actuales datos de
usuario. 

\begin{figure}[hbtp]
  \centering
  \includegraphics[width=0.8\textwidth]{apendice_manual_usuario/pantalla_perfil_usuario.png}
  \caption{}
  \label{fig:formulario-usuario}
\end{figure}

En este formulario podrá editar:

\begin{itemize}
\item Su \textbf{nombre} de usuario.
\item Su \textbf{dirección} de correo electrónico.
\item Si desea o no recibir \textbf{resúmenes diarios} sobre el estado de sus chequeos.
\end{itemize}

Una vez editados los datos que desee deberá pulsar el botón \textbf{Update
  details}. También podrá editar su contraseña pulsando el botón \textbf{Change
  password}, situado al pie del formulario a la izquierda.


\subsection{Creación de grupo de chequeos}

En SiteUp, los chequeos se organizan en \textbf{grupos de chequeos}, lo que
permite trabajar con ellos más fácilmente a la hora de activarlos, desactivarlos
o borrarlos. 

Para crear un grupo de chequeos deberá pulsar el botón \textbf{Add group},
situado en la zona superior derecha del panel principal. El botón puede verse en
la figura \ref{fig:crear-grupo}.

\begin{figure}[hbtp]
  \centering
  \includegraphics[width=0.3\textwidth]{apendice_manual_usuario/grupo_boton_nuevo.png}
  \caption{Botón para crear nuevo grupo de chequeos}
  \label{fig:crear-grupo}
\end{figure}

Tras pulsar el botón aparecerá el formulario de creación de nuevo grupo. En él
deberá introducir los datos del nuevo grupo. En particular:

\begin{itemize}
\item El \textbf{título} que identifica al grupo.
\end{itemize}

Tras ello, deberá pulsar el botón \textbf{Create group}, y el grupo será creado

\begin{figure}[H]
  \centering
  \includegraphics[width=0.8\textwidth]{apendice_manual_usuario/grupo_formulario_crear.png}
  \caption{Formulario de creación de grupo de chequeos}
  \label{fig:}
\end{figure}


\subsection{Operaciones con grupos de chequeos}

En la figura~\ref{fig:operaciones-grupo} se puede ver el menú de operaciones
disponibles para un grupo de chequeos. Este menú aparece en el panel principal,
al lado derecho junto al título de cada grupo de chequeos.

\begin{figure}[hbtp]
  \centering
  \includegraphics[width=0.5\textwidth]{apendice_manual_usuario/grupo_submenu_acciones.png}
  \caption{}
  \label{fig:operaciones-grupo}
\end{figure}

\subsubsection{Edición de grupo de chequeos}

En el menú de operaciones del grupo de chequeos, pulse el botón
\textbf{Edit}. Aparecerá un formulario similar al de creación de grupo, en el
que podrá editar el título del grupo de chequeos.

\subsubsection{Borrado de grupo de chequeos}

Para borrar un grupo de chequeos, pulse el botón \textbf{Delete} en el menú de
operaciones del grupo. Aparecerá un formulario como el de la figura
\ref{fig:formulario-borrado-grupo}, pidiéndole confirmación del borrado. 

\begin{figure}[hbtp]
  \centering
  \includegraphics[width=0.5\textwidth]{apendice_manual_usuario/formulario_borrado_grupo}
  \caption{Formulario de confirmación de borrado}
  \label{fig:formulario-borrado-grupo}
\end{figure}

Si pulsa en el botón \textbf{Go back}, volverá al panel principal, cancelando el
borrado. Si pulsa el botón \textbf{Delete}, confirmará la operación y el sistema
borrará el grupo de chequeos, así como todos los chequeos que contenga.

\subsubsection{Activación y desactivación de grupo de chequeos}

Es posible activar y desactivar de forma grupal todos los chequeos de un
grupo. Para ello, deberá dirigirse al menú de operaciones del grupo y pulsar en
\textbf{Enable all checks} o en \textbf{Disable all checks}, según sea su decisión.

\subsection{Creación de chequeo}

\begin{figure}[hbtp]
  \centering
  \includegraphics[width=0.8\textwidth]{apendice_manual_usuario/pantalla_seleccion_tipo.png}
  \caption{Pantalla de selección del tipo de chequeo a crear}
  \label{fig:seleccionar-tipo-check}
\end{figure}

Como se ha comentado previamente, el elemento principal de SiteUp es el
\textbf{chequeo}. El usuario da de alta chequeos dentro de un grupo, y estos
chequeos se lanzan periódicamente, generando datos.

Para crear un chequeo dentro de un grupo deberá ir al menú de operaciones del
grupo, que aparece en la figura \ref{fig:operaciones-grupo} y se sitúa a la
derecha del título de un grupo de chequeos. Dentro del menú, deberá pulsar la
opción \textbf{Add check}.



Tras esto aparecerá una pantalla en la que deberá seleccionar el tipo de chequeo
que desea crear. La pantalla es similar a la de la figura
\ref{fig:seleccionar-tipo-check}. Como se explicó en la introducción de este
manual, es posible crear cuatro tipo de chequeos: chequeos de ping, chequeos de
puerto, chequeos de registros DNS y chequeos por HTTP.

Si pulsa en el botón del tipo que desea crear, aparecerá un formulario con
campos para rellenar.

\subsubsection{Creación de un chequeo por ping}

\begin{figure}[hbtp]
  \centering
  \includegraphics[width=0.8\textwidth]{apendice_manual_usuario/pantalla_crear_ping.png}
  \caption{Formulario de creación de chequeo por ping}
  \label{fig:crear-ping}
\end{figure}

El formulario para crear un chequeo de tipo ping aparece en la
figura~\ref{fig:crear-ping}. Los campos que deberá rellenar son los siguientes:

\begin{itemize}
\item \textbf{Target}: es el objetivo a chequear. Debe ser un \textit{hostname} o una IP.
\item \textbf{Título} del chequeo.
\item \textbf{Descripción} opcional del chequeo.
\item \textbf{Check interval}: intervalo entre chequeos, definido en minutos. El
  valor mínimo es de 1 minuto.
\item \textbf{Notify email}: indica si los cambios de estado en este chequeo
  deberán ser notificados por correo electrónico.
\item \textbf{Notify Android}: indica si los cambios de estado en este chequeo
  deberán ser notificados mediante la aplicación Android.
\item \textbf{Should check timeout}: indica si hay que verificar que el ping
  enviado vuelve dentro de un tiempo máximo determinado.
\item \textbf{Timeout value}: si la anterior opción está activa, este campo
  define el número de milisegundos que puede tardar en responder el servidor al
  ping, como máximo.
\end{itemize}

\subsubsection{Creación de un chequeo de puerto}

\begin{figure}[hbtp]
  \centering
  \includegraphics[width=0.8\textwidth]{apendice_manual_usuario/pantalla_crear_port.png}
  \caption{Formulario de creación de chequeo de puerto}
  \label{fig:crear-port}
\end{figure}

El formulario para crear un chequeo de puertos aparece en la
figura~\ref{fig:crear-port}. Los campos que deberá rellenar son los siguientes:

\begin{itemize}
\item \textbf{Título} del chequeo.
\item \textbf{Descripción} opcional del chequeo.
\item \textbf{Target}: es el objetivo a chequear. Debe ser un \textit{hostname} o una IP.
\item \textbf{Target port}: es el puerto a chequear. Debe ser un entero.
\item \textbf{Check for string}: opcionalmente, se puede indicar una cadena que
  debe estar presente en la respuesta al conectar con el puerto.
\item \textbf{Check interval}: intervalo entre chequeos, definido en minutos. El
  valor mínimo es de 1 minuto.
\item \textbf{Notify email}: indica si los cambios de estado en este chequeo
  deberán ser notificados por correo electrónico.
\item \textbf{Notify Android}: indica si los cambios de estado en este chequeo
  deberán ser notificados mediante la aplicación Android.
\end{itemize}

\subsubsection{Creación de un chequeo de DNS}

\begin{figure}[hbtp]
  \centering
  \includegraphics[width=0.8\textwidth]{apendice_manual_usuario/pantalla_crear_dns.png}
  \caption{Formulario de creación de chequeo de registro DNS}
  \label{fig:crear-dns}
\end{figure}

El formulario para crear un chequeo de registro DNS aparece en la figura
\ref{fig:crear-dns}. Los campos que deberá rellenar son los siguientes:

\begin{itemize}
\item \textbf{Título} del chequeo.
\item \textbf{Descripción} opcional del chequeo.
\item \textbf{Target}: es el objetivo a chequear. Debe ser un \textit{hostname}.
\item \textbf{Record type}: indica el tipo de registro a chequear. Puede ser A, AAAA, CNAME, TXT y MX.
\item \textbf{Resolved address}: indica el contenido que debe tener el registro.
\item \textbf{Check interval}: intervalo entre chequeos, definido en minutos. El
  valor mínimo es de 1 minuto.
\item \textbf{Notify email}: indica si los cambios de estado en este chequeo
  deberán ser notificados por correo electrónico.
\item \textbf{Notify Android}: indica si los cambios de estado en este chequeo
  deberán ser notificados mediante la aplicación Android.
\end{itemize}

\subsubsection{Creación de un chequeo por HTTP}

\begin{figure}[hbtp]
  \centering
  \includegraphics[width=0.8\textwidth]{apendice_manual_usuario/pantalla_crear_http.png}
  \caption{Formulario de creación de chequeo por HTTP}
  \label{fig:crear-http}
\end{figure}

El formulario para crear un chequeo por HTTP aparece en la figura
\ref{fig:crear-http}. Los campos que deberá rellenar son los siguientes:

\begin{itemize}
\item \textbf{Título} del chequeo.
\item \textbf{Descripción} opcional del chequeo.
\item \textbf{Target}: es el objetivo a chequear. Debe ser una URL.
\item \textbf{Status code}: es el código de estado que deberá devolver la petición..
\item \textbf{Check for string}: opcionalmente, se puede indicar una cadena que
  debe estar presente en la respuesta a la petición.
\item \textbf{Check interval}: intervalo entre chequeos, definido en minutos. El
  valor mínimo es de 1 minuto.
\item \textbf{Notify email}: indica si los cambios de estado en este chequeo
  deberán ser notificados por correo electrónico.
\item \textbf{Notify Android}: indica si los cambios de estado en este chequeo
  deberán ser notificados mediante la aplicación Android.
\end{itemize}


\subsection{Operaciones con chequeos}

\begin{figure}[hbtp]
  \centering
  \includegraphics[width=\textwidth]{apendice_manual_usuario/chequeo_valido.png}
  \caption{}
  \label{fig:chequeo}
\end{figure}

Una vez creado un chequeo, aparecerá en el panel de control con el aspecto de la
figura~\ref{fig:chequeo}. Como se aprecia en la figura, se muestra mediante un
icono el estado del chequeo, los datos (título, tipo, descripción) y una gráfica
que ilustra el estado del chequeo en las últimas 24 horas.

\begin{figure}[hbtp]
  \centering
  \includegraphics[width=\textwidth]{apendice_manual_usuario/chequeo_invalido.png}
  \caption{}
  \label{fig:chequeo-invalido}
\end{figure}

Si un chequeo detecta un error, cambia de apariencia, pasando a tener el aspecto
que ofrece la figura \ref{fig:chequeo-invalido}. Además, según se haya
configurado en las opciones del chequeo, se envía una notificación al usuario
informando del cambio de estado del chequeo.

Cuando se ve en la lista general de chequeos, cada chequeo tiene una serie de
botones que permiten llevar a cabo operaciones diversas, que se detallan a
continuación.

\subsubsection{Activación y desactivación de un chequeo}

Para desactivar un chequeo, tiene que pulsar el botón \textbf{Disable} de la
botonera, que solo aparecerá si el chequeo está activado. De igual modo, si el
chequeo está desactivado, el botón que aparecerá tendrá la etiqueta
\textbf{Enable}, que podrá pulsar para activarlo.

\subsubsection{Borrado de un chequeo}

En la botonera se sitúa el botón \textbf{Delete}, que podrá pulsar para borrar
el chequeo. Aparecerá un formulario pidiéndole confirmación. Si confirma la
acción, se procederá al borrado del chequeo y todos los datos dependientes de
éste.

\subsubsection{Edición de un chequeo}

Si desea editar un chequeo debe pulsar el botón \textbf{Edit}, situado en la
botonera del chequeo. Aparecerá un formulario similar al de creación de chequeo,
pero con los datos actuales ya introducidos. Puede modificar los datos que desee
y confirmar los cambios pulsando el botón \textbf{Update check} del formulario.

\subsubsection{Vista de detalle del chequeo}

\begin{figure}[hbtp]
  \centering
  \includegraphics[width=\textwidth]{apendice_manual_usuario/pantalla_detalle_chequeo.png}
  \caption{}
  \label{fig:detalle-chequeo}
\end{figure}

En la botonera del chequeo, en el panel de control, verá un botón
\textbf{Details}. Si lo pulsa será llevado a la pantalla de vista de detalle del
chequeo, que se ilustra en la figura \ref{fig:detalle-chequeo}. Esta pantalla
ofrece mucha información sobre el chequeo. En particular:

\begin{itemize}
\item En la parte superior se sitúa la información general:
  \begin{itemize}
  \item Icono con el estado del chequeo.
  \item Título.
  \item Tipo de chequeo.
  \item Descripción.
  \item Botonera de acciones.
  \end{itemize}

\item En la parte media de la pantalla se sitúa un bloque con pestañas. Mediante
  estas pestañas es posible indicar sobre qué periodo de tiempo queremos
  visualizar los datos: últimas 24 horas o última semana. La información
  disponible es la siguiente:

  \begin{itemize}
  \item Porcentaje de \textit{uptime} del chequeo.
  \item \textit{(Solo chequeos tipo ping)} Tiempo de respuesta mínimo, máximo y promedio.
  \item Gráfica con el estado del chequeo.
  \end{itemize}

\item En la parte inferior de la pantalla se sitúa el panel de eventos del
  chequeo. Éste permite ver los últimos cambios de estado del chequeo:

  \begin{itemize}
  \item Tipo de cambio de estado (\textit{Up}, \textit{Down}).
  \item Duración del estado.
  \item Hora de comienzo del estado.
  \item Hora de fin del estado. Si es el estado en el que el chequeo se
    encuentra actualmente, no aparecerá ninguna fecha.
  \item Extra: muestra información adicional para estados de error.
  \end{itemize}
\end{itemize}

\begin{figure}[hbtp]
  \centering
  \includegraphics[width=\textwidth]{apendice_manual_usuario/pantalla_grafica_semanal.png}
  \caption{Detalle de la gráfica para un intervalo de tiempo de una semana}
  \label{fig:grafica-semanal}
\end{figure} 

\subsection{Recepción de notificaciones por correo electrónico}

\subsubsection{Notificaciones de cambio de estado}

Cuando el sistema detecte un cambio de estado en un chequeo, y esté activada la
notificación por correo electrónico, se enviará un email como el que aparece en
la figura \ref{fig:email-estado}.

\begin{figure}[hbtp]
  \centering
  \includegraphics[width=0.7\textwidth]{apendice_manual_usuario/email_estado.png}
  \caption{Email de notificación de cambio de estado}
  \label{fig:email-estado}
\end{figure}

En ese email se le indicará el título del chequeo, el nuevo estado, la fecha y
hora en la que ha pasado al nuevo estado, y un enlace directo al panel de
detalle del chequeo en SiteUp, de forma que pueda obtener más información.

\subsubsection{Notificaciones de resumen diarias}

Si ha activado el envío de resúmenes diarios sobre el estado de sus chequeos, el
sistema enviará todos los días un correo electrónico con información. El aspecto
del correo que recibirá se puede ver en la figura \ref{fig:email-resumen}

\begin{figure}[hbtp]
  \centering
  \includegraphics[width=0.7\textwidth]{apendice_manual_usuario/email_resumen.png}
  \caption{Email diario de resumen}
  \label{fig:email-resumen}
\end{figure}


\section{Uso de la aplicación móvil}

La aplicación móvil \textbf{SiteUp Client} es una aplicación para para el
sistema operativo Android que permite a sus usuarios revisar los chequeos que
tiene dados de alta en la plataforma web SiteUp, así como recibir de manera
instantánea notificaciones \textit{push} sobre el estado de sus chequeos.

\subsection{Instalación}

La aplicación móvil se instala copiando y ejecutando el fichero APK en el
dispositivo. En un futuro se podrá encontrar en la \textit{Play Store} de
Google.

\subsection{Lanzamiento}

En la figura~\ref{fig:manual-android-launcher} se puede observar el aspecto de
la aplicación \textbf{SiteUp Client} en el lanzador de Android. Para ejecutarla,
simplemente pulse sobre el icono de la aplicación.

\begin{figure}[htbp]
  \centering
  \includegraphics[width=0.4\textwidth]{5_diseno/android-1}
  \caption{Icono de SiteUp en el lanzador}
  \label{fig:manual-android-launcher}
\end{figure}

\subsection{Carga de la aplicación}

Tras lanzar la aplicación, el sistema cargará los datos necesarios y se
comunicará con los servidores de Google para obtener los datos necesarios. La
apariencia de la aplicación durante este proceso se puede ver en la figura
\ref{fig:manual-android-login}.

\begin{figure}[htbp]
  \centering
  \includegraphics[width=0.4\textwidth]{5_diseno/android-2}
  \caption{Pantalla de carga de la aplicación móvil}
  \label{fig:manual-android-loading}
\end{figure}

\subsection{Inicio de sesión}

Tras la carga, la primera vez que lance la aplicación será necesario que inicie
sesión con las credenciales que utilizó en la plataforma web. Para ello, tendrá
que introducir su nombre de usuario y contraseña en el formulario que aparece en
la figura~\ref{fig:manual-android-loading}, y pulsar el botón \textbf{Sign in}.

\begin{figure}[htbp]
  \centering
  \includegraphics[width=0.4\textwidth]{5_diseno/android-3}
  \caption{Pantalla de inicio de sesión de la aplicación móvil}
  \label{fig:manual-android-login}
\end{figure}

La aplicación móvil se comunicará con los servidores centrales de SiteUp,
validará sus credenciales de usuario y, si son correctas, le dará acceso al
sistema. Si las credenciales de usuario no son correctas, el formulario mostrará
un error.

\subsection{Revisión de chequeos}

Si ha iniciado sesión correctamente, tendrá la opción de revisar la lista de
chequeos que tiene dados de alta en la plataforma web \textbf{SiteUp}. La
apariencia de la lista de chequeos es igual a la de la
figura~\ref{fig:manual-android-dashboard}.

\begin{figure}[hbtp]
  \centering
  \includegraphics[width=0.4\textwidth]{5_diseno/android-4}
  \caption{Pantalla de chequeos de la aplicación móvil}
  \label{fig:manual-android-dashboard}
\end{figure}

Si pulsa en cualquiera de los chequeos, será redirigido a la plataforma web, que
podrá navegar sin problema gracias a que se encuentra preparada para su uso en
dispositivos móviles.

\subsection{Recepción de notificaciones}

La utilidad principal de la aplicación es la posibilidad de recibir
notificaciones en cualquier momento provenientes del servidor -- conocidas como
notificaciones \textit{push}. 

El procedimiento es transparente al usuario. Simplemente lance la aplicación e
inice sesión. Tras hacer esto, la aplicación estará preparada para recibir
notificaciones desde el servidor.

Cuando se reciba una notificación, la aplicación mostará un mensaje en el área
de notificaciones del móvil, con información resumida. Puede ver el aspecto de
la notificación en la figura \ref{fig:manual-android-notification}. Si pulsa en
la notificación será llevado a los detalles del chequeo para obtener más
detalles sobre el origen del aviso.

\begin{figure}[hbtp]
  \centering
  \includegraphics[width=0.4\textwidth]{5_diseno/android-5}
  \caption{Aspecto de las notificaciones}
  \label{fig:manual-android-notification}
\end{figure}



%%% Local Variables:
%%% mode: latex
%%% TeX-master: "../memoria"
%%% End:


\chapter{Planificación}
\label{chap:calendario}
En este capítulo se explicarán en detalle las instrucciones de uso de la
plataforma web.

\section{Introducción a SiteUp}

\textbf{SiteUp} es un servicio web creado para la monitorización de servicios de
Internet, especialmente pensado para dueños de sitios web y servidores que
deseen estar informados en caso de cualquier eventualidad.

El uso de internet en todos los ámbitos de la vida, desde las interacciones
sociales hasta la búsqueda de empleo, han dado lugar a la creación de servicios
accesibles a través de Internet y a nuevos modelos de negocio basados en
aquellos. Resulta por tanto imprescindible contar con sistemas para monitorizar
la disponibilidad de estos servicios y, en caso necesario, actuar de manera que
puedan subsanarse las causas de los problemas detectados.

SiteUp cubre precisamente estas necesidades, permitiendo a los usuarios:

\begin{itemize}
\item Crear chequeos mediante envío de paquetes ICMP (Ping), para revisar que un
  servidor está activo y conectado a la red.
\item Crear chequeos de registros DNS, de forma que el contenido de estos
  registros sea el correcto en todo momento.
\item Crear chequeos de puertos remotos, para vigilar que los servicios
  concretos dentro de un servidor son accesibles.
\item Crear chequeos mediante peticiones HTTP, para controlar que un servicio
  web responde de forma correcta y con el contenido adecuado.
\item Establecer las opciones de periodicidad de estos chequeos, de acuerdo a
  las necesidades particulares de cada caso.
\item Recibir notificaciones tanto a través de correo electrónico como mediante
  una aplicación móvil para el sistema operativo Android.
\end{itemize}

\section{Requisitos previos}

Para el uso correcto de la plataforma web SiteUp, el principal requisito es
contar con un dispositivo con acceso a internet que disponga de un navegador
reciente, como por ejemplo Google Chrome. Cualquier tipo de tablet, teléfono
móvil u ordenador personal podrá acceder a la plataforma web SiteUp, pero para
aprovechar mejor el potencial del sistema se recomienda usar una computadora
personal con las siguientes características:

\begin{itemize}
\item Navegador reciente. Se recomienda el uso de Google Chrome.
\item Pantalla con resolución mínima de 1280 píxeles en su lado longitudinal.
\item Conexión a internet con un caudal de descarga de al menos 512kb/s.
\end{itemize}

Si desea hacer uso de la aplicación móvil, será necesario que el usuario cuente
con un dispositivo con las siguientes características mínimas:

\begin{itemize}
\item Sistema operativo Android, versión mínima 2.3, versión recomendada 4.3.
\item Pantalla de al menos 3 pulgadas de diagonal y resolución mínima de 320x480
  píxeles.
\end{itemize}

\section{Uso de la plataforma web}

\subsection{Acceso}

A fecha de abril de 2014, la dirección de acceso a la plataforma web de SiteUp
es \url{http://siteup.josetomastocino.com}. Simplemente introduzca esa dirección
en su navegador y automáticamente será redirigido a la plataforma web, donde
verá la pantalla de bienvenida.

\subsection{Registro de usuario}

\begin{figure}[hbtp]
  \centering
  \includegraphics[width=0.3\textwidth]{apendice_manual_usuario/general_botones.png}
  \caption{Botones principales para el usuario anónimo}
  \label{fig:botones-1}
\end{figure}

Si es la primera vez que accede al sitio web deberá registrarse. Para ello, debe
pulsar en el botón \textbf{Sign up}, que se ve en la figura \ref{fig:botones-1}.

Tras ello, aparecerá el formulario de registro de usuario, que se ve en la
figura \ref{fig:registro}

\begin{figure}[hbtp]
  \centering
  \includegraphics[width=0.5\textwidth]{apendice_manual_usuario/formulario_registro.png}
  \caption{Formulario de registro de usuario}
  \label{fig:registro}
\end{figure}

En el formulario deberá introducir los siguientes datos:

\begin{itemize}
\item \textbf{Nombre de usuario}. Debe tener 30 caracteres como máximo, y solo
  pueden usarse letras, números y guiones.
\item \textbf{Dirección de correo electrónico}. Asegúrese de que es una
  dirección de correo correcta.
\item \textbf{Contraseña}.
\end{itemize}

Tras introducir los datos, deberá pulsar el botón \textbf{Sign up} del
formulario. Si el registro es correcto, el sistema dará de alta su usuario y
mostrará la pantalla \ref{fig:registro-correcto}. Si no, volverá a aparecer el
formulario con información sobre los errores en los datos.

\begin{figure}[hbtp]
  \centering
  \includegraphics[width=\textwidth]{apendice_manual_usuario/pantalla_register_completo.png}
  \caption{Mensaje indicando el registro correcto}
  \label{fig:registro-correcto}
\end{figure}

\subsection{Inicio de sesión}

Una vez que haya registrado su usuario, deberá pulsar en el botón \textbf{Log
  in}, como aparece en la imagen \ref{fig:botones-1}. Esto le llevará al
formulario de inicio de sesión que aparece en la figura \ref{fig:login}

\begin{figure}[hbtp]
  \centering
  \includegraphics[width=0.5\textwidth]{apendice_manual_usuario/general_formulario_login.png}
  \caption{Formulario de inicio de sesión}
  \label{fig:login}
\end{figure}

Deberá introducir su nombre de usuario y la contraseña que escribió al
registrarse, y posteriormente pulsar en el botón \textbf{Log in} del
formulario. Si los datos son correctos, será llevado al panel principal de la
web, donde podrá gestionar sus chequeos. Si no son correctos, volverá a aparecer
el formulario, con información sobre los errores encontrados.

\subsection{Recuperación de contraseña}

Si se registró como usuario pero no recuerda su contraseña es posible
recuperarla.

Para ello, vaya al formulario de inicio de sesión como ser ha explicado en el
paso anterior. En ese formulario, en la zona inferior, verá un enlace titulado
\textbf{Forgot password?} que deberá pulsar. Después de hacerlo se mostrará el
formulario de recuperación de contraseña, ilustrado en la figura
\ref{fig:formulario-reset}.

\begin{figure}[hbtp]
  \centering
  \includegraphics[width=0.5\textwidth]{apendice_manual_usuario/general_formulario_reset.png}
  \caption{Formulario de recuperación de contraseña}
  \label{fig:formulario-reset}
\end{figure}

En este formulario deberá introducir la dirección de correo electrónico que
utilizó al registrarse, y pulsar el botón \textbf{Send reset code}. Si lo hace
correctamente, aparecerá la pantalla de confirmación que se ve en la figura
\ref{fig:pantalla-reset-enviado}.

\begin{figure}[hbtp]
  \centering
  \includegraphics[width=0.5\textwidth]{apendice_manual_usuario/pantalla_reset_enviado.png}
  \caption{Pantalla de confirmación de envío de código de reinicio}
  \label{fig:pantalla-reset-enviado}
\end{figure}

El sistema habrá enviado un enlace a su dirección de correo electrónico. Deberá
pulsar el enlace que aparece en ese correo. Asegúrese de que recibe
correctamente el correo. Si no lo recibe, revise la carpeta de correo no
deseado. El enlace del correo le llevará a un formulario de reseteo de
contraseña como el que aparece en la figura \ref{fig:formulario-reset-2}.

\begin{figure}[hbtp]
  \centering
  \includegraphics[width=0.5\textwidth]{apendice_manual_usuario/general_formulario_reset_2.png}
  \caption{}
  \label{fig:formulario-reset-2}
\end{figure}

En este formulario deberá escribir la nueva contraseña de su cuenta dos veces,
para confirmar que la escribe correctamente. Tras ello, deberá pulsar el botón
\textbf{Set new password}, lo cual confirmará los cambios. 

Tras este proceso, podrá iniciar sesión como se ha descrito previamente, esta
vez utilizando la nueva contraseña.


\subsection{Pantalla de listado de chequeos}

\begin{figure}[hbtp]
  \centering
  \includegraphics[width=0.8\textwidth]{apendice_manual_usuario/pantalla-dashboard.png}
  \caption{Panel de control}
  \label{fig:dashboard}
\end{figure}

\begin{figure}[hbtp]
  \centering
  \includegraphics[width=0.4\textwidth]{apendice_manual_usuario/general_botones_2.png}
  \caption{Botones principales para el usuario logueado}
  \label{fig:botones-2}
\end{figure}

En la figura \ref{fig:dashboard} se puede ver el aspecto del panel de control de
la plataforma web SiteUp. Desde este panel podrá hacer una gran cantidad de
acciones, desde crear grupos de chequeos hasta editar sus detalles de usuario.

Es posible acceder a esta pantalla desde cualquier sitio de la web si está
logueado, pulsando en el botón \textbf{Dashboard} de la cabecera de la web, tal
y como se ve en la figura \ref{fig:botones-2}.

\subsection{Cerrar sesión}

Si ha iniciado sesión y desea cerrarla, debe dirigirse a los botones
principales, situados en el lado derecho de la cabecera de la web, tal y como
aparecen en la figura \ref{fig:botones-2}. Deberá pulsar el botón \textbf{Log
  out}, tras lo cual se cerrara su sesión y será redirigido a la pantalla
inicial de SiteUp.

\subsection{Edición de datos de usuario}

Si desea editar alguno de los datos de usuario que ingresó cuando hizo el
registro, deberá primero pulsar el botón \textbf{Your profile}, situado en la
cabecera de la web. Es el segundo de los botones que aparecen, tal y como se ve
en la figura \ref{fig:botones-2}.

Una vez hecho esto, aparecerá el formulario de edición de usuario, ilustrado en
la figura \ref{fig:formulario-usuario}, en el que verá sus actuales datos de
usuario. 

\begin{figure}[hbtp]
  \centering
  \includegraphics[width=0.8\textwidth]{apendice_manual_usuario/pantalla_perfil_usuario.png}
  \caption{}
  \label{fig:formulario-usuario}
\end{figure}

En este formulario podrá editar:

\begin{itemize}
\item Su \textbf{nombre} de usuario.
\item Su \textbf{dirección} de correo electrónico.
\item Si desea o no recibir \textbf{resúmenes diarios} sobre el estado de sus chequeos.
\end{itemize}

Una vez editados los datos que desee deberá pulsar el botón \textbf{Update
  details}. También podrá editar su contraseña pulsando el botón \textbf{Change
  password}, situado al pie del formulario a la izquierda.


\subsection{Creación de grupo de chequeos}

En SiteUp, los chequeos se organizan en \textbf{grupos de chequeos}, lo que
permite trabajar con ellos más fácilmente a la hora de activarlos, desactivarlos
o borrarlos. 

Para crear un grupo de chequeos deberá pulsar el botón \textbf{Add group},
situado en la zona superior derecha del panel principal. El botón puede verse en
la figura \ref{fig:crear-grupo}.

\begin{figure}[hbtp]
  \centering
  \includegraphics[width=0.3\textwidth]{apendice_manual_usuario/grupo_boton_nuevo.png}
  \caption{Botón para crear nuevo grupo de chequeos}
  \label{fig:crear-grupo}
\end{figure}

Tras pulsar el botón aparecerá el formulario de creación de nuevo grupo. En él
deberá introducir los datos del nuevo grupo. En particular:

\begin{itemize}
\item El \textbf{título} que identifica al grupo.
\end{itemize}

Tras ello, deberá pulsar el botón \textbf{Create group}, y el grupo será creado

\begin{figure}[H]
  \centering
  \includegraphics[width=0.8\textwidth]{apendice_manual_usuario/grupo_formulario_crear.png}
  \caption{Formulario de creación de grupo de chequeos}
  \label{fig:}
\end{figure}


\subsection{Operaciones con grupos de chequeos}

En la figura~\ref{fig:operaciones-grupo} se puede ver el menú de operaciones
disponibles para un grupo de chequeos. Este menú aparece en el panel principal,
al lado derecho junto al título de cada grupo de chequeos.

\begin{figure}[hbtp]
  \centering
  \includegraphics[width=0.5\textwidth]{apendice_manual_usuario/grupo_submenu_acciones.png}
  \caption{}
  \label{fig:operaciones-grupo}
\end{figure}

\subsubsection{Edición de grupo de chequeos}

En el menú de operaciones del grupo de chequeos, pulse el botón
\textbf{Edit}. Aparecerá un formulario similar al de creación de grupo, en el
que podrá editar el título del grupo de chequeos.

\subsubsection{Borrado de grupo de chequeos}

Para borrar un grupo de chequeos, pulse el botón \textbf{Delete} en el menú de
operaciones del grupo. Aparecerá un formulario como el de la figura
\ref{fig:formulario-borrado-grupo}, pidiéndole confirmación del borrado. 

\begin{figure}[hbtp]
  \centering
  \includegraphics[width=0.5\textwidth]{apendice_manual_usuario/formulario_borrado_grupo}
  \caption{Formulario de confirmación de borrado}
  \label{fig:formulario-borrado-grupo}
\end{figure}

Si pulsa en el botón \textbf{Go back}, volverá al panel principal, cancelando el
borrado. Si pulsa el botón \textbf{Delete}, confirmará la operación y el sistema
borrará el grupo de chequeos, así como todos los chequeos que contenga.

\subsubsection{Activación y desactivación de grupo de chequeos}

Es posible activar y desactivar de forma grupal todos los chequeos de un
grupo. Para ello, deberá dirigirse al menú de operaciones del grupo y pulsar en
\textbf{Enable all checks} o en \textbf{Disable all checks}, según sea su decisión.

\subsection{Creación de chequeo}

\begin{figure}[hbtp]
  \centering
  \includegraphics[width=0.8\textwidth]{apendice_manual_usuario/pantalla_seleccion_tipo.png}
  \caption{Pantalla de selección del tipo de chequeo a crear}
  \label{fig:seleccionar-tipo-check}
\end{figure}

Como se ha comentado previamente, el elemento principal de SiteUp es el
\textbf{chequeo}. El usuario da de alta chequeos dentro de un grupo, y estos
chequeos se lanzan periódicamente, generando datos.

Para crear un chequeo dentro de un grupo deberá ir al menú de operaciones del
grupo, que aparece en la figura \ref{fig:operaciones-grupo} y se sitúa a la
derecha del título de un grupo de chequeos. Dentro del menú, deberá pulsar la
opción \textbf{Add check}.



Tras esto aparecerá una pantalla en la que deberá seleccionar el tipo de chequeo
que desea crear. La pantalla es similar a la de la figura
\ref{fig:seleccionar-tipo-check}. Como se explicó en la introducción de este
manual, es posible crear cuatro tipo de chequeos: chequeos de ping, chequeos de
puerto, chequeos de registros DNS y chequeos por HTTP.

Si pulsa en el botón del tipo que desea crear, aparecerá un formulario con
campos para rellenar.

\subsubsection{Creación de un chequeo por ping}

\begin{figure}[hbtp]
  \centering
  \includegraphics[width=0.8\textwidth]{apendice_manual_usuario/pantalla_crear_ping.png}
  \caption{Formulario de creación de chequeo por ping}
  \label{fig:crear-ping}
\end{figure}

El formulario para crear un chequeo de tipo ping aparece en la
figura~\ref{fig:crear-ping}. Los campos que deberá rellenar son los siguientes:

\begin{itemize}
\item \textbf{Target}: es el objetivo a chequear. Debe ser un \textit{hostname} o una IP.
\item \textbf{Título} del chequeo.
\item \textbf{Descripción} opcional del chequeo.
\item \textbf{Check interval}: intervalo entre chequeos, definido en minutos. El
  valor mínimo es de 1 minuto.
\item \textbf{Notify email}: indica si los cambios de estado en este chequeo
  deberán ser notificados por correo electrónico.
\item \textbf{Notify Android}: indica si los cambios de estado en este chequeo
  deberán ser notificados mediante la aplicación Android.
\item \textbf{Should check timeout}: indica si hay que verificar que el ping
  enviado vuelve dentro de un tiempo máximo determinado.
\item \textbf{Timeout value}: si la anterior opción está activa, este campo
  define el número de milisegundos que puede tardar en responder el servidor al
  ping, como máximo.
\end{itemize}

\subsubsection{Creación de un chequeo de puerto}

\begin{figure}[hbtp]
  \centering
  \includegraphics[width=0.8\textwidth]{apendice_manual_usuario/pantalla_crear_port.png}
  \caption{Formulario de creación de chequeo de puerto}
  \label{fig:crear-port}
\end{figure}

El formulario para crear un chequeo de puertos aparece en la
figura~\ref{fig:crear-port}. Los campos que deberá rellenar son los siguientes:

\begin{itemize}
\item \textbf{Título} del chequeo.
\item \textbf{Descripción} opcional del chequeo.
\item \textbf{Target}: es el objetivo a chequear. Debe ser un \textit{hostname} o una IP.
\item \textbf{Target port}: es el puerto a chequear. Debe ser un entero.
\item \textbf{Check for string}: opcionalmente, se puede indicar una cadena que
  debe estar presente en la respuesta al conectar con el puerto.
\item \textbf{Check interval}: intervalo entre chequeos, definido en minutos. El
  valor mínimo es de 1 minuto.
\item \textbf{Notify email}: indica si los cambios de estado en este chequeo
  deberán ser notificados por correo electrónico.
\item \textbf{Notify Android}: indica si los cambios de estado en este chequeo
  deberán ser notificados mediante la aplicación Android.
\end{itemize}

\subsubsection{Creación de un chequeo de DNS}

\begin{figure}[hbtp]
  \centering
  \includegraphics[width=0.8\textwidth]{apendice_manual_usuario/pantalla_crear_dns.png}
  \caption{Formulario de creación de chequeo de registro DNS}
  \label{fig:crear-dns}
\end{figure}

El formulario para crear un chequeo de registro DNS aparece en la figura
\ref{fig:crear-dns}. Los campos que deberá rellenar son los siguientes:

\begin{itemize}
\item \textbf{Título} del chequeo.
\item \textbf{Descripción} opcional del chequeo.
\item \textbf{Target}: es el objetivo a chequear. Debe ser un \textit{hostname}.
\item \textbf{Record type}: indica el tipo de registro a chequear. Puede ser A, AAAA, CNAME, TXT y MX.
\item \textbf{Resolved address}: indica el contenido que debe tener el registro.
\item \textbf{Check interval}: intervalo entre chequeos, definido en minutos. El
  valor mínimo es de 1 minuto.
\item \textbf{Notify email}: indica si los cambios de estado en este chequeo
  deberán ser notificados por correo electrónico.
\item \textbf{Notify Android}: indica si los cambios de estado en este chequeo
  deberán ser notificados mediante la aplicación Android.
\end{itemize}

\subsubsection{Creación de un chequeo por HTTP}

\begin{figure}[hbtp]
  \centering
  \includegraphics[width=0.8\textwidth]{apendice_manual_usuario/pantalla_crear_http.png}
  \caption{Formulario de creación de chequeo por HTTP}
  \label{fig:crear-http}
\end{figure}

El formulario para crear un chequeo por HTTP aparece en la figura
\ref{fig:crear-http}. Los campos que deberá rellenar son los siguientes:

\begin{itemize}
\item \textbf{Título} del chequeo.
\item \textbf{Descripción} opcional del chequeo.
\item \textbf{Target}: es el objetivo a chequear. Debe ser una URL.
\item \textbf{Status code}: es el código de estado que deberá devolver la petición..
\item \textbf{Check for string}: opcionalmente, se puede indicar una cadena que
  debe estar presente en la respuesta a la petición.
\item \textbf{Check interval}: intervalo entre chequeos, definido en minutos. El
  valor mínimo es de 1 minuto.
\item \textbf{Notify email}: indica si los cambios de estado en este chequeo
  deberán ser notificados por correo electrónico.
\item \textbf{Notify Android}: indica si los cambios de estado en este chequeo
  deberán ser notificados mediante la aplicación Android.
\end{itemize}


\subsection{Operaciones con chequeos}

\begin{figure}[hbtp]
  \centering
  \includegraphics[width=\textwidth]{apendice_manual_usuario/chequeo_valido.png}
  \caption{}
  \label{fig:chequeo}
\end{figure}

Una vez creado un chequeo, aparecerá en el panel de control con el aspecto de la
figura~\ref{fig:chequeo}. Como se aprecia en la figura, se muestra mediante un
icono el estado del chequeo, los datos (título, tipo, descripción) y una gráfica
que ilustra el estado del chequeo en las últimas 24 horas.

\begin{figure}[hbtp]
  \centering
  \includegraphics[width=\textwidth]{apendice_manual_usuario/chequeo_invalido.png}
  \caption{}
  \label{fig:chequeo-invalido}
\end{figure}

Si un chequeo detecta un error, cambia de apariencia, pasando a tener el aspecto
que ofrece la figura \ref{fig:chequeo-invalido}. Además, según se haya
configurado en las opciones del chequeo, se envía una notificación al usuario
informando del cambio de estado del chequeo.

Cuando se ve en la lista general de chequeos, cada chequeo tiene una serie de
botones que permiten llevar a cabo operaciones diversas, que se detallan a
continuación.

\subsubsection{Activación y desactivación de un chequeo}

Para desactivar un chequeo, tiene que pulsar el botón \textbf{Disable} de la
botonera, que solo aparecerá si el chequeo está activado. De igual modo, si el
chequeo está desactivado, el botón que aparecerá tendrá la etiqueta
\textbf{Enable}, que podrá pulsar para activarlo.

\subsubsection{Borrado de un chequeo}

En la botonera se sitúa el botón \textbf{Delete}, que podrá pulsar para borrar
el chequeo. Aparecerá un formulario pidiéndole confirmación. Si confirma la
acción, se procederá al borrado del chequeo y todos los datos dependientes de
éste.

\subsubsection{Edición de un chequeo}

Si desea editar un chequeo debe pulsar el botón \textbf{Edit}, situado en la
botonera del chequeo. Aparecerá un formulario similar al de creación de chequeo,
pero con los datos actuales ya introducidos. Puede modificar los datos que desee
y confirmar los cambios pulsando el botón \textbf{Update check} del formulario.

\subsubsection{Vista de detalle del chequeo}

\begin{figure}[hbtp]
  \centering
  \includegraphics[width=\textwidth]{apendice_manual_usuario/pantalla_detalle_chequeo.png}
  \caption{}
  \label{fig:detalle-chequeo}
\end{figure}

En la botonera del chequeo, en el panel de control, verá un botón
\textbf{Details}. Si lo pulsa será llevado a la pantalla de vista de detalle del
chequeo, que se ilustra en la figura \ref{fig:detalle-chequeo}. Esta pantalla
ofrece mucha información sobre el chequeo. En particular:

\begin{itemize}
\item En la parte superior se sitúa la información general:
  \begin{itemize}
  \item Icono con el estado del chequeo.
  \item Título.
  \item Tipo de chequeo.
  \item Descripción.
  \item Botonera de acciones.
  \end{itemize}

\item En la parte media de la pantalla se sitúa un bloque con pestañas. Mediante
  estas pestañas es posible indicar sobre qué periodo de tiempo queremos
  visualizar los datos: últimas 24 horas o última semana. La información
  disponible es la siguiente:

  \begin{itemize}
  \item Porcentaje de \textit{uptime} del chequeo.
  \item \textit{(Solo chequeos tipo ping)} Tiempo de respuesta mínimo, máximo y promedio.
  \item Gráfica con el estado del chequeo.
  \end{itemize}

\item En la parte inferior de la pantalla se sitúa el panel de eventos del
  chequeo. Éste permite ver los últimos cambios de estado del chequeo:

  \begin{itemize}
  \item Tipo de cambio de estado (\textit{Up}, \textit{Down}).
  \item Duración del estado.
  \item Hora de comienzo del estado.
  \item Hora de fin del estado. Si es el estado en el que el chequeo se
    encuentra actualmente, no aparecerá ninguna fecha.
  \item Extra: muestra información adicional para estados de error.
  \end{itemize}
\end{itemize}

\begin{figure}[hbtp]
  \centering
  \includegraphics[width=\textwidth]{apendice_manual_usuario/pantalla_grafica_semanal.png}
  \caption{Detalle de la gráfica para un intervalo de tiempo de una semana}
  \label{fig:grafica-semanal}
\end{figure} 

\subsection{Recepción de notificaciones por correo electrónico}

\subsubsection{Notificaciones de cambio de estado}

Cuando el sistema detecte un cambio de estado en un chequeo, y esté activada la
notificación por correo electrónico, se enviará un email como el que aparece en
la figura \ref{fig:email-estado}.

\begin{figure}[hbtp]
  \centering
  \includegraphics[width=0.7\textwidth]{apendice_manual_usuario/email_estado.png}
  \caption{Email de notificación de cambio de estado}
  \label{fig:email-estado}
\end{figure}

En ese email se le indicará el título del chequeo, el nuevo estado, la fecha y
hora en la que ha pasado al nuevo estado, y un enlace directo al panel de
detalle del chequeo en SiteUp, de forma que pueda obtener más información.

\subsubsection{Notificaciones de resumen diarias}

Si ha activado el envío de resúmenes diarios sobre el estado de sus chequeos, el
sistema enviará todos los días un correo electrónico con información. El aspecto
del correo que recibirá se puede ver en la figura \ref{fig:email-resumen}

\begin{figure}[hbtp]
  \centering
  \includegraphics[width=0.7\textwidth]{apendice_manual_usuario/email_resumen.png}
  \caption{Email diario de resumen}
  \label{fig:email-resumen}
\end{figure}


\section{Uso de la aplicación móvil}

La aplicación móvil \textbf{SiteUp Client} es una aplicación para para el
sistema operativo Android que permite a sus usuarios revisar los chequeos que
tiene dados de alta en la plataforma web SiteUp, así como recibir de manera
instantánea notificaciones \textit{push} sobre el estado de sus chequeos.

\subsection{Instalación}

La aplicación móvil se instala copiando y ejecutando el fichero APK en el
dispositivo. En un futuro se podrá encontrar en la \textit{Play Store} de
Google.

\subsection{Lanzamiento}

En la figura~\ref{fig:manual-android-launcher} se puede observar el aspecto de
la aplicación \textbf{SiteUp Client} en el lanzador de Android. Para ejecutarla,
simplemente pulse sobre el icono de la aplicación.

\begin{figure}[htbp]
  \centering
  \includegraphics[width=0.4\textwidth]{5_diseno/android-1}
  \caption{Icono de SiteUp en el lanzador}
  \label{fig:manual-android-launcher}
\end{figure}

\subsection{Carga de la aplicación}

Tras lanzar la aplicación, el sistema cargará los datos necesarios y se
comunicará con los servidores de Google para obtener los datos necesarios. La
apariencia de la aplicación durante este proceso se puede ver en la figura
\ref{fig:manual-android-login}.

\begin{figure}[htbp]
  \centering
  \includegraphics[width=0.4\textwidth]{5_diseno/android-2}
  \caption{Pantalla de carga de la aplicación móvil}
  \label{fig:manual-android-loading}
\end{figure}

\subsection{Inicio de sesión}

Tras la carga, la primera vez que lance la aplicación será necesario que inicie
sesión con las credenciales que utilizó en la plataforma web. Para ello, tendrá
que introducir su nombre de usuario y contraseña en el formulario que aparece en
la figura~\ref{fig:manual-android-loading}, y pulsar el botón \textbf{Sign in}.

\begin{figure}[htbp]
  \centering
  \includegraphics[width=0.4\textwidth]{5_diseno/android-3}
  \caption{Pantalla de inicio de sesión de la aplicación móvil}
  \label{fig:manual-android-login}
\end{figure}

La aplicación móvil se comunicará con los servidores centrales de SiteUp,
validará sus credenciales de usuario y, si son correctas, le dará acceso al
sistema. Si las credenciales de usuario no son correctas, el formulario mostrará
un error.

\subsection{Revisión de chequeos}

Si ha iniciado sesión correctamente, tendrá la opción de revisar la lista de
chequeos que tiene dados de alta en la plataforma web \textbf{SiteUp}. La
apariencia de la lista de chequeos es igual a la de la
figura~\ref{fig:manual-android-dashboard}.

\begin{figure}[hbtp]
  \centering
  \includegraphics[width=0.4\textwidth]{5_diseno/android-4}
  \caption{Pantalla de chequeos de la aplicación móvil}
  \label{fig:manual-android-dashboard}
\end{figure}

Si pulsa en cualquiera de los chequeos, será redirigido a la plataforma web, que
podrá navegar sin problema gracias a que se encuentra preparada para su uso en
dispositivos móviles.

\subsection{Recepción de notificaciones}

La utilidad principal de la aplicación es la posibilidad de recibir
notificaciones en cualquier momento provenientes del servidor -- conocidas como
notificaciones \textit{push}. 

El procedimiento es transparente al usuario. Simplemente lance la aplicación e
inice sesión. Tras hacer esto, la aplicación estará preparada para recibir
notificaciones desde el servidor.

Cuando se reciba una notificación, la aplicación mostará un mensaje en el área
de notificaciones del móvil, con información resumida. Puede ver el aspecto de
la notificación en la figura \ref{fig:manual-android-notification}. Si pulsa en
la notificación será llevado a los detalles del chequeo para obtener más
detalles sobre el origen del aviso.

\begin{figure}[hbtp]
  \centering
  \includegraphics[width=0.4\textwidth]{5_diseno/android-5}
  \caption{Aspecto de las notificaciones}
  \label{fig:manual-android-notification}
\end{figure}



%%% Local Variables:
%%% mode: latex
%%% TeX-master: "../memoria"
%%% End:


\chapter{Requisitos del sistema}
\label{chap:requisitos}
En este capítulo se explicarán en detalle las instrucciones de uso de la
plataforma web.

\section{Introducción a SiteUp}

\textbf{SiteUp} es un servicio web creado para la monitorización de servicios de
Internet, especialmente pensado para dueños de sitios web y servidores que
deseen estar informados en caso de cualquier eventualidad.

El uso de internet en todos los ámbitos de la vida, desde las interacciones
sociales hasta la búsqueda de empleo, han dado lugar a la creación de servicios
accesibles a través de Internet y a nuevos modelos de negocio basados en
aquellos. Resulta por tanto imprescindible contar con sistemas para monitorizar
la disponibilidad de estos servicios y, en caso necesario, actuar de manera que
puedan subsanarse las causas de los problemas detectados.

SiteUp cubre precisamente estas necesidades, permitiendo a los usuarios:

\begin{itemize}
\item Crear chequeos mediante envío de paquetes ICMP (Ping), para revisar que un
  servidor está activo y conectado a la red.
\item Crear chequeos de registros DNS, de forma que el contenido de estos
  registros sea el correcto en todo momento.
\item Crear chequeos de puertos remotos, para vigilar que los servicios
  concretos dentro de un servidor son accesibles.
\item Crear chequeos mediante peticiones HTTP, para controlar que un servicio
  web responde de forma correcta y con el contenido adecuado.
\item Establecer las opciones de periodicidad de estos chequeos, de acuerdo a
  las necesidades particulares de cada caso.
\item Recibir notificaciones tanto a través de correo electrónico como mediante
  una aplicación móvil para el sistema operativo Android.
\end{itemize}

\section{Requisitos previos}

Para el uso correcto de la plataforma web SiteUp, el principal requisito es
contar con un dispositivo con acceso a internet que disponga de un navegador
reciente, como por ejemplo Google Chrome. Cualquier tipo de tablet, teléfono
móvil u ordenador personal podrá acceder a la plataforma web SiteUp, pero para
aprovechar mejor el potencial del sistema se recomienda usar una computadora
personal con las siguientes características:

\begin{itemize}
\item Navegador reciente. Se recomienda el uso de Google Chrome.
\item Pantalla con resolución mínima de 1280 píxeles en su lado longitudinal.
\item Conexión a internet con un caudal de descarga de al menos 512kb/s.
\end{itemize}

Si desea hacer uso de la aplicación móvil, será necesario que el usuario cuente
con un dispositivo con las siguientes características mínimas:

\begin{itemize}
\item Sistema operativo Android, versión mínima 2.3, versión recomendada 4.3.
\item Pantalla de al menos 3 pulgadas de diagonal y resolución mínima de 320x480
  píxeles.
\end{itemize}

\section{Uso de la plataforma web}

\subsection{Acceso}

A fecha de abril de 2014, la dirección de acceso a la plataforma web de SiteUp
es \url{http://siteup.josetomastocino.com}. Simplemente introduzca esa dirección
en su navegador y automáticamente será redirigido a la plataforma web, donde
verá la pantalla de bienvenida.

\subsection{Registro de usuario}

\begin{figure}[hbtp]
  \centering
  \includegraphics[width=0.3\textwidth]{apendice_manual_usuario/general_botones.png}
  \caption{Botones principales para el usuario anónimo}
  \label{fig:botones-1}
\end{figure}

Si es la primera vez que accede al sitio web deberá registrarse. Para ello, debe
pulsar en el botón \textbf{Sign up}, que se ve en la figura \ref{fig:botones-1}.

Tras ello, aparecerá el formulario de registro de usuario, que se ve en la
figura \ref{fig:registro}

\begin{figure}[hbtp]
  \centering
  \includegraphics[width=0.5\textwidth]{apendice_manual_usuario/formulario_registro.png}
  \caption{Formulario de registro de usuario}
  \label{fig:registro}
\end{figure}

En el formulario deberá introducir los siguientes datos:

\begin{itemize}
\item \textbf{Nombre de usuario}. Debe tener 30 caracteres como máximo, y solo
  pueden usarse letras, números y guiones.
\item \textbf{Dirección de correo electrónico}. Asegúrese de que es una
  dirección de correo correcta.
\item \textbf{Contraseña}.
\end{itemize}

Tras introducir los datos, deberá pulsar el botón \textbf{Sign up} del
formulario. Si el registro es correcto, el sistema dará de alta su usuario y
mostrará la pantalla \ref{fig:registro-correcto}. Si no, volverá a aparecer el
formulario con información sobre los errores en los datos.

\begin{figure}[hbtp]
  \centering
  \includegraphics[width=\textwidth]{apendice_manual_usuario/pantalla_register_completo.png}
  \caption{Mensaje indicando el registro correcto}
  \label{fig:registro-correcto}
\end{figure}

\subsection{Inicio de sesión}

Una vez que haya registrado su usuario, deberá pulsar en el botón \textbf{Log
  in}, como aparece en la imagen \ref{fig:botones-1}. Esto le llevará al
formulario de inicio de sesión que aparece en la figura \ref{fig:login}

\begin{figure}[hbtp]
  \centering
  \includegraphics[width=0.5\textwidth]{apendice_manual_usuario/general_formulario_login.png}
  \caption{Formulario de inicio de sesión}
  \label{fig:login}
\end{figure}

Deberá introducir su nombre de usuario y la contraseña que escribió al
registrarse, y posteriormente pulsar en el botón \textbf{Log in} del
formulario. Si los datos son correctos, será llevado al panel principal de la
web, donde podrá gestionar sus chequeos. Si no son correctos, volverá a aparecer
el formulario, con información sobre los errores encontrados.

\subsection{Recuperación de contraseña}

Si se registró como usuario pero no recuerda su contraseña es posible
recuperarla.

Para ello, vaya al formulario de inicio de sesión como ser ha explicado en el
paso anterior. En ese formulario, en la zona inferior, verá un enlace titulado
\textbf{Forgot password?} que deberá pulsar. Después de hacerlo se mostrará el
formulario de recuperación de contraseña, ilustrado en la figura
\ref{fig:formulario-reset}.

\begin{figure}[hbtp]
  \centering
  \includegraphics[width=0.5\textwidth]{apendice_manual_usuario/general_formulario_reset.png}
  \caption{Formulario de recuperación de contraseña}
  \label{fig:formulario-reset}
\end{figure}

En este formulario deberá introducir la dirección de correo electrónico que
utilizó al registrarse, y pulsar el botón \textbf{Send reset code}. Si lo hace
correctamente, aparecerá la pantalla de confirmación que se ve en la figura
\ref{fig:pantalla-reset-enviado}.

\begin{figure}[hbtp]
  \centering
  \includegraphics[width=0.5\textwidth]{apendice_manual_usuario/pantalla_reset_enviado.png}
  \caption{Pantalla de confirmación de envío de código de reinicio}
  \label{fig:pantalla-reset-enviado}
\end{figure}

El sistema habrá enviado un enlace a su dirección de correo electrónico. Deberá
pulsar el enlace que aparece en ese correo. Asegúrese de que recibe
correctamente el correo. Si no lo recibe, revise la carpeta de correo no
deseado. El enlace del correo le llevará a un formulario de reseteo de
contraseña como el que aparece en la figura \ref{fig:formulario-reset-2}.

\begin{figure}[hbtp]
  \centering
  \includegraphics[width=0.5\textwidth]{apendice_manual_usuario/general_formulario_reset_2.png}
  \caption{}
  \label{fig:formulario-reset-2}
\end{figure}

En este formulario deberá escribir la nueva contraseña de su cuenta dos veces,
para confirmar que la escribe correctamente. Tras ello, deberá pulsar el botón
\textbf{Set new password}, lo cual confirmará los cambios. 

Tras este proceso, podrá iniciar sesión como se ha descrito previamente, esta
vez utilizando la nueva contraseña.


\subsection{Pantalla de listado de chequeos}

\begin{figure}[hbtp]
  \centering
  \includegraphics[width=0.8\textwidth]{apendice_manual_usuario/pantalla-dashboard.png}
  \caption{Panel de control}
  \label{fig:dashboard}
\end{figure}

\begin{figure}[hbtp]
  \centering
  \includegraphics[width=0.4\textwidth]{apendice_manual_usuario/general_botones_2.png}
  \caption{Botones principales para el usuario logueado}
  \label{fig:botones-2}
\end{figure}

En la figura \ref{fig:dashboard} se puede ver el aspecto del panel de control de
la plataforma web SiteUp. Desde este panel podrá hacer una gran cantidad de
acciones, desde crear grupos de chequeos hasta editar sus detalles de usuario.

Es posible acceder a esta pantalla desde cualquier sitio de la web si está
logueado, pulsando en el botón \textbf{Dashboard} de la cabecera de la web, tal
y como se ve en la figura \ref{fig:botones-2}.

\subsection{Cerrar sesión}

Si ha iniciado sesión y desea cerrarla, debe dirigirse a los botones
principales, situados en el lado derecho de la cabecera de la web, tal y como
aparecen en la figura \ref{fig:botones-2}. Deberá pulsar el botón \textbf{Log
  out}, tras lo cual se cerrara su sesión y será redirigido a la pantalla
inicial de SiteUp.

\subsection{Edición de datos de usuario}

Si desea editar alguno de los datos de usuario que ingresó cuando hizo el
registro, deberá primero pulsar el botón \textbf{Your profile}, situado en la
cabecera de la web. Es el segundo de los botones que aparecen, tal y como se ve
en la figura \ref{fig:botones-2}.

Una vez hecho esto, aparecerá el formulario de edición de usuario, ilustrado en
la figura \ref{fig:formulario-usuario}, en el que verá sus actuales datos de
usuario. 

\begin{figure}[hbtp]
  \centering
  \includegraphics[width=0.8\textwidth]{apendice_manual_usuario/pantalla_perfil_usuario.png}
  \caption{}
  \label{fig:formulario-usuario}
\end{figure}

En este formulario podrá editar:

\begin{itemize}
\item Su \textbf{nombre} de usuario.
\item Su \textbf{dirección} de correo electrónico.
\item Si desea o no recibir \textbf{resúmenes diarios} sobre el estado de sus chequeos.
\end{itemize}

Una vez editados los datos que desee deberá pulsar el botón \textbf{Update
  details}. También podrá editar su contraseña pulsando el botón \textbf{Change
  password}, situado al pie del formulario a la izquierda.


\subsection{Creación de grupo de chequeos}

En SiteUp, los chequeos se organizan en \textbf{grupos de chequeos}, lo que
permite trabajar con ellos más fácilmente a la hora de activarlos, desactivarlos
o borrarlos. 

Para crear un grupo de chequeos deberá pulsar el botón \textbf{Add group},
situado en la zona superior derecha del panel principal. El botón puede verse en
la figura \ref{fig:crear-grupo}.

\begin{figure}[hbtp]
  \centering
  \includegraphics[width=0.3\textwidth]{apendice_manual_usuario/grupo_boton_nuevo.png}
  \caption{Botón para crear nuevo grupo de chequeos}
  \label{fig:crear-grupo}
\end{figure}

Tras pulsar el botón aparecerá el formulario de creación de nuevo grupo. En él
deberá introducir los datos del nuevo grupo. En particular:

\begin{itemize}
\item El \textbf{título} que identifica al grupo.
\end{itemize}

Tras ello, deberá pulsar el botón \textbf{Create group}, y el grupo será creado

\begin{figure}[H]
  \centering
  \includegraphics[width=0.8\textwidth]{apendice_manual_usuario/grupo_formulario_crear.png}
  \caption{Formulario de creación de grupo de chequeos}
  \label{fig:}
\end{figure}


\subsection{Operaciones con grupos de chequeos}

En la figura~\ref{fig:operaciones-grupo} se puede ver el menú de operaciones
disponibles para un grupo de chequeos. Este menú aparece en el panel principal,
al lado derecho junto al título de cada grupo de chequeos.

\begin{figure}[hbtp]
  \centering
  \includegraphics[width=0.5\textwidth]{apendice_manual_usuario/grupo_submenu_acciones.png}
  \caption{}
  \label{fig:operaciones-grupo}
\end{figure}

\subsubsection{Edición de grupo de chequeos}

En el menú de operaciones del grupo de chequeos, pulse el botón
\textbf{Edit}. Aparecerá un formulario similar al de creación de grupo, en el
que podrá editar el título del grupo de chequeos.

\subsubsection{Borrado de grupo de chequeos}

Para borrar un grupo de chequeos, pulse el botón \textbf{Delete} en el menú de
operaciones del grupo. Aparecerá un formulario como el de la figura
\ref{fig:formulario-borrado-grupo}, pidiéndole confirmación del borrado. 

\begin{figure}[hbtp]
  \centering
  \includegraphics[width=0.5\textwidth]{apendice_manual_usuario/formulario_borrado_grupo}
  \caption{Formulario de confirmación de borrado}
  \label{fig:formulario-borrado-grupo}
\end{figure}

Si pulsa en el botón \textbf{Go back}, volverá al panel principal, cancelando el
borrado. Si pulsa el botón \textbf{Delete}, confirmará la operación y el sistema
borrará el grupo de chequeos, así como todos los chequeos que contenga.

\subsubsection{Activación y desactivación de grupo de chequeos}

Es posible activar y desactivar de forma grupal todos los chequeos de un
grupo. Para ello, deberá dirigirse al menú de operaciones del grupo y pulsar en
\textbf{Enable all checks} o en \textbf{Disable all checks}, según sea su decisión.

\subsection{Creación de chequeo}

\begin{figure}[hbtp]
  \centering
  \includegraphics[width=0.8\textwidth]{apendice_manual_usuario/pantalla_seleccion_tipo.png}
  \caption{Pantalla de selección del tipo de chequeo a crear}
  \label{fig:seleccionar-tipo-check}
\end{figure}

Como se ha comentado previamente, el elemento principal de SiteUp es el
\textbf{chequeo}. El usuario da de alta chequeos dentro de un grupo, y estos
chequeos se lanzan periódicamente, generando datos.

Para crear un chequeo dentro de un grupo deberá ir al menú de operaciones del
grupo, que aparece en la figura \ref{fig:operaciones-grupo} y se sitúa a la
derecha del título de un grupo de chequeos. Dentro del menú, deberá pulsar la
opción \textbf{Add check}.



Tras esto aparecerá una pantalla en la que deberá seleccionar el tipo de chequeo
que desea crear. La pantalla es similar a la de la figura
\ref{fig:seleccionar-tipo-check}. Como se explicó en la introducción de este
manual, es posible crear cuatro tipo de chequeos: chequeos de ping, chequeos de
puerto, chequeos de registros DNS y chequeos por HTTP.

Si pulsa en el botón del tipo que desea crear, aparecerá un formulario con
campos para rellenar.

\subsubsection{Creación de un chequeo por ping}

\begin{figure}[hbtp]
  \centering
  \includegraphics[width=0.8\textwidth]{apendice_manual_usuario/pantalla_crear_ping.png}
  \caption{Formulario de creación de chequeo por ping}
  \label{fig:crear-ping}
\end{figure}

El formulario para crear un chequeo de tipo ping aparece en la
figura~\ref{fig:crear-ping}. Los campos que deberá rellenar son los siguientes:

\begin{itemize}
\item \textbf{Target}: es el objetivo a chequear. Debe ser un \textit{hostname} o una IP.
\item \textbf{Título} del chequeo.
\item \textbf{Descripción} opcional del chequeo.
\item \textbf{Check interval}: intervalo entre chequeos, definido en minutos. El
  valor mínimo es de 1 minuto.
\item \textbf{Notify email}: indica si los cambios de estado en este chequeo
  deberán ser notificados por correo electrónico.
\item \textbf{Notify Android}: indica si los cambios de estado en este chequeo
  deberán ser notificados mediante la aplicación Android.
\item \textbf{Should check timeout}: indica si hay que verificar que el ping
  enviado vuelve dentro de un tiempo máximo determinado.
\item \textbf{Timeout value}: si la anterior opción está activa, este campo
  define el número de milisegundos que puede tardar en responder el servidor al
  ping, como máximo.
\end{itemize}

\subsubsection{Creación de un chequeo de puerto}

\begin{figure}[hbtp]
  \centering
  \includegraphics[width=0.8\textwidth]{apendice_manual_usuario/pantalla_crear_port.png}
  \caption{Formulario de creación de chequeo de puerto}
  \label{fig:crear-port}
\end{figure}

El formulario para crear un chequeo de puertos aparece en la
figura~\ref{fig:crear-port}. Los campos que deberá rellenar son los siguientes:

\begin{itemize}
\item \textbf{Título} del chequeo.
\item \textbf{Descripción} opcional del chequeo.
\item \textbf{Target}: es el objetivo a chequear. Debe ser un \textit{hostname} o una IP.
\item \textbf{Target port}: es el puerto a chequear. Debe ser un entero.
\item \textbf{Check for string}: opcionalmente, se puede indicar una cadena que
  debe estar presente en la respuesta al conectar con el puerto.
\item \textbf{Check interval}: intervalo entre chequeos, definido en minutos. El
  valor mínimo es de 1 minuto.
\item \textbf{Notify email}: indica si los cambios de estado en este chequeo
  deberán ser notificados por correo electrónico.
\item \textbf{Notify Android}: indica si los cambios de estado en este chequeo
  deberán ser notificados mediante la aplicación Android.
\end{itemize}

\subsubsection{Creación de un chequeo de DNS}

\begin{figure}[hbtp]
  \centering
  \includegraphics[width=0.8\textwidth]{apendice_manual_usuario/pantalla_crear_dns.png}
  \caption{Formulario de creación de chequeo de registro DNS}
  \label{fig:crear-dns}
\end{figure}

El formulario para crear un chequeo de registro DNS aparece en la figura
\ref{fig:crear-dns}. Los campos que deberá rellenar son los siguientes:

\begin{itemize}
\item \textbf{Título} del chequeo.
\item \textbf{Descripción} opcional del chequeo.
\item \textbf{Target}: es el objetivo a chequear. Debe ser un \textit{hostname}.
\item \textbf{Record type}: indica el tipo de registro a chequear. Puede ser A, AAAA, CNAME, TXT y MX.
\item \textbf{Resolved address}: indica el contenido que debe tener el registro.
\item \textbf{Check interval}: intervalo entre chequeos, definido en minutos. El
  valor mínimo es de 1 minuto.
\item \textbf{Notify email}: indica si los cambios de estado en este chequeo
  deberán ser notificados por correo electrónico.
\item \textbf{Notify Android}: indica si los cambios de estado en este chequeo
  deberán ser notificados mediante la aplicación Android.
\end{itemize}

\subsubsection{Creación de un chequeo por HTTP}

\begin{figure}[hbtp]
  \centering
  \includegraphics[width=0.8\textwidth]{apendice_manual_usuario/pantalla_crear_http.png}
  \caption{Formulario de creación de chequeo por HTTP}
  \label{fig:crear-http}
\end{figure}

El formulario para crear un chequeo por HTTP aparece en la figura
\ref{fig:crear-http}. Los campos que deberá rellenar son los siguientes:

\begin{itemize}
\item \textbf{Título} del chequeo.
\item \textbf{Descripción} opcional del chequeo.
\item \textbf{Target}: es el objetivo a chequear. Debe ser una URL.
\item \textbf{Status code}: es el código de estado que deberá devolver la petición..
\item \textbf{Check for string}: opcionalmente, se puede indicar una cadena que
  debe estar presente en la respuesta a la petición.
\item \textbf{Check interval}: intervalo entre chequeos, definido en minutos. El
  valor mínimo es de 1 minuto.
\item \textbf{Notify email}: indica si los cambios de estado en este chequeo
  deberán ser notificados por correo electrónico.
\item \textbf{Notify Android}: indica si los cambios de estado en este chequeo
  deberán ser notificados mediante la aplicación Android.
\end{itemize}


\subsection{Operaciones con chequeos}

\begin{figure}[hbtp]
  \centering
  \includegraphics[width=\textwidth]{apendice_manual_usuario/chequeo_valido.png}
  \caption{}
  \label{fig:chequeo}
\end{figure}

Una vez creado un chequeo, aparecerá en el panel de control con el aspecto de la
figura~\ref{fig:chequeo}. Como se aprecia en la figura, se muestra mediante un
icono el estado del chequeo, los datos (título, tipo, descripción) y una gráfica
que ilustra el estado del chequeo en las últimas 24 horas.

\begin{figure}[hbtp]
  \centering
  \includegraphics[width=\textwidth]{apendice_manual_usuario/chequeo_invalido.png}
  \caption{}
  \label{fig:chequeo-invalido}
\end{figure}

Si un chequeo detecta un error, cambia de apariencia, pasando a tener el aspecto
que ofrece la figura \ref{fig:chequeo-invalido}. Además, según se haya
configurado en las opciones del chequeo, se envía una notificación al usuario
informando del cambio de estado del chequeo.

Cuando se ve en la lista general de chequeos, cada chequeo tiene una serie de
botones que permiten llevar a cabo operaciones diversas, que se detallan a
continuación.

\subsubsection{Activación y desactivación de un chequeo}

Para desactivar un chequeo, tiene que pulsar el botón \textbf{Disable} de la
botonera, que solo aparecerá si el chequeo está activado. De igual modo, si el
chequeo está desactivado, el botón que aparecerá tendrá la etiqueta
\textbf{Enable}, que podrá pulsar para activarlo.

\subsubsection{Borrado de un chequeo}

En la botonera se sitúa el botón \textbf{Delete}, que podrá pulsar para borrar
el chequeo. Aparecerá un formulario pidiéndole confirmación. Si confirma la
acción, se procederá al borrado del chequeo y todos los datos dependientes de
éste.

\subsubsection{Edición de un chequeo}

Si desea editar un chequeo debe pulsar el botón \textbf{Edit}, situado en la
botonera del chequeo. Aparecerá un formulario similar al de creación de chequeo,
pero con los datos actuales ya introducidos. Puede modificar los datos que desee
y confirmar los cambios pulsando el botón \textbf{Update check} del formulario.

\subsubsection{Vista de detalle del chequeo}

\begin{figure}[hbtp]
  \centering
  \includegraphics[width=\textwidth]{apendice_manual_usuario/pantalla_detalle_chequeo.png}
  \caption{}
  \label{fig:detalle-chequeo}
\end{figure}

En la botonera del chequeo, en el panel de control, verá un botón
\textbf{Details}. Si lo pulsa será llevado a la pantalla de vista de detalle del
chequeo, que se ilustra en la figura \ref{fig:detalle-chequeo}. Esta pantalla
ofrece mucha información sobre el chequeo. En particular:

\begin{itemize}
\item En la parte superior se sitúa la información general:
  \begin{itemize}
  \item Icono con el estado del chequeo.
  \item Título.
  \item Tipo de chequeo.
  \item Descripción.
  \item Botonera de acciones.
  \end{itemize}

\item En la parte media de la pantalla se sitúa un bloque con pestañas. Mediante
  estas pestañas es posible indicar sobre qué periodo de tiempo queremos
  visualizar los datos: últimas 24 horas o última semana. La información
  disponible es la siguiente:

  \begin{itemize}
  \item Porcentaje de \textit{uptime} del chequeo.
  \item \textit{(Solo chequeos tipo ping)} Tiempo de respuesta mínimo, máximo y promedio.
  \item Gráfica con el estado del chequeo.
  \end{itemize}

\item En la parte inferior de la pantalla se sitúa el panel de eventos del
  chequeo. Éste permite ver los últimos cambios de estado del chequeo:

  \begin{itemize}
  \item Tipo de cambio de estado (\textit{Up}, \textit{Down}).
  \item Duración del estado.
  \item Hora de comienzo del estado.
  \item Hora de fin del estado. Si es el estado en el que el chequeo se
    encuentra actualmente, no aparecerá ninguna fecha.
  \item Extra: muestra información adicional para estados de error.
  \end{itemize}
\end{itemize}

\begin{figure}[hbtp]
  \centering
  \includegraphics[width=\textwidth]{apendice_manual_usuario/pantalla_grafica_semanal.png}
  \caption{Detalle de la gráfica para un intervalo de tiempo de una semana}
  \label{fig:grafica-semanal}
\end{figure} 

\subsection{Recepción de notificaciones por correo electrónico}

\subsubsection{Notificaciones de cambio de estado}

Cuando el sistema detecte un cambio de estado en un chequeo, y esté activada la
notificación por correo electrónico, se enviará un email como el que aparece en
la figura \ref{fig:email-estado}.

\begin{figure}[hbtp]
  \centering
  \includegraphics[width=0.7\textwidth]{apendice_manual_usuario/email_estado.png}
  \caption{Email de notificación de cambio de estado}
  \label{fig:email-estado}
\end{figure}

En ese email se le indicará el título del chequeo, el nuevo estado, la fecha y
hora en la que ha pasado al nuevo estado, y un enlace directo al panel de
detalle del chequeo en SiteUp, de forma que pueda obtener más información.

\subsubsection{Notificaciones de resumen diarias}

Si ha activado el envío de resúmenes diarios sobre el estado de sus chequeos, el
sistema enviará todos los días un correo electrónico con información. El aspecto
del correo que recibirá se puede ver en la figura \ref{fig:email-resumen}

\begin{figure}[hbtp]
  \centering
  \includegraphics[width=0.7\textwidth]{apendice_manual_usuario/email_resumen.png}
  \caption{Email diario de resumen}
  \label{fig:email-resumen}
\end{figure}


\section{Uso de la aplicación móvil}

La aplicación móvil \textbf{SiteUp Client} es una aplicación para para el
sistema operativo Android que permite a sus usuarios revisar los chequeos que
tiene dados de alta en la plataforma web SiteUp, así como recibir de manera
instantánea notificaciones \textit{push} sobre el estado de sus chequeos.

\subsection{Instalación}

La aplicación móvil se instala copiando y ejecutando el fichero APK en el
dispositivo. En un futuro se podrá encontrar en la \textit{Play Store} de
Google.

\subsection{Lanzamiento}

En la figura~\ref{fig:manual-android-launcher} se puede observar el aspecto de
la aplicación \textbf{SiteUp Client} en el lanzador de Android. Para ejecutarla,
simplemente pulse sobre el icono de la aplicación.

\begin{figure}[htbp]
  \centering
  \includegraphics[width=0.4\textwidth]{5_diseno/android-1}
  \caption{Icono de SiteUp en el lanzador}
  \label{fig:manual-android-launcher}
\end{figure}

\subsection{Carga de la aplicación}

Tras lanzar la aplicación, el sistema cargará los datos necesarios y se
comunicará con los servidores de Google para obtener los datos necesarios. La
apariencia de la aplicación durante este proceso se puede ver en la figura
\ref{fig:manual-android-login}.

\begin{figure}[htbp]
  \centering
  \includegraphics[width=0.4\textwidth]{5_diseno/android-2}
  \caption{Pantalla de carga de la aplicación móvil}
  \label{fig:manual-android-loading}
\end{figure}

\subsection{Inicio de sesión}

Tras la carga, la primera vez que lance la aplicación será necesario que inicie
sesión con las credenciales que utilizó en la plataforma web. Para ello, tendrá
que introducir su nombre de usuario y contraseña en el formulario que aparece en
la figura~\ref{fig:manual-android-loading}, y pulsar el botón \textbf{Sign in}.

\begin{figure}[htbp]
  \centering
  \includegraphics[width=0.4\textwidth]{5_diseno/android-3}
  \caption{Pantalla de inicio de sesión de la aplicación móvil}
  \label{fig:manual-android-login}
\end{figure}

La aplicación móvil se comunicará con los servidores centrales de SiteUp,
validará sus credenciales de usuario y, si son correctas, le dará acceso al
sistema. Si las credenciales de usuario no son correctas, el formulario mostrará
un error.

\subsection{Revisión de chequeos}

Si ha iniciado sesión correctamente, tendrá la opción de revisar la lista de
chequeos que tiene dados de alta en la plataforma web \textbf{SiteUp}. La
apariencia de la lista de chequeos es igual a la de la
figura~\ref{fig:manual-android-dashboard}.

\begin{figure}[hbtp]
  \centering
  \includegraphics[width=0.4\textwidth]{5_diseno/android-4}
  \caption{Pantalla de chequeos de la aplicación móvil}
  \label{fig:manual-android-dashboard}
\end{figure}

Si pulsa en cualquiera de los chequeos, será redirigido a la plataforma web, que
podrá navegar sin problema gracias a que se encuentra preparada para su uso en
dispositivos móviles.

\subsection{Recepción de notificaciones}

La utilidad principal de la aplicación es la posibilidad de recibir
notificaciones en cualquier momento provenientes del servidor -- conocidas como
notificaciones \textit{push}. 

El procedimiento es transparente al usuario. Simplemente lance la aplicación e
inice sesión. Tras hacer esto, la aplicación estará preparada para recibir
notificaciones desde el servidor.

Cuando se reciba una notificación, la aplicación mostará un mensaje en el área
de notificaciones del móvil, con información resumida. Puede ver el aspecto de
la notificación en la figura \ref{fig:manual-android-notification}. Si pulsa en
la notificación será llevado a los detalles del chequeo para obtener más
detalles sobre el origen del aviso.

\begin{figure}[hbtp]
  \centering
  \includegraphics[width=0.4\textwidth]{5_diseno/android-5}
  \caption{Aspecto de las notificaciones}
  \label{fig:manual-android-notification}
\end{figure}



%%% Local Variables:
%%% mode: latex
%%% TeX-master: "../memoria"
%%% End:


\chapter{Análisis}
\label{chap:analisis}
En este capítulo se explicarán en detalle las instrucciones de uso de la
plataforma web.

\section{Introducción a SiteUp}

\textbf{SiteUp} es un servicio web creado para la monitorización de servicios de
Internet, especialmente pensado para dueños de sitios web y servidores que
deseen estar informados en caso de cualquier eventualidad.

El uso de internet en todos los ámbitos de la vida, desde las interacciones
sociales hasta la búsqueda de empleo, han dado lugar a la creación de servicios
accesibles a través de Internet y a nuevos modelos de negocio basados en
aquellos. Resulta por tanto imprescindible contar con sistemas para monitorizar
la disponibilidad de estos servicios y, en caso necesario, actuar de manera que
puedan subsanarse las causas de los problemas detectados.

SiteUp cubre precisamente estas necesidades, permitiendo a los usuarios:

\begin{itemize}
\item Crear chequeos mediante envío de paquetes ICMP (Ping), para revisar que un
  servidor está activo y conectado a la red.
\item Crear chequeos de registros DNS, de forma que el contenido de estos
  registros sea el correcto en todo momento.
\item Crear chequeos de puertos remotos, para vigilar que los servicios
  concretos dentro de un servidor son accesibles.
\item Crear chequeos mediante peticiones HTTP, para controlar que un servicio
  web responde de forma correcta y con el contenido adecuado.
\item Establecer las opciones de periodicidad de estos chequeos, de acuerdo a
  las necesidades particulares de cada caso.
\item Recibir notificaciones tanto a través de correo electrónico como mediante
  una aplicación móvil para el sistema operativo Android.
\end{itemize}

\section{Requisitos previos}

Para el uso correcto de la plataforma web SiteUp, el principal requisito es
contar con un dispositivo con acceso a internet que disponga de un navegador
reciente, como por ejemplo Google Chrome. Cualquier tipo de tablet, teléfono
móvil u ordenador personal podrá acceder a la plataforma web SiteUp, pero para
aprovechar mejor el potencial del sistema se recomienda usar una computadora
personal con las siguientes características:

\begin{itemize}
\item Navegador reciente. Se recomienda el uso de Google Chrome.
\item Pantalla con resolución mínima de 1280 píxeles en su lado longitudinal.
\item Conexión a internet con un caudal de descarga de al menos 512kb/s.
\end{itemize}

Si desea hacer uso de la aplicación móvil, será necesario que el usuario cuente
con un dispositivo con las siguientes características mínimas:

\begin{itemize}
\item Sistema operativo Android, versión mínima 2.3, versión recomendada 4.3.
\item Pantalla de al menos 3 pulgadas de diagonal y resolución mínima de 320x480
  píxeles.
\end{itemize}

\section{Uso de la plataforma web}

\subsection{Acceso}

A fecha de abril de 2014, la dirección de acceso a la plataforma web de SiteUp
es \url{http://siteup.josetomastocino.com}. Simplemente introduzca esa dirección
en su navegador y automáticamente será redirigido a la plataforma web, donde
verá la pantalla de bienvenida.

\subsection{Registro de usuario}

\begin{figure}[hbtp]
  \centering
  \includegraphics[width=0.3\textwidth]{apendice_manual_usuario/general_botones.png}
  \caption{Botones principales para el usuario anónimo}
  \label{fig:botones-1}
\end{figure}

Si es la primera vez que accede al sitio web deberá registrarse. Para ello, debe
pulsar en el botón \textbf{Sign up}, que se ve en la figura \ref{fig:botones-1}.

Tras ello, aparecerá el formulario de registro de usuario, que se ve en la
figura \ref{fig:registro}

\begin{figure}[hbtp]
  \centering
  \includegraphics[width=0.5\textwidth]{apendice_manual_usuario/formulario_registro.png}
  \caption{Formulario de registro de usuario}
  \label{fig:registro}
\end{figure}

En el formulario deberá introducir los siguientes datos:

\begin{itemize}
\item \textbf{Nombre de usuario}. Debe tener 30 caracteres como máximo, y solo
  pueden usarse letras, números y guiones.
\item \textbf{Dirección de correo electrónico}. Asegúrese de que es una
  dirección de correo correcta.
\item \textbf{Contraseña}.
\end{itemize}

Tras introducir los datos, deberá pulsar el botón \textbf{Sign up} del
formulario. Si el registro es correcto, el sistema dará de alta su usuario y
mostrará la pantalla \ref{fig:registro-correcto}. Si no, volverá a aparecer el
formulario con información sobre los errores en los datos.

\begin{figure}[hbtp]
  \centering
  \includegraphics[width=\textwidth]{apendice_manual_usuario/pantalla_register_completo.png}
  \caption{Mensaje indicando el registro correcto}
  \label{fig:registro-correcto}
\end{figure}

\subsection{Inicio de sesión}

Una vez que haya registrado su usuario, deberá pulsar en el botón \textbf{Log
  in}, como aparece en la imagen \ref{fig:botones-1}. Esto le llevará al
formulario de inicio de sesión que aparece en la figura \ref{fig:login}

\begin{figure}[hbtp]
  \centering
  \includegraphics[width=0.5\textwidth]{apendice_manual_usuario/general_formulario_login.png}
  \caption{Formulario de inicio de sesión}
  \label{fig:login}
\end{figure}

Deberá introducir su nombre de usuario y la contraseña que escribió al
registrarse, y posteriormente pulsar en el botón \textbf{Log in} del
formulario. Si los datos son correctos, será llevado al panel principal de la
web, donde podrá gestionar sus chequeos. Si no son correctos, volverá a aparecer
el formulario, con información sobre los errores encontrados.

\subsection{Recuperación de contraseña}

Si se registró como usuario pero no recuerda su contraseña es posible
recuperarla.

Para ello, vaya al formulario de inicio de sesión como ser ha explicado en el
paso anterior. En ese formulario, en la zona inferior, verá un enlace titulado
\textbf{Forgot password?} que deberá pulsar. Después de hacerlo se mostrará el
formulario de recuperación de contraseña, ilustrado en la figura
\ref{fig:formulario-reset}.

\begin{figure}[hbtp]
  \centering
  \includegraphics[width=0.5\textwidth]{apendice_manual_usuario/general_formulario_reset.png}
  \caption{Formulario de recuperación de contraseña}
  \label{fig:formulario-reset}
\end{figure}

En este formulario deberá introducir la dirección de correo electrónico que
utilizó al registrarse, y pulsar el botón \textbf{Send reset code}. Si lo hace
correctamente, aparecerá la pantalla de confirmación que se ve en la figura
\ref{fig:pantalla-reset-enviado}.

\begin{figure}[hbtp]
  \centering
  \includegraphics[width=0.5\textwidth]{apendice_manual_usuario/pantalla_reset_enviado.png}
  \caption{Pantalla de confirmación de envío de código de reinicio}
  \label{fig:pantalla-reset-enviado}
\end{figure}

El sistema habrá enviado un enlace a su dirección de correo electrónico. Deberá
pulsar el enlace que aparece en ese correo. Asegúrese de que recibe
correctamente el correo. Si no lo recibe, revise la carpeta de correo no
deseado. El enlace del correo le llevará a un formulario de reseteo de
contraseña como el que aparece en la figura \ref{fig:formulario-reset-2}.

\begin{figure}[hbtp]
  \centering
  \includegraphics[width=0.5\textwidth]{apendice_manual_usuario/general_formulario_reset_2.png}
  \caption{}
  \label{fig:formulario-reset-2}
\end{figure}

En este formulario deberá escribir la nueva contraseña de su cuenta dos veces,
para confirmar que la escribe correctamente. Tras ello, deberá pulsar el botón
\textbf{Set new password}, lo cual confirmará los cambios. 

Tras este proceso, podrá iniciar sesión como se ha descrito previamente, esta
vez utilizando la nueva contraseña.


\subsection{Pantalla de listado de chequeos}

\begin{figure}[hbtp]
  \centering
  \includegraphics[width=0.8\textwidth]{apendice_manual_usuario/pantalla-dashboard.png}
  \caption{Panel de control}
  \label{fig:dashboard}
\end{figure}

\begin{figure}[hbtp]
  \centering
  \includegraphics[width=0.4\textwidth]{apendice_manual_usuario/general_botones_2.png}
  \caption{Botones principales para el usuario logueado}
  \label{fig:botones-2}
\end{figure}

En la figura \ref{fig:dashboard} se puede ver el aspecto del panel de control de
la plataforma web SiteUp. Desde este panel podrá hacer una gran cantidad de
acciones, desde crear grupos de chequeos hasta editar sus detalles de usuario.

Es posible acceder a esta pantalla desde cualquier sitio de la web si está
logueado, pulsando en el botón \textbf{Dashboard} de la cabecera de la web, tal
y como se ve en la figura \ref{fig:botones-2}.

\subsection{Cerrar sesión}

Si ha iniciado sesión y desea cerrarla, debe dirigirse a los botones
principales, situados en el lado derecho de la cabecera de la web, tal y como
aparecen en la figura \ref{fig:botones-2}. Deberá pulsar el botón \textbf{Log
  out}, tras lo cual se cerrara su sesión y será redirigido a la pantalla
inicial de SiteUp.

\subsection{Edición de datos de usuario}

Si desea editar alguno de los datos de usuario que ingresó cuando hizo el
registro, deberá primero pulsar el botón \textbf{Your profile}, situado en la
cabecera de la web. Es el segundo de los botones que aparecen, tal y como se ve
en la figura \ref{fig:botones-2}.

Una vez hecho esto, aparecerá el formulario de edición de usuario, ilustrado en
la figura \ref{fig:formulario-usuario}, en el que verá sus actuales datos de
usuario. 

\begin{figure}[hbtp]
  \centering
  \includegraphics[width=0.8\textwidth]{apendice_manual_usuario/pantalla_perfil_usuario.png}
  \caption{}
  \label{fig:formulario-usuario}
\end{figure}

En este formulario podrá editar:

\begin{itemize}
\item Su \textbf{nombre} de usuario.
\item Su \textbf{dirección} de correo electrónico.
\item Si desea o no recibir \textbf{resúmenes diarios} sobre el estado de sus chequeos.
\end{itemize}

Una vez editados los datos que desee deberá pulsar el botón \textbf{Update
  details}. También podrá editar su contraseña pulsando el botón \textbf{Change
  password}, situado al pie del formulario a la izquierda.


\subsection{Creación de grupo de chequeos}

En SiteUp, los chequeos se organizan en \textbf{grupos de chequeos}, lo que
permite trabajar con ellos más fácilmente a la hora de activarlos, desactivarlos
o borrarlos. 

Para crear un grupo de chequeos deberá pulsar el botón \textbf{Add group},
situado en la zona superior derecha del panel principal. El botón puede verse en
la figura \ref{fig:crear-grupo}.

\begin{figure}[hbtp]
  \centering
  \includegraphics[width=0.3\textwidth]{apendice_manual_usuario/grupo_boton_nuevo.png}
  \caption{Botón para crear nuevo grupo de chequeos}
  \label{fig:crear-grupo}
\end{figure}

Tras pulsar el botón aparecerá el formulario de creación de nuevo grupo. En él
deberá introducir los datos del nuevo grupo. En particular:

\begin{itemize}
\item El \textbf{título} que identifica al grupo.
\end{itemize}

Tras ello, deberá pulsar el botón \textbf{Create group}, y el grupo será creado

\begin{figure}[H]
  \centering
  \includegraphics[width=0.8\textwidth]{apendice_manual_usuario/grupo_formulario_crear.png}
  \caption{Formulario de creación de grupo de chequeos}
  \label{fig:}
\end{figure}


\subsection{Operaciones con grupos de chequeos}

En la figura~\ref{fig:operaciones-grupo} se puede ver el menú de operaciones
disponibles para un grupo de chequeos. Este menú aparece en el panel principal,
al lado derecho junto al título de cada grupo de chequeos.

\begin{figure}[hbtp]
  \centering
  \includegraphics[width=0.5\textwidth]{apendice_manual_usuario/grupo_submenu_acciones.png}
  \caption{}
  \label{fig:operaciones-grupo}
\end{figure}

\subsubsection{Edición de grupo de chequeos}

En el menú de operaciones del grupo de chequeos, pulse el botón
\textbf{Edit}. Aparecerá un formulario similar al de creación de grupo, en el
que podrá editar el título del grupo de chequeos.

\subsubsection{Borrado de grupo de chequeos}

Para borrar un grupo de chequeos, pulse el botón \textbf{Delete} en el menú de
operaciones del grupo. Aparecerá un formulario como el de la figura
\ref{fig:formulario-borrado-grupo}, pidiéndole confirmación del borrado. 

\begin{figure}[hbtp]
  \centering
  \includegraphics[width=0.5\textwidth]{apendice_manual_usuario/formulario_borrado_grupo}
  \caption{Formulario de confirmación de borrado}
  \label{fig:formulario-borrado-grupo}
\end{figure}

Si pulsa en el botón \textbf{Go back}, volverá al panel principal, cancelando el
borrado. Si pulsa el botón \textbf{Delete}, confirmará la operación y el sistema
borrará el grupo de chequeos, así como todos los chequeos que contenga.

\subsubsection{Activación y desactivación de grupo de chequeos}

Es posible activar y desactivar de forma grupal todos los chequeos de un
grupo. Para ello, deberá dirigirse al menú de operaciones del grupo y pulsar en
\textbf{Enable all checks} o en \textbf{Disable all checks}, según sea su decisión.

\subsection{Creación de chequeo}

\begin{figure}[hbtp]
  \centering
  \includegraphics[width=0.8\textwidth]{apendice_manual_usuario/pantalla_seleccion_tipo.png}
  \caption{Pantalla de selección del tipo de chequeo a crear}
  \label{fig:seleccionar-tipo-check}
\end{figure}

Como se ha comentado previamente, el elemento principal de SiteUp es el
\textbf{chequeo}. El usuario da de alta chequeos dentro de un grupo, y estos
chequeos se lanzan periódicamente, generando datos.

Para crear un chequeo dentro de un grupo deberá ir al menú de operaciones del
grupo, que aparece en la figura \ref{fig:operaciones-grupo} y se sitúa a la
derecha del título de un grupo de chequeos. Dentro del menú, deberá pulsar la
opción \textbf{Add check}.



Tras esto aparecerá una pantalla en la que deberá seleccionar el tipo de chequeo
que desea crear. La pantalla es similar a la de la figura
\ref{fig:seleccionar-tipo-check}. Como se explicó en la introducción de este
manual, es posible crear cuatro tipo de chequeos: chequeos de ping, chequeos de
puerto, chequeos de registros DNS y chequeos por HTTP.

Si pulsa en el botón del tipo que desea crear, aparecerá un formulario con
campos para rellenar.

\subsubsection{Creación de un chequeo por ping}

\begin{figure}[hbtp]
  \centering
  \includegraphics[width=0.8\textwidth]{apendice_manual_usuario/pantalla_crear_ping.png}
  \caption{Formulario de creación de chequeo por ping}
  \label{fig:crear-ping}
\end{figure}

El formulario para crear un chequeo de tipo ping aparece en la
figura~\ref{fig:crear-ping}. Los campos que deberá rellenar son los siguientes:

\begin{itemize}
\item \textbf{Target}: es el objetivo a chequear. Debe ser un \textit{hostname} o una IP.
\item \textbf{Título} del chequeo.
\item \textbf{Descripción} opcional del chequeo.
\item \textbf{Check interval}: intervalo entre chequeos, definido en minutos. El
  valor mínimo es de 1 minuto.
\item \textbf{Notify email}: indica si los cambios de estado en este chequeo
  deberán ser notificados por correo electrónico.
\item \textbf{Notify Android}: indica si los cambios de estado en este chequeo
  deberán ser notificados mediante la aplicación Android.
\item \textbf{Should check timeout}: indica si hay que verificar que el ping
  enviado vuelve dentro de un tiempo máximo determinado.
\item \textbf{Timeout value}: si la anterior opción está activa, este campo
  define el número de milisegundos que puede tardar en responder el servidor al
  ping, como máximo.
\end{itemize}

\subsubsection{Creación de un chequeo de puerto}

\begin{figure}[hbtp]
  \centering
  \includegraphics[width=0.8\textwidth]{apendice_manual_usuario/pantalla_crear_port.png}
  \caption{Formulario de creación de chequeo de puerto}
  \label{fig:crear-port}
\end{figure}

El formulario para crear un chequeo de puertos aparece en la
figura~\ref{fig:crear-port}. Los campos que deberá rellenar son los siguientes:

\begin{itemize}
\item \textbf{Título} del chequeo.
\item \textbf{Descripción} opcional del chequeo.
\item \textbf{Target}: es el objetivo a chequear. Debe ser un \textit{hostname} o una IP.
\item \textbf{Target port}: es el puerto a chequear. Debe ser un entero.
\item \textbf{Check for string}: opcionalmente, se puede indicar una cadena que
  debe estar presente en la respuesta al conectar con el puerto.
\item \textbf{Check interval}: intervalo entre chequeos, definido en minutos. El
  valor mínimo es de 1 minuto.
\item \textbf{Notify email}: indica si los cambios de estado en este chequeo
  deberán ser notificados por correo electrónico.
\item \textbf{Notify Android}: indica si los cambios de estado en este chequeo
  deberán ser notificados mediante la aplicación Android.
\end{itemize}

\subsubsection{Creación de un chequeo de DNS}

\begin{figure}[hbtp]
  \centering
  \includegraphics[width=0.8\textwidth]{apendice_manual_usuario/pantalla_crear_dns.png}
  \caption{Formulario de creación de chequeo de registro DNS}
  \label{fig:crear-dns}
\end{figure}

El formulario para crear un chequeo de registro DNS aparece en la figura
\ref{fig:crear-dns}. Los campos que deberá rellenar son los siguientes:

\begin{itemize}
\item \textbf{Título} del chequeo.
\item \textbf{Descripción} opcional del chequeo.
\item \textbf{Target}: es el objetivo a chequear. Debe ser un \textit{hostname}.
\item \textbf{Record type}: indica el tipo de registro a chequear. Puede ser A, AAAA, CNAME, TXT y MX.
\item \textbf{Resolved address}: indica el contenido que debe tener el registro.
\item \textbf{Check interval}: intervalo entre chequeos, definido en minutos. El
  valor mínimo es de 1 minuto.
\item \textbf{Notify email}: indica si los cambios de estado en este chequeo
  deberán ser notificados por correo electrónico.
\item \textbf{Notify Android}: indica si los cambios de estado en este chequeo
  deberán ser notificados mediante la aplicación Android.
\end{itemize}

\subsubsection{Creación de un chequeo por HTTP}

\begin{figure}[hbtp]
  \centering
  \includegraphics[width=0.8\textwidth]{apendice_manual_usuario/pantalla_crear_http.png}
  \caption{Formulario de creación de chequeo por HTTP}
  \label{fig:crear-http}
\end{figure}

El formulario para crear un chequeo por HTTP aparece en la figura
\ref{fig:crear-http}. Los campos que deberá rellenar son los siguientes:

\begin{itemize}
\item \textbf{Título} del chequeo.
\item \textbf{Descripción} opcional del chequeo.
\item \textbf{Target}: es el objetivo a chequear. Debe ser una URL.
\item \textbf{Status code}: es el código de estado que deberá devolver la petición..
\item \textbf{Check for string}: opcionalmente, se puede indicar una cadena que
  debe estar presente en la respuesta a la petición.
\item \textbf{Check interval}: intervalo entre chequeos, definido en minutos. El
  valor mínimo es de 1 minuto.
\item \textbf{Notify email}: indica si los cambios de estado en este chequeo
  deberán ser notificados por correo electrónico.
\item \textbf{Notify Android}: indica si los cambios de estado en este chequeo
  deberán ser notificados mediante la aplicación Android.
\end{itemize}


\subsection{Operaciones con chequeos}

\begin{figure}[hbtp]
  \centering
  \includegraphics[width=\textwidth]{apendice_manual_usuario/chequeo_valido.png}
  \caption{}
  \label{fig:chequeo}
\end{figure}

Una vez creado un chequeo, aparecerá en el panel de control con el aspecto de la
figura~\ref{fig:chequeo}. Como se aprecia en la figura, se muestra mediante un
icono el estado del chequeo, los datos (título, tipo, descripción) y una gráfica
que ilustra el estado del chequeo en las últimas 24 horas.

\begin{figure}[hbtp]
  \centering
  \includegraphics[width=\textwidth]{apendice_manual_usuario/chequeo_invalido.png}
  \caption{}
  \label{fig:chequeo-invalido}
\end{figure}

Si un chequeo detecta un error, cambia de apariencia, pasando a tener el aspecto
que ofrece la figura \ref{fig:chequeo-invalido}. Además, según se haya
configurado en las opciones del chequeo, se envía una notificación al usuario
informando del cambio de estado del chequeo.

Cuando se ve en la lista general de chequeos, cada chequeo tiene una serie de
botones que permiten llevar a cabo operaciones diversas, que se detallan a
continuación.

\subsubsection{Activación y desactivación de un chequeo}

Para desactivar un chequeo, tiene que pulsar el botón \textbf{Disable} de la
botonera, que solo aparecerá si el chequeo está activado. De igual modo, si el
chequeo está desactivado, el botón que aparecerá tendrá la etiqueta
\textbf{Enable}, que podrá pulsar para activarlo.

\subsubsection{Borrado de un chequeo}

En la botonera se sitúa el botón \textbf{Delete}, que podrá pulsar para borrar
el chequeo. Aparecerá un formulario pidiéndole confirmación. Si confirma la
acción, se procederá al borrado del chequeo y todos los datos dependientes de
éste.

\subsubsection{Edición de un chequeo}

Si desea editar un chequeo debe pulsar el botón \textbf{Edit}, situado en la
botonera del chequeo. Aparecerá un formulario similar al de creación de chequeo,
pero con los datos actuales ya introducidos. Puede modificar los datos que desee
y confirmar los cambios pulsando el botón \textbf{Update check} del formulario.

\subsubsection{Vista de detalle del chequeo}

\begin{figure}[hbtp]
  \centering
  \includegraphics[width=\textwidth]{apendice_manual_usuario/pantalla_detalle_chequeo.png}
  \caption{}
  \label{fig:detalle-chequeo}
\end{figure}

En la botonera del chequeo, en el panel de control, verá un botón
\textbf{Details}. Si lo pulsa será llevado a la pantalla de vista de detalle del
chequeo, que se ilustra en la figura \ref{fig:detalle-chequeo}. Esta pantalla
ofrece mucha información sobre el chequeo. En particular:

\begin{itemize}
\item En la parte superior se sitúa la información general:
  \begin{itemize}
  \item Icono con el estado del chequeo.
  \item Título.
  \item Tipo de chequeo.
  \item Descripción.
  \item Botonera de acciones.
  \end{itemize}

\item En la parte media de la pantalla se sitúa un bloque con pestañas. Mediante
  estas pestañas es posible indicar sobre qué periodo de tiempo queremos
  visualizar los datos: últimas 24 horas o última semana. La información
  disponible es la siguiente:

  \begin{itemize}
  \item Porcentaje de \textit{uptime} del chequeo.
  \item \textit{(Solo chequeos tipo ping)} Tiempo de respuesta mínimo, máximo y promedio.
  \item Gráfica con el estado del chequeo.
  \end{itemize}

\item En la parte inferior de la pantalla se sitúa el panel de eventos del
  chequeo. Éste permite ver los últimos cambios de estado del chequeo:

  \begin{itemize}
  \item Tipo de cambio de estado (\textit{Up}, \textit{Down}).
  \item Duración del estado.
  \item Hora de comienzo del estado.
  \item Hora de fin del estado. Si es el estado en el que el chequeo se
    encuentra actualmente, no aparecerá ninguna fecha.
  \item Extra: muestra información adicional para estados de error.
  \end{itemize}
\end{itemize}

\begin{figure}[hbtp]
  \centering
  \includegraphics[width=\textwidth]{apendice_manual_usuario/pantalla_grafica_semanal.png}
  \caption{Detalle de la gráfica para un intervalo de tiempo de una semana}
  \label{fig:grafica-semanal}
\end{figure} 

\subsection{Recepción de notificaciones por correo electrónico}

\subsubsection{Notificaciones de cambio de estado}

Cuando el sistema detecte un cambio de estado en un chequeo, y esté activada la
notificación por correo electrónico, se enviará un email como el que aparece en
la figura \ref{fig:email-estado}.

\begin{figure}[hbtp]
  \centering
  \includegraphics[width=0.7\textwidth]{apendice_manual_usuario/email_estado.png}
  \caption{Email de notificación de cambio de estado}
  \label{fig:email-estado}
\end{figure}

En ese email se le indicará el título del chequeo, el nuevo estado, la fecha y
hora en la que ha pasado al nuevo estado, y un enlace directo al panel de
detalle del chequeo en SiteUp, de forma que pueda obtener más información.

\subsubsection{Notificaciones de resumen diarias}

Si ha activado el envío de resúmenes diarios sobre el estado de sus chequeos, el
sistema enviará todos los días un correo electrónico con información. El aspecto
del correo que recibirá se puede ver en la figura \ref{fig:email-resumen}

\begin{figure}[hbtp]
  \centering
  \includegraphics[width=0.7\textwidth]{apendice_manual_usuario/email_resumen.png}
  \caption{Email diario de resumen}
  \label{fig:email-resumen}
\end{figure}


\section{Uso de la aplicación móvil}

La aplicación móvil \textbf{SiteUp Client} es una aplicación para para el
sistema operativo Android que permite a sus usuarios revisar los chequeos que
tiene dados de alta en la plataforma web SiteUp, así como recibir de manera
instantánea notificaciones \textit{push} sobre el estado de sus chequeos.

\subsection{Instalación}

La aplicación móvil se instala copiando y ejecutando el fichero APK en el
dispositivo. En un futuro se podrá encontrar en la \textit{Play Store} de
Google.

\subsection{Lanzamiento}

En la figura~\ref{fig:manual-android-launcher} se puede observar el aspecto de
la aplicación \textbf{SiteUp Client} en el lanzador de Android. Para ejecutarla,
simplemente pulse sobre el icono de la aplicación.

\begin{figure}[htbp]
  \centering
  \includegraphics[width=0.4\textwidth]{5_diseno/android-1}
  \caption{Icono de SiteUp en el lanzador}
  \label{fig:manual-android-launcher}
\end{figure}

\subsection{Carga de la aplicación}

Tras lanzar la aplicación, el sistema cargará los datos necesarios y se
comunicará con los servidores de Google para obtener los datos necesarios. La
apariencia de la aplicación durante este proceso se puede ver en la figura
\ref{fig:manual-android-login}.

\begin{figure}[htbp]
  \centering
  \includegraphics[width=0.4\textwidth]{5_diseno/android-2}
  \caption{Pantalla de carga de la aplicación móvil}
  \label{fig:manual-android-loading}
\end{figure}

\subsection{Inicio de sesión}

Tras la carga, la primera vez que lance la aplicación será necesario que inicie
sesión con las credenciales que utilizó en la plataforma web. Para ello, tendrá
que introducir su nombre de usuario y contraseña en el formulario que aparece en
la figura~\ref{fig:manual-android-loading}, y pulsar el botón \textbf{Sign in}.

\begin{figure}[htbp]
  \centering
  \includegraphics[width=0.4\textwidth]{5_diseno/android-3}
  \caption{Pantalla de inicio de sesión de la aplicación móvil}
  \label{fig:manual-android-login}
\end{figure}

La aplicación móvil se comunicará con los servidores centrales de SiteUp,
validará sus credenciales de usuario y, si son correctas, le dará acceso al
sistema. Si las credenciales de usuario no son correctas, el formulario mostrará
un error.

\subsection{Revisión de chequeos}

Si ha iniciado sesión correctamente, tendrá la opción de revisar la lista de
chequeos que tiene dados de alta en la plataforma web \textbf{SiteUp}. La
apariencia de la lista de chequeos es igual a la de la
figura~\ref{fig:manual-android-dashboard}.

\begin{figure}[hbtp]
  \centering
  \includegraphics[width=0.4\textwidth]{5_diseno/android-4}
  \caption{Pantalla de chequeos de la aplicación móvil}
  \label{fig:manual-android-dashboard}
\end{figure}

Si pulsa en cualquiera de los chequeos, será redirigido a la plataforma web, que
podrá navegar sin problema gracias a que se encuentra preparada para su uso en
dispositivos móviles.

\subsection{Recepción de notificaciones}

La utilidad principal de la aplicación es la posibilidad de recibir
notificaciones en cualquier momento provenientes del servidor -- conocidas como
notificaciones \textit{push}. 

El procedimiento es transparente al usuario. Simplemente lance la aplicación e
inice sesión. Tras hacer esto, la aplicación estará preparada para recibir
notificaciones desde el servidor.

Cuando se reciba una notificación, la aplicación mostará un mensaje en el área
de notificaciones del móvil, con información resumida. Puede ver el aspecto de
la notificación en la figura \ref{fig:manual-android-notification}. Si pulsa en
la notificación será llevado a los detalles del chequeo para obtener más
detalles sobre el origen del aviso.

\begin{figure}[hbtp]
  \centering
  \includegraphics[width=0.4\textwidth]{5_diseno/android-5}
  \caption{Aspecto de las notificaciones}
  \label{fig:manual-android-notification}
\end{figure}



%%% Local Variables:
%%% mode: latex
%%% TeX-master: "../memoria"
%%% End:


\chapter{Diseño}
\label{chap:diseno}
En este capítulo se explicarán en detalle las instrucciones de uso de la
plataforma web.

\section{Introducción a SiteUp}

\textbf{SiteUp} es un servicio web creado para la monitorización de servicios de
Internet, especialmente pensado para dueños de sitios web y servidores que
deseen estar informados en caso de cualquier eventualidad.

El uso de internet en todos los ámbitos de la vida, desde las interacciones
sociales hasta la búsqueda de empleo, han dado lugar a la creación de servicios
accesibles a través de Internet y a nuevos modelos de negocio basados en
aquellos. Resulta por tanto imprescindible contar con sistemas para monitorizar
la disponibilidad de estos servicios y, en caso necesario, actuar de manera que
puedan subsanarse las causas de los problemas detectados.

SiteUp cubre precisamente estas necesidades, permitiendo a los usuarios:

\begin{itemize}
\item Crear chequeos mediante envío de paquetes ICMP (Ping), para revisar que un
  servidor está activo y conectado a la red.
\item Crear chequeos de registros DNS, de forma que el contenido de estos
  registros sea el correcto en todo momento.
\item Crear chequeos de puertos remotos, para vigilar que los servicios
  concretos dentro de un servidor son accesibles.
\item Crear chequeos mediante peticiones HTTP, para controlar que un servicio
  web responde de forma correcta y con el contenido adecuado.
\item Establecer las opciones de periodicidad de estos chequeos, de acuerdo a
  las necesidades particulares de cada caso.
\item Recibir notificaciones tanto a través de correo electrónico como mediante
  una aplicación móvil para el sistema operativo Android.
\end{itemize}

\section{Requisitos previos}

Para el uso correcto de la plataforma web SiteUp, el principal requisito es
contar con un dispositivo con acceso a internet que disponga de un navegador
reciente, como por ejemplo Google Chrome. Cualquier tipo de tablet, teléfono
móvil u ordenador personal podrá acceder a la plataforma web SiteUp, pero para
aprovechar mejor el potencial del sistema se recomienda usar una computadora
personal con las siguientes características:

\begin{itemize}
\item Navegador reciente. Se recomienda el uso de Google Chrome.
\item Pantalla con resolución mínima de 1280 píxeles en su lado longitudinal.
\item Conexión a internet con un caudal de descarga de al menos 512kb/s.
\end{itemize}

Si desea hacer uso de la aplicación móvil, será necesario que el usuario cuente
con un dispositivo con las siguientes características mínimas:

\begin{itemize}
\item Sistema operativo Android, versión mínima 2.3, versión recomendada 4.3.
\item Pantalla de al menos 3 pulgadas de diagonal y resolución mínima de 320x480
  píxeles.
\end{itemize}

\section{Uso de la plataforma web}

\subsection{Acceso}

A fecha de abril de 2014, la dirección de acceso a la plataforma web de SiteUp
es \url{http://siteup.josetomastocino.com}. Simplemente introduzca esa dirección
en su navegador y automáticamente será redirigido a la plataforma web, donde
verá la pantalla de bienvenida.

\subsection{Registro de usuario}

\begin{figure}[hbtp]
  \centering
  \includegraphics[width=0.3\textwidth]{apendice_manual_usuario/general_botones.png}
  \caption{Botones principales para el usuario anónimo}
  \label{fig:botones-1}
\end{figure}

Si es la primera vez que accede al sitio web deberá registrarse. Para ello, debe
pulsar en el botón \textbf{Sign up}, que se ve en la figura \ref{fig:botones-1}.

Tras ello, aparecerá el formulario de registro de usuario, que se ve en la
figura \ref{fig:registro}

\begin{figure}[hbtp]
  \centering
  \includegraphics[width=0.5\textwidth]{apendice_manual_usuario/formulario_registro.png}
  \caption{Formulario de registro de usuario}
  \label{fig:registro}
\end{figure}

En el formulario deberá introducir los siguientes datos:

\begin{itemize}
\item \textbf{Nombre de usuario}. Debe tener 30 caracteres como máximo, y solo
  pueden usarse letras, números y guiones.
\item \textbf{Dirección de correo electrónico}. Asegúrese de que es una
  dirección de correo correcta.
\item \textbf{Contraseña}.
\end{itemize}

Tras introducir los datos, deberá pulsar el botón \textbf{Sign up} del
formulario. Si el registro es correcto, el sistema dará de alta su usuario y
mostrará la pantalla \ref{fig:registro-correcto}. Si no, volverá a aparecer el
formulario con información sobre los errores en los datos.

\begin{figure}[hbtp]
  \centering
  \includegraphics[width=\textwidth]{apendice_manual_usuario/pantalla_register_completo.png}
  \caption{Mensaje indicando el registro correcto}
  \label{fig:registro-correcto}
\end{figure}

\subsection{Inicio de sesión}

Una vez que haya registrado su usuario, deberá pulsar en el botón \textbf{Log
  in}, como aparece en la imagen \ref{fig:botones-1}. Esto le llevará al
formulario de inicio de sesión que aparece en la figura \ref{fig:login}

\begin{figure}[hbtp]
  \centering
  \includegraphics[width=0.5\textwidth]{apendice_manual_usuario/general_formulario_login.png}
  \caption{Formulario de inicio de sesión}
  \label{fig:login}
\end{figure}

Deberá introducir su nombre de usuario y la contraseña que escribió al
registrarse, y posteriormente pulsar en el botón \textbf{Log in} del
formulario. Si los datos son correctos, será llevado al panel principal de la
web, donde podrá gestionar sus chequeos. Si no son correctos, volverá a aparecer
el formulario, con información sobre los errores encontrados.

\subsection{Recuperación de contraseña}

Si se registró como usuario pero no recuerda su contraseña es posible
recuperarla.

Para ello, vaya al formulario de inicio de sesión como ser ha explicado en el
paso anterior. En ese formulario, en la zona inferior, verá un enlace titulado
\textbf{Forgot password?} que deberá pulsar. Después de hacerlo se mostrará el
formulario de recuperación de contraseña, ilustrado en la figura
\ref{fig:formulario-reset}.

\begin{figure}[hbtp]
  \centering
  \includegraphics[width=0.5\textwidth]{apendice_manual_usuario/general_formulario_reset.png}
  \caption{Formulario de recuperación de contraseña}
  \label{fig:formulario-reset}
\end{figure}

En este formulario deberá introducir la dirección de correo electrónico que
utilizó al registrarse, y pulsar el botón \textbf{Send reset code}. Si lo hace
correctamente, aparecerá la pantalla de confirmación que se ve en la figura
\ref{fig:pantalla-reset-enviado}.

\begin{figure}[hbtp]
  \centering
  \includegraphics[width=0.5\textwidth]{apendice_manual_usuario/pantalla_reset_enviado.png}
  \caption{Pantalla de confirmación de envío de código de reinicio}
  \label{fig:pantalla-reset-enviado}
\end{figure}

El sistema habrá enviado un enlace a su dirección de correo electrónico. Deberá
pulsar el enlace que aparece en ese correo. Asegúrese de que recibe
correctamente el correo. Si no lo recibe, revise la carpeta de correo no
deseado. El enlace del correo le llevará a un formulario de reseteo de
contraseña como el que aparece en la figura \ref{fig:formulario-reset-2}.

\begin{figure}[hbtp]
  \centering
  \includegraphics[width=0.5\textwidth]{apendice_manual_usuario/general_formulario_reset_2.png}
  \caption{}
  \label{fig:formulario-reset-2}
\end{figure}

En este formulario deberá escribir la nueva contraseña de su cuenta dos veces,
para confirmar que la escribe correctamente. Tras ello, deberá pulsar el botón
\textbf{Set new password}, lo cual confirmará los cambios. 

Tras este proceso, podrá iniciar sesión como se ha descrito previamente, esta
vez utilizando la nueva contraseña.


\subsection{Pantalla de listado de chequeos}

\begin{figure}[hbtp]
  \centering
  \includegraphics[width=0.8\textwidth]{apendice_manual_usuario/pantalla-dashboard.png}
  \caption{Panel de control}
  \label{fig:dashboard}
\end{figure}

\begin{figure}[hbtp]
  \centering
  \includegraphics[width=0.4\textwidth]{apendice_manual_usuario/general_botones_2.png}
  \caption{Botones principales para el usuario logueado}
  \label{fig:botones-2}
\end{figure}

En la figura \ref{fig:dashboard} se puede ver el aspecto del panel de control de
la plataforma web SiteUp. Desde este panel podrá hacer una gran cantidad de
acciones, desde crear grupos de chequeos hasta editar sus detalles de usuario.

Es posible acceder a esta pantalla desde cualquier sitio de la web si está
logueado, pulsando en el botón \textbf{Dashboard} de la cabecera de la web, tal
y como se ve en la figura \ref{fig:botones-2}.

\subsection{Cerrar sesión}

Si ha iniciado sesión y desea cerrarla, debe dirigirse a los botones
principales, situados en el lado derecho de la cabecera de la web, tal y como
aparecen en la figura \ref{fig:botones-2}. Deberá pulsar el botón \textbf{Log
  out}, tras lo cual se cerrara su sesión y será redirigido a la pantalla
inicial de SiteUp.

\subsection{Edición de datos de usuario}

Si desea editar alguno de los datos de usuario que ingresó cuando hizo el
registro, deberá primero pulsar el botón \textbf{Your profile}, situado en la
cabecera de la web. Es el segundo de los botones que aparecen, tal y como se ve
en la figura \ref{fig:botones-2}.

Una vez hecho esto, aparecerá el formulario de edición de usuario, ilustrado en
la figura \ref{fig:formulario-usuario}, en el que verá sus actuales datos de
usuario. 

\begin{figure}[hbtp]
  \centering
  \includegraphics[width=0.8\textwidth]{apendice_manual_usuario/pantalla_perfil_usuario.png}
  \caption{}
  \label{fig:formulario-usuario}
\end{figure}

En este formulario podrá editar:

\begin{itemize}
\item Su \textbf{nombre} de usuario.
\item Su \textbf{dirección} de correo electrónico.
\item Si desea o no recibir \textbf{resúmenes diarios} sobre el estado de sus chequeos.
\end{itemize}

Una vez editados los datos que desee deberá pulsar el botón \textbf{Update
  details}. También podrá editar su contraseña pulsando el botón \textbf{Change
  password}, situado al pie del formulario a la izquierda.


\subsection{Creación de grupo de chequeos}

En SiteUp, los chequeos se organizan en \textbf{grupos de chequeos}, lo que
permite trabajar con ellos más fácilmente a la hora de activarlos, desactivarlos
o borrarlos. 

Para crear un grupo de chequeos deberá pulsar el botón \textbf{Add group},
situado en la zona superior derecha del panel principal. El botón puede verse en
la figura \ref{fig:crear-grupo}.

\begin{figure}[hbtp]
  \centering
  \includegraphics[width=0.3\textwidth]{apendice_manual_usuario/grupo_boton_nuevo.png}
  \caption{Botón para crear nuevo grupo de chequeos}
  \label{fig:crear-grupo}
\end{figure}

Tras pulsar el botón aparecerá el formulario de creación de nuevo grupo. En él
deberá introducir los datos del nuevo grupo. En particular:

\begin{itemize}
\item El \textbf{título} que identifica al grupo.
\end{itemize}

Tras ello, deberá pulsar el botón \textbf{Create group}, y el grupo será creado

\begin{figure}[H]
  \centering
  \includegraphics[width=0.8\textwidth]{apendice_manual_usuario/grupo_formulario_crear.png}
  \caption{Formulario de creación de grupo de chequeos}
  \label{fig:}
\end{figure}


\subsection{Operaciones con grupos de chequeos}

En la figura~\ref{fig:operaciones-grupo} se puede ver el menú de operaciones
disponibles para un grupo de chequeos. Este menú aparece en el panel principal,
al lado derecho junto al título de cada grupo de chequeos.

\begin{figure}[hbtp]
  \centering
  \includegraphics[width=0.5\textwidth]{apendice_manual_usuario/grupo_submenu_acciones.png}
  \caption{}
  \label{fig:operaciones-grupo}
\end{figure}

\subsubsection{Edición de grupo de chequeos}

En el menú de operaciones del grupo de chequeos, pulse el botón
\textbf{Edit}. Aparecerá un formulario similar al de creación de grupo, en el
que podrá editar el título del grupo de chequeos.

\subsubsection{Borrado de grupo de chequeos}

Para borrar un grupo de chequeos, pulse el botón \textbf{Delete} en el menú de
operaciones del grupo. Aparecerá un formulario como el de la figura
\ref{fig:formulario-borrado-grupo}, pidiéndole confirmación del borrado. 

\begin{figure}[hbtp]
  \centering
  \includegraphics[width=0.5\textwidth]{apendice_manual_usuario/formulario_borrado_grupo}
  \caption{Formulario de confirmación de borrado}
  \label{fig:formulario-borrado-grupo}
\end{figure}

Si pulsa en el botón \textbf{Go back}, volverá al panel principal, cancelando el
borrado. Si pulsa el botón \textbf{Delete}, confirmará la operación y el sistema
borrará el grupo de chequeos, así como todos los chequeos que contenga.

\subsubsection{Activación y desactivación de grupo de chequeos}

Es posible activar y desactivar de forma grupal todos los chequeos de un
grupo. Para ello, deberá dirigirse al menú de operaciones del grupo y pulsar en
\textbf{Enable all checks} o en \textbf{Disable all checks}, según sea su decisión.

\subsection{Creación de chequeo}

\begin{figure}[hbtp]
  \centering
  \includegraphics[width=0.8\textwidth]{apendice_manual_usuario/pantalla_seleccion_tipo.png}
  \caption{Pantalla de selección del tipo de chequeo a crear}
  \label{fig:seleccionar-tipo-check}
\end{figure}

Como se ha comentado previamente, el elemento principal de SiteUp es el
\textbf{chequeo}. El usuario da de alta chequeos dentro de un grupo, y estos
chequeos se lanzan periódicamente, generando datos.

Para crear un chequeo dentro de un grupo deberá ir al menú de operaciones del
grupo, que aparece en la figura \ref{fig:operaciones-grupo} y se sitúa a la
derecha del título de un grupo de chequeos. Dentro del menú, deberá pulsar la
opción \textbf{Add check}.



Tras esto aparecerá una pantalla en la que deberá seleccionar el tipo de chequeo
que desea crear. La pantalla es similar a la de la figura
\ref{fig:seleccionar-tipo-check}. Como se explicó en la introducción de este
manual, es posible crear cuatro tipo de chequeos: chequeos de ping, chequeos de
puerto, chequeos de registros DNS y chequeos por HTTP.

Si pulsa en el botón del tipo que desea crear, aparecerá un formulario con
campos para rellenar.

\subsubsection{Creación de un chequeo por ping}

\begin{figure}[hbtp]
  \centering
  \includegraphics[width=0.8\textwidth]{apendice_manual_usuario/pantalla_crear_ping.png}
  \caption{Formulario de creación de chequeo por ping}
  \label{fig:crear-ping}
\end{figure}

El formulario para crear un chequeo de tipo ping aparece en la
figura~\ref{fig:crear-ping}. Los campos que deberá rellenar son los siguientes:

\begin{itemize}
\item \textbf{Target}: es el objetivo a chequear. Debe ser un \textit{hostname} o una IP.
\item \textbf{Título} del chequeo.
\item \textbf{Descripción} opcional del chequeo.
\item \textbf{Check interval}: intervalo entre chequeos, definido en minutos. El
  valor mínimo es de 1 minuto.
\item \textbf{Notify email}: indica si los cambios de estado en este chequeo
  deberán ser notificados por correo electrónico.
\item \textbf{Notify Android}: indica si los cambios de estado en este chequeo
  deberán ser notificados mediante la aplicación Android.
\item \textbf{Should check timeout}: indica si hay que verificar que el ping
  enviado vuelve dentro de un tiempo máximo determinado.
\item \textbf{Timeout value}: si la anterior opción está activa, este campo
  define el número de milisegundos que puede tardar en responder el servidor al
  ping, como máximo.
\end{itemize}

\subsubsection{Creación de un chequeo de puerto}

\begin{figure}[hbtp]
  \centering
  \includegraphics[width=0.8\textwidth]{apendice_manual_usuario/pantalla_crear_port.png}
  \caption{Formulario de creación de chequeo de puerto}
  \label{fig:crear-port}
\end{figure}

El formulario para crear un chequeo de puertos aparece en la
figura~\ref{fig:crear-port}. Los campos que deberá rellenar son los siguientes:

\begin{itemize}
\item \textbf{Título} del chequeo.
\item \textbf{Descripción} opcional del chequeo.
\item \textbf{Target}: es el objetivo a chequear. Debe ser un \textit{hostname} o una IP.
\item \textbf{Target port}: es el puerto a chequear. Debe ser un entero.
\item \textbf{Check for string}: opcionalmente, se puede indicar una cadena que
  debe estar presente en la respuesta al conectar con el puerto.
\item \textbf{Check interval}: intervalo entre chequeos, definido en minutos. El
  valor mínimo es de 1 minuto.
\item \textbf{Notify email}: indica si los cambios de estado en este chequeo
  deberán ser notificados por correo electrónico.
\item \textbf{Notify Android}: indica si los cambios de estado en este chequeo
  deberán ser notificados mediante la aplicación Android.
\end{itemize}

\subsubsection{Creación de un chequeo de DNS}

\begin{figure}[hbtp]
  \centering
  \includegraphics[width=0.8\textwidth]{apendice_manual_usuario/pantalla_crear_dns.png}
  \caption{Formulario de creación de chequeo de registro DNS}
  \label{fig:crear-dns}
\end{figure}

El formulario para crear un chequeo de registro DNS aparece en la figura
\ref{fig:crear-dns}. Los campos que deberá rellenar son los siguientes:

\begin{itemize}
\item \textbf{Título} del chequeo.
\item \textbf{Descripción} opcional del chequeo.
\item \textbf{Target}: es el objetivo a chequear. Debe ser un \textit{hostname}.
\item \textbf{Record type}: indica el tipo de registro a chequear. Puede ser A, AAAA, CNAME, TXT y MX.
\item \textbf{Resolved address}: indica el contenido que debe tener el registro.
\item \textbf{Check interval}: intervalo entre chequeos, definido en minutos. El
  valor mínimo es de 1 minuto.
\item \textbf{Notify email}: indica si los cambios de estado en este chequeo
  deberán ser notificados por correo electrónico.
\item \textbf{Notify Android}: indica si los cambios de estado en este chequeo
  deberán ser notificados mediante la aplicación Android.
\end{itemize}

\subsubsection{Creación de un chequeo por HTTP}

\begin{figure}[hbtp]
  \centering
  \includegraphics[width=0.8\textwidth]{apendice_manual_usuario/pantalla_crear_http.png}
  \caption{Formulario de creación de chequeo por HTTP}
  \label{fig:crear-http}
\end{figure}

El formulario para crear un chequeo por HTTP aparece en la figura
\ref{fig:crear-http}. Los campos que deberá rellenar son los siguientes:

\begin{itemize}
\item \textbf{Título} del chequeo.
\item \textbf{Descripción} opcional del chequeo.
\item \textbf{Target}: es el objetivo a chequear. Debe ser una URL.
\item \textbf{Status code}: es el código de estado que deberá devolver la petición..
\item \textbf{Check for string}: opcionalmente, se puede indicar una cadena que
  debe estar presente en la respuesta a la petición.
\item \textbf{Check interval}: intervalo entre chequeos, definido en minutos. El
  valor mínimo es de 1 minuto.
\item \textbf{Notify email}: indica si los cambios de estado en este chequeo
  deberán ser notificados por correo electrónico.
\item \textbf{Notify Android}: indica si los cambios de estado en este chequeo
  deberán ser notificados mediante la aplicación Android.
\end{itemize}


\subsection{Operaciones con chequeos}

\begin{figure}[hbtp]
  \centering
  \includegraphics[width=\textwidth]{apendice_manual_usuario/chequeo_valido.png}
  \caption{}
  \label{fig:chequeo}
\end{figure}

Una vez creado un chequeo, aparecerá en el panel de control con el aspecto de la
figura~\ref{fig:chequeo}. Como se aprecia en la figura, se muestra mediante un
icono el estado del chequeo, los datos (título, tipo, descripción) y una gráfica
que ilustra el estado del chequeo en las últimas 24 horas.

\begin{figure}[hbtp]
  \centering
  \includegraphics[width=\textwidth]{apendice_manual_usuario/chequeo_invalido.png}
  \caption{}
  \label{fig:chequeo-invalido}
\end{figure}

Si un chequeo detecta un error, cambia de apariencia, pasando a tener el aspecto
que ofrece la figura \ref{fig:chequeo-invalido}. Además, según se haya
configurado en las opciones del chequeo, se envía una notificación al usuario
informando del cambio de estado del chequeo.

Cuando se ve en la lista general de chequeos, cada chequeo tiene una serie de
botones que permiten llevar a cabo operaciones diversas, que se detallan a
continuación.

\subsubsection{Activación y desactivación de un chequeo}

Para desactivar un chequeo, tiene que pulsar el botón \textbf{Disable} de la
botonera, que solo aparecerá si el chequeo está activado. De igual modo, si el
chequeo está desactivado, el botón que aparecerá tendrá la etiqueta
\textbf{Enable}, que podrá pulsar para activarlo.

\subsubsection{Borrado de un chequeo}

En la botonera se sitúa el botón \textbf{Delete}, que podrá pulsar para borrar
el chequeo. Aparecerá un formulario pidiéndole confirmación. Si confirma la
acción, se procederá al borrado del chequeo y todos los datos dependientes de
éste.

\subsubsection{Edición de un chequeo}

Si desea editar un chequeo debe pulsar el botón \textbf{Edit}, situado en la
botonera del chequeo. Aparecerá un formulario similar al de creación de chequeo,
pero con los datos actuales ya introducidos. Puede modificar los datos que desee
y confirmar los cambios pulsando el botón \textbf{Update check} del formulario.

\subsubsection{Vista de detalle del chequeo}

\begin{figure}[hbtp]
  \centering
  \includegraphics[width=\textwidth]{apendice_manual_usuario/pantalla_detalle_chequeo.png}
  \caption{}
  \label{fig:detalle-chequeo}
\end{figure}

En la botonera del chequeo, en el panel de control, verá un botón
\textbf{Details}. Si lo pulsa será llevado a la pantalla de vista de detalle del
chequeo, que se ilustra en la figura \ref{fig:detalle-chequeo}. Esta pantalla
ofrece mucha información sobre el chequeo. En particular:

\begin{itemize}
\item En la parte superior se sitúa la información general:
  \begin{itemize}
  \item Icono con el estado del chequeo.
  \item Título.
  \item Tipo de chequeo.
  \item Descripción.
  \item Botonera de acciones.
  \end{itemize}

\item En la parte media de la pantalla se sitúa un bloque con pestañas. Mediante
  estas pestañas es posible indicar sobre qué periodo de tiempo queremos
  visualizar los datos: últimas 24 horas o última semana. La información
  disponible es la siguiente:

  \begin{itemize}
  \item Porcentaje de \textit{uptime} del chequeo.
  \item \textit{(Solo chequeos tipo ping)} Tiempo de respuesta mínimo, máximo y promedio.
  \item Gráfica con el estado del chequeo.
  \end{itemize}

\item En la parte inferior de la pantalla se sitúa el panel de eventos del
  chequeo. Éste permite ver los últimos cambios de estado del chequeo:

  \begin{itemize}
  \item Tipo de cambio de estado (\textit{Up}, \textit{Down}).
  \item Duración del estado.
  \item Hora de comienzo del estado.
  \item Hora de fin del estado. Si es el estado en el que el chequeo se
    encuentra actualmente, no aparecerá ninguna fecha.
  \item Extra: muestra información adicional para estados de error.
  \end{itemize}
\end{itemize}

\begin{figure}[hbtp]
  \centering
  \includegraphics[width=\textwidth]{apendice_manual_usuario/pantalla_grafica_semanal.png}
  \caption{Detalle de la gráfica para un intervalo de tiempo de una semana}
  \label{fig:grafica-semanal}
\end{figure} 

\subsection{Recepción de notificaciones por correo electrónico}

\subsubsection{Notificaciones de cambio de estado}

Cuando el sistema detecte un cambio de estado en un chequeo, y esté activada la
notificación por correo electrónico, se enviará un email como el que aparece en
la figura \ref{fig:email-estado}.

\begin{figure}[hbtp]
  \centering
  \includegraphics[width=0.7\textwidth]{apendice_manual_usuario/email_estado.png}
  \caption{Email de notificación de cambio de estado}
  \label{fig:email-estado}
\end{figure}

En ese email se le indicará el título del chequeo, el nuevo estado, la fecha y
hora en la que ha pasado al nuevo estado, y un enlace directo al panel de
detalle del chequeo en SiteUp, de forma que pueda obtener más información.

\subsubsection{Notificaciones de resumen diarias}

Si ha activado el envío de resúmenes diarios sobre el estado de sus chequeos, el
sistema enviará todos los días un correo electrónico con información. El aspecto
del correo que recibirá se puede ver en la figura \ref{fig:email-resumen}

\begin{figure}[hbtp]
  \centering
  \includegraphics[width=0.7\textwidth]{apendice_manual_usuario/email_resumen.png}
  \caption{Email diario de resumen}
  \label{fig:email-resumen}
\end{figure}


\section{Uso de la aplicación móvil}

La aplicación móvil \textbf{SiteUp Client} es una aplicación para para el
sistema operativo Android que permite a sus usuarios revisar los chequeos que
tiene dados de alta en la plataforma web SiteUp, así como recibir de manera
instantánea notificaciones \textit{push} sobre el estado de sus chequeos.

\subsection{Instalación}

La aplicación móvil se instala copiando y ejecutando el fichero APK en el
dispositivo. En un futuro se podrá encontrar en la \textit{Play Store} de
Google.

\subsection{Lanzamiento}

En la figura~\ref{fig:manual-android-launcher} se puede observar el aspecto de
la aplicación \textbf{SiteUp Client} en el lanzador de Android. Para ejecutarla,
simplemente pulse sobre el icono de la aplicación.

\begin{figure}[htbp]
  \centering
  \includegraphics[width=0.4\textwidth]{5_diseno/android-1}
  \caption{Icono de SiteUp en el lanzador}
  \label{fig:manual-android-launcher}
\end{figure}

\subsection{Carga de la aplicación}

Tras lanzar la aplicación, el sistema cargará los datos necesarios y se
comunicará con los servidores de Google para obtener los datos necesarios. La
apariencia de la aplicación durante este proceso se puede ver en la figura
\ref{fig:manual-android-login}.

\begin{figure}[htbp]
  \centering
  \includegraphics[width=0.4\textwidth]{5_diseno/android-2}
  \caption{Pantalla de carga de la aplicación móvil}
  \label{fig:manual-android-loading}
\end{figure}

\subsection{Inicio de sesión}

Tras la carga, la primera vez que lance la aplicación será necesario que inicie
sesión con las credenciales que utilizó en la plataforma web. Para ello, tendrá
que introducir su nombre de usuario y contraseña en el formulario que aparece en
la figura~\ref{fig:manual-android-loading}, y pulsar el botón \textbf{Sign in}.

\begin{figure}[htbp]
  \centering
  \includegraphics[width=0.4\textwidth]{5_diseno/android-3}
  \caption{Pantalla de inicio de sesión de la aplicación móvil}
  \label{fig:manual-android-login}
\end{figure}

La aplicación móvil se comunicará con los servidores centrales de SiteUp,
validará sus credenciales de usuario y, si son correctas, le dará acceso al
sistema. Si las credenciales de usuario no son correctas, el formulario mostrará
un error.

\subsection{Revisión de chequeos}

Si ha iniciado sesión correctamente, tendrá la opción de revisar la lista de
chequeos que tiene dados de alta en la plataforma web \textbf{SiteUp}. La
apariencia de la lista de chequeos es igual a la de la
figura~\ref{fig:manual-android-dashboard}.

\begin{figure}[hbtp]
  \centering
  \includegraphics[width=0.4\textwidth]{5_diseno/android-4}
  \caption{Pantalla de chequeos de la aplicación móvil}
  \label{fig:manual-android-dashboard}
\end{figure}

Si pulsa en cualquiera de los chequeos, será redirigido a la plataforma web, que
podrá navegar sin problema gracias a que se encuentra preparada para su uso en
dispositivos móviles.

\subsection{Recepción de notificaciones}

La utilidad principal de la aplicación es la posibilidad de recibir
notificaciones en cualquier momento provenientes del servidor -- conocidas como
notificaciones \textit{push}. 

El procedimiento es transparente al usuario. Simplemente lance la aplicación e
inice sesión. Tras hacer esto, la aplicación estará preparada para recibir
notificaciones desde el servidor.

Cuando se reciba una notificación, la aplicación mostará un mensaje en el área
de notificaciones del móvil, con información resumida. Puede ver el aspecto de
la notificación en la figura \ref{fig:manual-android-notification}. Si pulsa en
la notificación será llevado a los detalles del chequeo para obtener más
detalles sobre el origen del aviso.

\begin{figure}[hbtp]
  \centering
  \includegraphics[width=0.4\textwidth]{5_diseno/android-5}
  \caption{Aspecto de las notificaciones}
  \label{fig:manual-android-notification}
\end{figure}



%%% Local Variables:
%%% mode: latex
%%% TeX-master: "../memoria"
%%% End:


\chapter{Implementación}
\label{chap:implementacion}
En este capítulo se explicarán en detalle las instrucciones de uso de la
plataforma web.

\section{Introducción a SiteUp}

\textbf{SiteUp} es un servicio web creado para la monitorización de servicios de
Internet, especialmente pensado para dueños de sitios web y servidores que
deseen estar informados en caso de cualquier eventualidad.

El uso de internet en todos los ámbitos de la vida, desde las interacciones
sociales hasta la búsqueda de empleo, han dado lugar a la creación de servicios
accesibles a través de Internet y a nuevos modelos de negocio basados en
aquellos. Resulta por tanto imprescindible contar con sistemas para monitorizar
la disponibilidad de estos servicios y, en caso necesario, actuar de manera que
puedan subsanarse las causas de los problemas detectados.

SiteUp cubre precisamente estas necesidades, permitiendo a los usuarios:

\begin{itemize}
\item Crear chequeos mediante envío de paquetes ICMP (Ping), para revisar que un
  servidor está activo y conectado a la red.
\item Crear chequeos de registros DNS, de forma que el contenido de estos
  registros sea el correcto en todo momento.
\item Crear chequeos de puertos remotos, para vigilar que los servicios
  concretos dentro de un servidor son accesibles.
\item Crear chequeos mediante peticiones HTTP, para controlar que un servicio
  web responde de forma correcta y con el contenido adecuado.
\item Establecer las opciones de periodicidad de estos chequeos, de acuerdo a
  las necesidades particulares de cada caso.
\item Recibir notificaciones tanto a través de correo electrónico como mediante
  una aplicación móvil para el sistema operativo Android.
\end{itemize}

\section{Requisitos previos}

Para el uso correcto de la plataforma web SiteUp, el principal requisito es
contar con un dispositivo con acceso a internet que disponga de un navegador
reciente, como por ejemplo Google Chrome. Cualquier tipo de tablet, teléfono
móvil u ordenador personal podrá acceder a la plataforma web SiteUp, pero para
aprovechar mejor el potencial del sistema se recomienda usar una computadora
personal con las siguientes características:

\begin{itemize}
\item Navegador reciente. Se recomienda el uso de Google Chrome.
\item Pantalla con resolución mínima de 1280 píxeles en su lado longitudinal.
\item Conexión a internet con un caudal de descarga de al menos 512kb/s.
\end{itemize}

Si desea hacer uso de la aplicación móvil, será necesario que el usuario cuente
con un dispositivo con las siguientes características mínimas:

\begin{itemize}
\item Sistema operativo Android, versión mínima 2.3, versión recomendada 4.3.
\item Pantalla de al menos 3 pulgadas de diagonal y resolución mínima de 320x480
  píxeles.
\end{itemize}

\section{Uso de la plataforma web}

\subsection{Acceso}

A fecha de abril de 2014, la dirección de acceso a la plataforma web de SiteUp
es \url{http://siteup.josetomastocino.com}. Simplemente introduzca esa dirección
en su navegador y automáticamente será redirigido a la plataforma web, donde
verá la pantalla de bienvenida.

\subsection{Registro de usuario}

\begin{figure}[hbtp]
  \centering
  \includegraphics[width=0.3\textwidth]{apendice_manual_usuario/general_botones.png}
  \caption{Botones principales para el usuario anónimo}
  \label{fig:botones-1}
\end{figure}

Si es la primera vez que accede al sitio web deberá registrarse. Para ello, debe
pulsar en el botón \textbf{Sign up}, que se ve en la figura \ref{fig:botones-1}.

Tras ello, aparecerá el formulario de registro de usuario, que se ve en la
figura \ref{fig:registro}

\begin{figure}[hbtp]
  \centering
  \includegraphics[width=0.5\textwidth]{apendice_manual_usuario/formulario_registro.png}
  \caption{Formulario de registro de usuario}
  \label{fig:registro}
\end{figure}

En el formulario deberá introducir los siguientes datos:

\begin{itemize}
\item \textbf{Nombre de usuario}. Debe tener 30 caracteres como máximo, y solo
  pueden usarse letras, números y guiones.
\item \textbf{Dirección de correo electrónico}. Asegúrese de que es una
  dirección de correo correcta.
\item \textbf{Contraseña}.
\end{itemize}

Tras introducir los datos, deberá pulsar el botón \textbf{Sign up} del
formulario. Si el registro es correcto, el sistema dará de alta su usuario y
mostrará la pantalla \ref{fig:registro-correcto}. Si no, volverá a aparecer el
formulario con información sobre los errores en los datos.

\begin{figure}[hbtp]
  \centering
  \includegraphics[width=\textwidth]{apendice_manual_usuario/pantalla_register_completo.png}
  \caption{Mensaje indicando el registro correcto}
  \label{fig:registro-correcto}
\end{figure}

\subsection{Inicio de sesión}

Una vez que haya registrado su usuario, deberá pulsar en el botón \textbf{Log
  in}, como aparece en la imagen \ref{fig:botones-1}. Esto le llevará al
formulario de inicio de sesión que aparece en la figura \ref{fig:login}

\begin{figure}[hbtp]
  \centering
  \includegraphics[width=0.5\textwidth]{apendice_manual_usuario/general_formulario_login.png}
  \caption{Formulario de inicio de sesión}
  \label{fig:login}
\end{figure}

Deberá introducir su nombre de usuario y la contraseña que escribió al
registrarse, y posteriormente pulsar en el botón \textbf{Log in} del
formulario. Si los datos son correctos, será llevado al panel principal de la
web, donde podrá gestionar sus chequeos. Si no son correctos, volverá a aparecer
el formulario, con información sobre los errores encontrados.

\subsection{Recuperación de contraseña}

Si se registró como usuario pero no recuerda su contraseña es posible
recuperarla.

Para ello, vaya al formulario de inicio de sesión como ser ha explicado en el
paso anterior. En ese formulario, en la zona inferior, verá un enlace titulado
\textbf{Forgot password?} que deberá pulsar. Después de hacerlo se mostrará el
formulario de recuperación de contraseña, ilustrado en la figura
\ref{fig:formulario-reset}.

\begin{figure}[hbtp]
  \centering
  \includegraphics[width=0.5\textwidth]{apendice_manual_usuario/general_formulario_reset.png}
  \caption{Formulario de recuperación de contraseña}
  \label{fig:formulario-reset}
\end{figure}

En este formulario deberá introducir la dirección de correo electrónico que
utilizó al registrarse, y pulsar el botón \textbf{Send reset code}. Si lo hace
correctamente, aparecerá la pantalla de confirmación que se ve en la figura
\ref{fig:pantalla-reset-enviado}.

\begin{figure}[hbtp]
  \centering
  \includegraphics[width=0.5\textwidth]{apendice_manual_usuario/pantalla_reset_enviado.png}
  \caption{Pantalla de confirmación de envío de código de reinicio}
  \label{fig:pantalla-reset-enviado}
\end{figure}

El sistema habrá enviado un enlace a su dirección de correo electrónico. Deberá
pulsar el enlace que aparece en ese correo. Asegúrese de que recibe
correctamente el correo. Si no lo recibe, revise la carpeta de correo no
deseado. El enlace del correo le llevará a un formulario de reseteo de
contraseña como el que aparece en la figura \ref{fig:formulario-reset-2}.

\begin{figure}[hbtp]
  \centering
  \includegraphics[width=0.5\textwidth]{apendice_manual_usuario/general_formulario_reset_2.png}
  \caption{}
  \label{fig:formulario-reset-2}
\end{figure}

En este formulario deberá escribir la nueva contraseña de su cuenta dos veces,
para confirmar que la escribe correctamente. Tras ello, deberá pulsar el botón
\textbf{Set new password}, lo cual confirmará los cambios. 

Tras este proceso, podrá iniciar sesión como se ha descrito previamente, esta
vez utilizando la nueva contraseña.


\subsection{Pantalla de listado de chequeos}

\begin{figure}[hbtp]
  \centering
  \includegraphics[width=0.8\textwidth]{apendice_manual_usuario/pantalla-dashboard.png}
  \caption{Panel de control}
  \label{fig:dashboard}
\end{figure}

\begin{figure}[hbtp]
  \centering
  \includegraphics[width=0.4\textwidth]{apendice_manual_usuario/general_botones_2.png}
  \caption{Botones principales para el usuario logueado}
  \label{fig:botones-2}
\end{figure}

En la figura \ref{fig:dashboard} se puede ver el aspecto del panel de control de
la plataforma web SiteUp. Desde este panel podrá hacer una gran cantidad de
acciones, desde crear grupos de chequeos hasta editar sus detalles de usuario.

Es posible acceder a esta pantalla desde cualquier sitio de la web si está
logueado, pulsando en el botón \textbf{Dashboard} de la cabecera de la web, tal
y como se ve en la figura \ref{fig:botones-2}.

\subsection{Cerrar sesión}

Si ha iniciado sesión y desea cerrarla, debe dirigirse a los botones
principales, situados en el lado derecho de la cabecera de la web, tal y como
aparecen en la figura \ref{fig:botones-2}. Deberá pulsar el botón \textbf{Log
  out}, tras lo cual se cerrara su sesión y será redirigido a la pantalla
inicial de SiteUp.

\subsection{Edición de datos de usuario}

Si desea editar alguno de los datos de usuario que ingresó cuando hizo el
registro, deberá primero pulsar el botón \textbf{Your profile}, situado en la
cabecera de la web. Es el segundo de los botones que aparecen, tal y como se ve
en la figura \ref{fig:botones-2}.

Una vez hecho esto, aparecerá el formulario de edición de usuario, ilustrado en
la figura \ref{fig:formulario-usuario}, en el que verá sus actuales datos de
usuario. 

\begin{figure}[hbtp]
  \centering
  \includegraphics[width=0.8\textwidth]{apendice_manual_usuario/pantalla_perfil_usuario.png}
  \caption{}
  \label{fig:formulario-usuario}
\end{figure}

En este formulario podrá editar:

\begin{itemize}
\item Su \textbf{nombre} de usuario.
\item Su \textbf{dirección} de correo electrónico.
\item Si desea o no recibir \textbf{resúmenes diarios} sobre el estado de sus chequeos.
\end{itemize}

Una vez editados los datos que desee deberá pulsar el botón \textbf{Update
  details}. También podrá editar su contraseña pulsando el botón \textbf{Change
  password}, situado al pie del formulario a la izquierda.


\subsection{Creación de grupo de chequeos}

En SiteUp, los chequeos se organizan en \textbf{grupos de chequeos}, lo que
permite trabajar con ellos más fácilmente a la hora de activarlos, desactivarlos
o borrarlos. 

Para crear un grupo de chequeos deberá pulsar el botón \textbf{Add group},
situado en la zona superior derecha del panel principal. El botón puede verse en
la figura \ref{fig:crear-grupo}.

\begin{figure}[hbtp]
  \centering
  \includegraphics[width=0.3\textwidth]{apendice_manual_usuario/grupo_boton_nuevo.png}
  \caption{Botón para crear nuevo grupo de chequeos}
  \label{fig:crear-grupo}
\end{figure}

Tras pulsar el botón aparecerá el formulario de creación de nuevo grupo. En él
deberá introducir los datos del nuevo grupo. En particular:

\begin{itemize}
\item El \textbf{título} que identifica al grupo.
\end{itemize}

Tras ello, deberá pulsar el botón \textbf{Create group}, y el grupo será creado

\begin{figure}[H]
  \centering
  \includegraphics[width=0.8\textwidth]{apendice_manual_usuario/grupo_formulario_crear.png}
  \caption{Formulario de creación de grupo de chequeos}
  \label{fig:}
\end{figure}


\subsection{Operaciones con grupos de chequeos}

En la figura~\ref{fig:operaciones-grupo} se puede ver el menú de operaciones
disponibles para un grupo de chequeos. Este menú aparece en el panel principal,
al lado derecho junto al título de cada grupo de chequeos.

\begin{figure}[hbtp]
  \centering
  \includegraphics[width=0.5\textwidth]{apendice_manual_usuario/grupo_submenu_acciones.png}
  \caption{}
  \label{fig:operaciones-grupo}
\end{figure}

\subsubsection{Edición de grupo de chequeos}

En el menú de operaciones del grupo de chequeos, pulse el botón
\textbf{Edit}. Aparecerá un formulario similar al de creación de grupo, en el
que podrá editar el título del grupo de chequeos.

\subsubsection{Borrado de grupo de chequeos}

Para borrar un grupo de chequeos, pulse el botón \textbf{Delete} en el menú de
operaciones del grupo. Aparecerá un formulario como el de la figura
\ref{fig:formulario-borrado-grupo}, pidiéndole confirmación del borrado. 

\begin{figure}[hbtp]
  \centering
  \includegraphics[width=0.5\textwidth]{apendice_manual_usuario/formulario_borrado_grupo}
  \caption{Formulario de confirmación de borrado}
  \label{fig:formulario-borrado-grupo}
\end{figure}

Si pulsa en el botón \textbf{Go back}, volverá al panel principal, cancelando el
borrado. Si pulsa el botón \textbf{Delete}, confirmará la operación y el sistema
borrará el grupo de chequeos, así como todos los chequeos que contenga.

\subsubsection{Activación y desactivación de grupo de chequeos}

Es posible activar y desactivar de forma grupal todos los chequeos de un
grupo. Para ello, deberá dirigirse al menú de operaciones del grupo y pulsar en
\textbf{Enable all checks} o en \textbf{Disable all checks}, según sea su decisión.

\subsection{Creación de chequeo}

\begin{figure}[hbtp]
  \centering
  \includegraphics[width=0.8\textwidth]{apendice_manual_usuario/pantalla_seleccion_tipo.png}
  \caption{Pantalla de selección del tipo de chequeo a crear}
  \label{fig:seleccionar-tipo-check}
\end{figure}

Como se ha comentado previamente, el elemento principal de SiteUp es el
\textbf{chequeo}. El usuario da de alta chequeos dentro de un grupo, y estos
chequeos se lanzan periódicamente, generando datos.

Para crear un chequeo dentro de un grupo deberá ir al menú de operaciones del
grupo, que aparece en la figura \ref{fig:operaciones-grupo} y se sitúa a la
derecha del título de un grupo de chequeos. Dentro del menú, deberá pulsar la
opción \textbf{Add check}.



Tras esto aparecerá una pantalla en la que deberá seleccionar el tipo de chequeo
que desea crear. La pantalla es similar a la de la figura
\ref{fig:seleccionar-tipo-check}. Como se explicó en la introducción de este
manual, es posible crear cuatro tipo de chequeos: chequeos de ping, chequeos de
puerto, chequeos de registros DNS y chequeos por HTTP.

Si pulsa en el botón del tipo que desea crear, aparecerá un formulario con
campos para rellenar.

\subsubsection{Creación de un chequeo por ping}

\begin{figure}[hbtp]
  \centering
  \includegraphics[width=0.8\textwidth]{apendice_manual_usuario/pantalla_crear_ping.png}
  \caption{Formulario de creación de chequeo por ping}
  \label{fig:crear-ping}
\end{figure}

El formulario para crear un chequeo de tipo ping aparece en la
figura~\ref{fig:crear-ping}. Los campos que deberá rellenar son los siguientes:

\begin{itemize}
\item \textbf{Target}: es el objetivo a chequear. Debe ser un \textit{hostname} o una IP.
\item \textbf{Título} del chequeo.
\item \textbf{Descripción} opcional del chequeo.
\item \textbf{Check interval}: intervalo entre chequeos, definido en minutos. El
  valor mínimo es de 1 minuto.
\item \textbf{Notify email}: indica si los cambios de estado en este chequeo
  deberán ser notificados por correo electrónico.
\item \textbf{Notify Android}: indica si los cambios de estado en este chequeo
  deberán ser notificados mediante la aplicación Android.
\item \textbf{Should check timeout}: indica si hay que verificar que el ping
  enviado vuelve dentro de un tiempo máximo determinado.
\item \textbf{Timeout value}: si la anterior opción está activa, este campo
  define el número de milisegundos que puede tardar en responder el servidor al
  ping, como máximo.
\end{itemize}

\subsubsection{Creación de un chequeo de puerto}

\begin{figure}[hbtp]
  \centering
  \includegraphics[width=0.8\textwidth]{apendice_manual_usuario/pantalla_crear_port.png}
  \caption{Formulario de creación de chequeo de puerto}
  \label{fig:crear-port}
\end{figure}

El formulario para crear un chequeo de puertos aparece en la
figura~\ref{fig:crear-port}. Los campos que deberá rellenar son los siguientes:

\begin{itemize}
\item \textbf{Título} del chequeo.
\item \textbf{Descripción} opcional del chequeo.
\item \textbf{Target}: es el objetivo a chequear. Debe ser un \textit{hostname} o una IP.
\item \textbf{Target port}: es el puerto a chequear. Debe ser un entero.
\item \textbf{Check for string}: opcionalmente, se puede indicar una cadena que
  debe estar presente en la respuesta al conectar con el puerto.
\item \textbf{Check interval}: intervalo entre chequeos, definido en minutos. El
  valor mínimo es de 1 minuto.
\item \textbf{Notify email}: indica si los cambios de estado en este chequeo
  deberán ser notificados por correo electrónico.
\item \textbf{Notify Android}: indica si los cambios de estado en este chequeo
  deberán ser notificados mediante la aplicación Android.
\end{itemize}

\subsubsection{Creación de un chequeo de DNS}

\begin{figure}[hbtp]
  \centering
  \includegraphics[width=0.8\textwidth]{apendice_manual_usuario/pantalla_crear_dns.png}
  \caption{Formulario de creación de chequeo de registro DNS}
  \label{fig:crear-dns}
\end{figure}

El formulario para crear un chequeo de registro DNS aparece en la figura
\ref{fig:crear-dns}. Los campos que deberá rellenar son los siguientes:

\begin{itemize}
\item \textbf{Título} del chequeo.
\item \textbf{Descripción} opcional del chequeo.
\item \textbf{Target}: es el objetivo a chequear. Debe ser un \textit{hostname}.
\item \textbf{Record type}: indica el tipo de registro a chequear. Puede ser A, AAAA, CNAME, TXT y MX.
\item \textbf{Resolved address}: indica el contenido que debe tener el registro.
\item \textbf{Check interval}: intervalo entre chequeos, definido en minutos. El
  valor mínimo es de 1 minuto.
\item \textbf{Notify email}: indica si los cambios de estado en este chequeo
  deberán ser notificados por correo electrónico.
\item \textbf{Notify Android}: indica si los cambios de estado en este chequeo
  deberán ser notificados mediante la aplicación Android.
\end{itemize}

\subsubsection{Creación de un chequeo por HTTP}

\begin{figure}[hbtp]
  \centering
  \includegraphics[width=0.8\textwidth]{apendice_manual_usuario/pantalla_crear_http.png}
  \caption{Formulario de creación de chequeo por HTTP}
  \label{fig:crear-http}
\end{figure}

El formulario para crear un chequeo por HTTP aparece en la figura
\ref{fig:crear-http}. Los campos que deberá rellenar son los siguientes:

\begin{itemize}
\item \textbf{Título} del chequeo.
\item \textbf{Descripción} opcional del chequeo.
\item \textbf{Target}: es el objetivo a chequear. Debe ser una URL.
\item \textbf{Status code}: es el código de estado que deberá devolver la petición..
\item \textbf{Check for string}: opcionalmente, se puede indicar una cadena que
  debe estar presente en la respuesta a la petición.
\item \textbf{Check interval}: intervalo entre chequeos, definido en minutos. El
  valor mínimo es de 1 minuto.
\item \textbf{Notify email}: indica si los cambios de estado en este chequeo
  deberán ser notificados por correo electrónico.
\item \textbf{Notify Android}: indica si los cambios de estado en este chequeo
  deberán ser notificados mediante la aplicación Android.
\end{itemize}


\subsection{Operaciones con chequeos}

\begin{figure}[hbtp]
  \centering
  \includegraphics[width=\textwidth]{apendice_manual_usuario/chequeo_valido.png}
  \caption{}
  \label{fig:chequeo}
\end{figure}

Una vez creado un chequeo, aparecerá en el panel de control con el aspecto de la
figura~\ref{fig:chequeo}. Como se aprecia en la figura, se muestra mediante un
icono el estado del chequeo, los datos (título, tipo, descripción) y una gráfica
que ilustra el estado del chequeo en las últimas 24 horas.

\begin{figure}[hbtp]
  \centering
  \includegraphics[width=\textwidth]{apendice_manual_usuario/chequeo_invalido.png}
  \caption{}
  \label{fig:chequeo-invalido}
\end{figure}

Si un chequeo detecta un error, cambia de apariencia, pasando a tener el aspecto
que ofrece la figura \ref{fig:chequeo-invalido}. Además, según se haya
configurado en las opciones del chequeo, se envía una notificación al usuario
informando del cambio de estado del chequeo.

Cuando se ve en la lista general de chequeos, cada chequeo tiene una serie de
botones que permiten llevar a cabo operaciones diversas, que se detallan a
continuación.

\subsubsection{Activación y desactivación de un chequeo}

Para desactivar un chequeo, tiene que pulsar el botón \textbf{Disable} de la
botonera, que solo aparecerá si el chequeo está activado. De igual modo, si el
chequeo está desactivado, el botón que aparecerá tendrá la etiqueta
\textbf{Enable}, que podrá pulsar para activarlo.

\subsubsection{Borrado de un chequeo}

En la botonera se sitúa el botón \textbf{Delete}, que podrá pulsar para borrar
el chequeo. Aparecerá un formulario pidiéndole confirmación. Si confirma la
acción, se procederá al borrado del chequeo y todos los datos dependientes de
éste.

\subsubsection{Edición de un chequeo}

Si desea editar un chequeo debe pulsar el botón \textbf{Edit}, situado en la
botonera del chequeo. Aparecerá un formulario similar al de creación de chequeo,
pero con los datos actuales ya introducidos. Puede modificar los datos que desee
y confirmar los cambios pulsando el botón \textbf{Update check} del formulario.

\subsubsection{Vista de detalle del chequeo}

\begin{figure}[hbtp]
  \centering
  \includegraphics[width=\textwidth]{apendice_manual_usuario/pantalla_detalle_chequeo.png}
  \caption{}
  \label{fig:detalle-chequeo}
\end{figure}

En la botonera del chequeo, en el panel de control, verá un botón
\textbf{Details}. Si lo pulsa será llevado a la pantalla de vista de detalle del
chequeo, que se ilustra en la figura \ref{fig:detalle-chequeo}. Esta pantalla
ofrece mucha información sobre el chequeo. En particular:

\begin{itemize}
\item En la parte superior se sitúa la información general:
  \begin{itemize}
  \item Icono con el estado del chequeo.
  \item Título.
  \item Tipo de chequeo.
  \item Descripción.
  \item Botonera de acciones.
  \end{itemize}

\item En la parte media de la pantalla se sitúa un bloque con pestañas. Mediante
  estas pestañas es posible indicar sobre qué periodo de tiempo queremos
  visualizar los datos: últimas 24 horas o última semana. La información
  disponible es la siguiente:

  \begin{itemize}
  \item Porcentaje de \textit{uptime} del chequeo.
  \item \textit{(Solo chequeos tipo ping)} Tiempo de respuesta mínimo, máximo y promedio.
  \item Gráfica con el estado del chequeo.
  \end{itemize}

\item En la parte inferior de la pantalla se sitúa el panel de eventos del
  chequeo. Éste permite ver los últimos cambios de estado del chequeo:

  \begin{itemize}
  \item Tipo de cambio de estado (\textit{Up}, \textit{Down}).
  \item Duración del estado.
  \item Hora de comienzo del estado.
  \item Hora de fin del estado. Si es el estado en el que el chequeo se
    encuentra actualmente, no aparecerá ninguna fecha.
  \item Extra: muestra información adicional para estados de error.
  \end{itemize}
\end{itemize}

\begin{figure}[hbtp]
  \centering
  \includegraphics[width=\textwidth]{apendice_manual_usuario/pantalla_grafica_semanal.png}
  \caption{Detalle de la gráfica para un intervalo de tiempo de una semana}
  \label{fig:grafica-semanal}
\end{figure} 

\subsection{Recepción de notificaciones por correo electrónico}

\subsubsection{Notificaciones de cambio de estado}

Cuando el sistema detecte un cambio de estado en un chequeo, y esté activada la
notificación por correo electrónico, se enviará un email como el que aparece en
la figura \ref{fig:email-estado}.

\begin{figure}[hbtp]
  \centering
  \includegraphics[width=0.7\textwidth]{apendice_manual_usuario/email_estado.png}
  \caption{Email de notificación de cambio de estado}
  \label{fig:email-estado}
\end{figure}

En ese email se le indicará el título del chequeo, el nuevo estado, la fecha y
hora en la que ha pasado al nuevo estado, y un enlace directo al panel de
detalle del chequeo en SiteUp, de forma que pueda obtener más información.

\subsubsection{Notificaciones de resumen diarias}

Si ha activado el envío de resúmenes diarios sobre el estado de sus chequeos, el
sistema enviará todos los días un correo electrónico con información. El aspecto
del correo que recibirá se puede ver en la figura \ref{fig:email-resumen}

\begin{figure}[hbtp]
  \centering
  \includegraphics[width=0.7\textwidth]{apendice_manual_usuario/email_resumen.png}
  \caption{Email diario de resumen}
  \label{fig:email-resumen}
\end{figure}


\section{Uso de la aplicación móvil}

La aplicación móvil \textbf{SiteUp Client} es una aplicación para para el
sistema operativo Android que permite a sus usuarios revisar los chequeos que
tiene dados de alta en la plataforma web SiteUp, así como recibir de manera
instantánea notificaciones \textit{push} sobre el estado de sus chequeos.

\subsection{Instalación}

La aplicación móvil se instala copiando y ejecutando el fichero APK en el
dispositivo. En un futuro se podrá encontrar en la \textit{Play Store} de
Google.

\subsection{Lanzamiento}

En la figura~\ref{fig:manual-android-launcher} se puede observar el aspecto de
la aplicación \textbf{SiteUp Client} en el lanzador de Android. Para ejecutarla,
simplemente pulse sobre el icono de la aplicación.

\begin{figure}[htbp]
  \centering
  \includegraphics[width=0.4\textwidth]{5_diseno/android-1}
  \caption{Icono de SiteUp en el lanzador}
  \label{fig:manual-android-launcher}
\end{figure}

\subsection{Carga de la aplicación}

Tras lanzar la aplicación, el sistema cargará los datos necesarios y se
comunicará con los servidores de Google para obtener los datos necesarios. La
apariencia de la aplicación durante este proceso se puede ver en la figura
\ref{fig:manual-android-login}.

\begin{figure}[htbp]
  \centering
  \includegraphics[width=0.4\textwidth]{5_diseno/android-2}
  \caption{Pantalla de carga de la aplicación móvil}
  \label{fig:manual-android-loading}
\end{figure}

\subsection{Inicio de sesión}

Tras la carga, la primera vez que lance la aplicación será necesario que inicie
sesión con las credenciales que utilizó en la plataforma web. Para ello, tendrá
que introducir su nombre de usuario y contraseña en el formulario que aparece en
la figura~\ref{fig:manual-android-loading}, y pulsar el botón \textbf{Sign in}.

\begin{figure}[htbp]
  \centering
  \includegraphics[width=0.4\textwidth]{5_diseno/android-3}
  \caption{Pantalla de inicio de sesión de la aplicación móvil}
  \label{fig:manual-android-login}
\end{figure}

La aplicación móvil se comunicará con los servidores centrales de SiteUp,
validará sus credenciales de usuario y, si son correctas, le dará acceso al
sistema. Si las credenciales de usuario no son correctas, el formulario mostrará
un error.

\subsection{Revisión de chequeos}

Si ha iniciado sesión correctamente, tendrá la opción de revisar la lista de
chequeos que tiene dados de alta en la plataforma web \textbf{SiteUp}. La
apariencia de la lista de chequeos es igual a la de la
figura~\ref{fig:manual-android-dashboard}.

\begin{figure}[hbtp]
  \centering
  \includegraphics[width=0.4\textwidth]{5_diseno/android-4}
  \caption{Pantalla de chequeos de la aplicación móvil}
  \label{fig:manual-android-dashboard}
\end{figure}

Si pulsa en cualquiera de los chequeos, será redirigido a la plataforma web, que
podrá navegar sin problema gracias a que se encuentra preparada para su uso en
dispositivos móviles.

\subsection{Recepción de notificaciones}

La utilidad principal de la aplicación es la posibilidad de recibir
notificaciones en cualquier momento provenientes del servidor -- conocidas como
notificaciones \textit{push}. 

El procedimiento es transparente al usuario. Simplemente lance la aplicación e
inice sesión. Tras hacer esto, la aplicación estará preparada para recibir
notificaciones desde el servidor.

Cuando se reciba una notificación, la aplicación mostará un mensaje en el área
de notificaciones del móvil, con información resumida. Puede ver el aspecto de
la notificación en la figura \ref{fig:manual-android-notification}. Si pulsa en
la notificación será llevado a los detalles del chequeo para obtener más
detalles sobre el origen del aviso.

\begin{figure}[hbtp]
  \centering
  \includegraphics[width=0.4\textwidth]{5_diseno/android-5}
  \caption{Aspecto de las notificaciones}
  \label{fig:manual-android-notification}
\end{figure}



%%% Local Variables:
%%% mode: latex
%%% TeX-master: "../memoria"
%%% End:


\chapter{Pruebas}
\label{chap:pruebas}
En este capítulo se explicarán en detalle las instrucciones de uso de la
plataforma web.

\section{Introducción a SiteUp}

\textbf{SiteUp} es un servicio web creado para la monitorización de servicios de
Internet, especialmente pensado para dueños de sitios web y servidores que
deseen estar informados en caso de cualquier eventualidad.

El uso de internet en todos los ámbitos de la vida, desde las interacciones
sociales hasta la búsqueda de empleo, han dado lugar a la creación de servicios
accesibles a través de Internet y a nuevos modelos de negocio basados en
aquellos. Resulta por tanto imprescindible contar con sistemas para monitorizar
la disponibilidad de estos servicios y, en caso necesario, actuar de manera que
puedan subsanarse las causas de los problemas detectados.

SiteUp cubre precisamente estas necesidades, permitiendo a los usuarios:

\begin{itemize}
\item Crear chequeos mediante envío de paquetes ICMP (Ping), para revisar que un
  servidor está activo y conectado a la red.
\item Crear chequeos de registros DNS, de forma que el contenido de estos
  registros sea el correcto en todo momento.
\item Crear chequeos de puertos remotos, para vigilar que los servicios
  concretos dentro de un servidor son accesibles.
\item Crear chequeos mediante peticiones HTTP, para controlar que un servicio
  web responde de forma correcta y con el contenido adecuado.
\item Establecer las opciones de periodicidad de estos chequeos, de acuerdo a
  las necesidades particulares de cada caso.
\item Recibir notificaciones tanto a través de correo electrónico como mediante
  una aplicación móvil para el sistema operativo Android.
\end{itemize}

\section{Requisitos previos}

Para el uso correcto de la plataforma web SiteUp, el principal requisito es
contar con un dispositivo con acceso a internet que disponga de un navegador
reciente, como por ejemplo Google Chrome. Cualquier tipo de tablet, teléfono
móvil u ordenador personal podrá acceder a la plataforma web SiteUp, pero para
aprovechar mejor el potencial del sistema se recomienda usar una computadora
personal con las siguientes características:

\begin{itemize}
\item Navegador reciente. Se recomienda el uso de Google Chrome.
\item Pantalla con resolución mínima de 1280 píxeles en su lado longitudinal.
\item Conexión a internet con un caudal de descarga de al menos 512kb/s.
\end{itemize}

Si desea hacer uso de la aplicación móvil, será necesario que el usuario cuente
con un dispositivo con las siguientes características mínimas:

\begin{itemize}
\item Sistema operativo Android, versión mínima 2.3, versión recomendada 4.3.
\item Pantalla de al menos 3 pulgadas de diagonal y resolución mínima de 320x480
  píxeles.
\end{itemize}

\section{Uso de la plataforma web}

\subsection{Acceso}

A fecha de abril de 2014, la dirección de acceso a la plataforma web de SiteUp
es \url{http://siteup.josetomastocino.com}. Simplemente introduzca esa dirección
en su navegador y automáticamente será redirigido a la plataforma web, donde
verá la pantalla de bienvenida.

\subsection{Registro de usuario}

\begin{figure}[hbtp]
  \centering
  \includegraphics[width=0.3\textwidth]{apendice_manual_usuario/general_botones.png}
  \caption{Botones principales para el usuario anónimo}
  \label{fig:botones-1}
\end{figure}

Si es la primera vez que accede al sitio web deberá registrarse. Para ello, debe
pulsar en el botón \textbf{Sign up}, que se ve en la figura \ref{fig:botones-1}.

Tras ello, aparecerá el formulario de registro de usuario, que se ve en la
figura \ref{fig:registro}

\begin{figure}[hbtp]
  \centering
  \includegraphics[width=0.5\textwidth]{apendice_manual_usuario/formulario_registro.png}
  \caption{Formulario de registro de usuario}
  \label{fig:registro}
\end{figure}

En el formulario deberá introducir los siguientes datos:

\begin{itemize}
\item \textbf{Nombre de usuario}. Debe tener 30 caracteres como máximo, y solo
  pueden usarse letras, números y guiones.
\item \textbf{Dirección de correo electrónico}. Asegúrese de que es una
  dirección de correo correcta.
\item \textbf{Contraseña}.
\end{itemize}

Tras introducir los datos, deberá pulsar el botón \textbf{Sign up} del
formulario. Si el registro es correcto, el sistema dará de alta su usuario y
mostrará la pantalla \ref{fig:registro-correcto}. Si no, volverá a aparecer el
formulario con información sobre los errores en los datos.

\begin{figure}[hbtp]
  \centering
  \includegraphics[width=\textwidth]{apendice_manual_usuario/pantalla_register_completo.png}
  \caption{Mensaje indicando el registro correcto}
  \label{fig:registro-correcto}
\end{figure}

\subsection{Inicio de sesión}

Una vez que haya registrado su usuario, deberá pulsar en el botón \textbf{Log
  in}, como aparece en la imagen \ref{fig:botones-1}. Esto le llevará al
formulario de inicio de sesión que aparece en la figura \ref{fig:login}

\begin{figure}[hbtp]
  \centering
  \includegraphics[width=0.5\textwidth]{apendice_manual_usuario/general_formulario_login.png}
  \caption{Formulario de inicio de sesión}
  \label{fig:login}
\end{figure}

Deberá introducir su nombre de usuario y la contraseña que escribió al
registrarse, y posteriormente pulsar en el botón \textbf{Log in} del
formulario. Si los datos son correctos, será llevado al panel principal de la
web, donde podrá gestionar sus chequeos. Si no son correctos, volverá a aparecer
el formulario, con información sobre los errores encontrados.

\subsection{Recuperación de contraseña}

Si se registró como usuario pero no recuerda su contraseña es posible
recuperarla.

Para ello, vaya al formulario de inicio de sesión como ser ha explicado en el
paso anterior. En ese formulario, en la zona inferior, verá un enlace titulado
\textbf{Forgot password?} que deberá pulsar. Después de hacerlo se mostrará el
formulario de recuperación de contraseña, ilustrado en la figura
\ref{fig:formulario-reset}.

\begin{figure}[hbtp]
  \centering
  \includegraphics[width=0.5\textwidth]{apendice_manual_usuario/general_formulario_reset.png}
  \caption{Formulario de recuperación de contraseña}
  \label{fig:formulario-reset}
\end{figure}

En este formulario deberá introducir la dirección de correo electrónico que
utilizó al registrarse, y pulsar el botón \textbf{Send reset code}. Si lo hace
correctamente, aparecerá la pantalla de confirmación que se ve en la figura
\ref{fig:pantalla-reset-enviado}.

\begin{figure}[hbtp]
  \centering
  \includegraphics[width=0.5\textwidth]{apendice_manual_usuario/pantalla_reset_enviado.png}
  \caption{Pantalla de confirmación de envío de código de reinicio}
  \label{fig:pantalla-reset-enviado}
\end{figure}

El sistema habrá enviado un enlace a su dirección de correo electrónico. Deberá
pulsar el enlace que aparece en ese correo. Asegúrese de que recibe
correctamente el correo. Si no lo recibe, revise la carpeta de correo no
deseado. El enlace del correo le llevará a un formulario de reseteo de
contraseña como el que aparece en la figura \ref{fig:formulario-reset-2}.

\begin{figure}[hbtp]
  \centering
  \includegraphics[width=0.5\textwidth]{apendice_manual_usuario/general_formulario_reset_2.png}
  \caption{}
  \label{fig:formulario-reset-2}
\end{figure}

En este formulario deberá escribir la nueva contraseña de su cuenta dos veces,
para confirmar que la escribe correctamente. Tras ello, deberá pulsar el botón
\textbf{Set new password}, lo cual confirmará los cambios. 

Tras este proceso, podrá iniciar sesión como se ha descrito previamente, esta
vez utilizando la nueva contraseña.


\subsection{Pantalla de listado de chequeos}

\begin{figure}[hbtp]
  \centering
  \includegraphics[width=0.8\textwidth]{apendice_manual_usuario/pantalla-dashboard.png}
  \caption{Panel de control}
  \label{fig:dashboard}
\end{figure}

\begin{figure}[hbtp]
  \centering
  \includegraphics[width=0.4\textwidth]{apendice_manual_usuario/general_botones_2.png}
  \caption{Botones principales para el usuario logueado}
  \label{fig:botones-2}
\end{figure}

En la figura \ref{fig:dashboard} se puede ver el aspecto del panel de control de
la plataforma web SiteUp. Desde este panel podrá hacer una gran cantidad de
acciones, desde crear grupos de chequeos hasta editar sus detalles de usuario.

Es posible acceder a esta pantalla desde cualquier sitio de la web si está
logueado, pulsando en el botón \textbf{Dashboard} de la cabecera de la web, tal
y como se ve en la figura \ref{fig:botones-2}.

\subsection{Cerrar sesión}

Si ha iniciado sesión y desea cerrarla, debe dirigirse a los botones
principales, situados en el lado derecho de la cabecera de la web, tal y como
aparecen en la figura \ref{fig:botones-2}. Deberá pulsar el botón \textbf{Log
  out}, tras lo cual se cerrara su sesión y será redirigido a la pantalla
inicial de SiteUp.

\subsection{Edición de datos de usuario}

Si desea editar alguno de los datos de usuario que ingresó cuando hizo el
registro, deberá primero pulsar el botón \textbf{Your profile}, situado en la
cabecera de la web. Es el segundo de los botones que aparecen, tal y como se ve
en la figura \ref{fig:botones-2}.

Una vez hecho esto, aparecerá el formulario de edición de usuario, ilustrado en
la figura \ref{fig:formulario-usuario}, en el que verá sus actuales datos de
usuario. 

\begin{figure}[hbtp]
  \centering
  \includegraphics[width=0.8\textwidth]{apendice_manual_usuario/pantalla_perfil_usuario.png}
  \caption{}
  \label{fig:formulario-usuario}
\end{figure}

En este formulario podrá editar:

\begin{itemize}
\item Su \textbf{nombre} de usuario.
\item Su \textbf{dirección} de correo electrónico.
\item Si desea o no recibir \textbf{resúmenes diarios} sobre el estado de sus chequeos.
\end{itemize}

Una vez editados los datos que desee deberá pulsar el botón \textbf{Update
  details}. También podrá editar su contraseña pulsando el botón \textbf{Change
  password}, situado al pie del formulario a la izquierda.


\subsection{Creación de grupo de chequeos}

En SiteUp, los chequeos se organizan en \textbf{grupos de chequeos}, lo que
permite trabajar con ellos más fácilmente a la hora de activarlos, desactivarlos
o borrarlos. 

Para crear un grupo de chequeos deberá pulsar el botón \textbf{Add group},
situado en la zona superior derecha del panel principal. El botón puede verse en
la figura \ref{fig:crear-grupo}.

\begin{figure}[hbtp]
  \centering
  \includegraphics[width=0.3\textwidth]{apendice_manual_usuario/grupo_boton_nuevo.png}
  \caption{Botón para crear nuevo grupo de chequeos}
  \label{fig:crear-grupo}
\end{figure}

Tras pulsar el botón aparecerá el formulario de creación de nuevo grupo. En él
deberá introducir los datos del nuevo grupo. En particular:

\begin{itemize}
\item El \textbf{título} que identifica al grupo.
\end{itemize}

Tras ello, deberá pulsar el botón \textbf{Create group}, y el grupo será creado

\begin{figure}[H]
  \centering
  \includegraphics[width=0.8\textwidth]{apendice_manual_usuario/grupo_formulario_crear.png}
  \caption{Formulario de creación de grupo de chequeos}
  \label{fig:}
\end{figure}


\subsection{Operaciones con grupos de chequeos}

En la figura~\ref{fig:operaciones-grupo} se puede ver el menú de operaciones
disponibles para un grupo de chequeos. Este menú aparece en el panel principal,
al lado derecho junto al título de cada grupo de chequeos.

\begin{figure}[hbtp]
  \centering
  \includegraphics[width=0.5\textwidth]{apendice_manual_usuario/grupo_submenu_acciones.png}
  \caption{}
  \label{fig:operaciones-grupo}
\end{figure}

\subsubsection{Edición de grupo de chequeos}

En el menú de operaciones del grupo de chequeos, pulse el botón
\textbf{Edit}. Aparecerá un formulario similar al de creación de grupo, en el
que podrá editar el título del grupo de chequeos.

\subsubsection{Borrado de grupo de chequeos}

Para borrar un grupo de chequeos, pulse el botón \textbf{Delete} en el menú de
operaciones del grupo. Aparecerá un formulario como el de la figura
\ref{fig:formulario-borrado-grupo}, pidiéndole confirmación del borrado. 

\begin{figure}[hbtp]
  \centering
  \includegraphics[width=0.5\textwidth]{apendice_manual_usuario/formulario_borrado_grupo}
  \caption{Formulario de confirmación de borrado}
  \label{fig:formulario-borrado-grupo}
\end{figure}

Si pulsa en el botón \textbf{Go back}, volverá al panel principal, cancelando el
borrado. Si pulsa el botón \textbf{Delete}, confirmará la operación y el sistema
borrará el grupo de chequeos, así como todos los chequeos que contenga.

\subsubsection{Activación y desactivación de grupo de chequeos}

Es posible activar y desactivar de forma grupal todos los chequeos de un
grupo. Para ello, deberá dirigirse al menú de operaciones del grupo y pulsar en
\textbf{Enable all checks} o en \textbf{Disable all checks}, según sea su decisión.

\subsection{Creación de chequeo}

\begin{figure}[hbtp]
  \centering
  \includegraphics[width=0.8\textwidth]{apendice_manual_usuario/pantalla_seleccion_tipo.png}
  \caption{Pantalla de selección del tipo de chequeo a crear}
  \label{fig:seleccionar-tipo-check}
\end{figure}

Como se ha comentado previamente, el elemento principal de SiteUp es el
\textbf{chequeo}. El usuario da de alta chequeos dentro de un grupo, y estos
chequeos se lanzan periódicamente, generando datos.

Para crear un chequeo dentro de un grupo deberá ir al menú de operaciones del
grupo, que aparece en la figura \ref{fig:operaciones-grupo} y se sitúa a la
derecha del título de un grupo de chequeos. Dentro del menú, deberá pulsar la
opción \textbf{Add check}.



Tras esto aparecerá una pantalla en la que deberá seleccionar el tipo de chequeo
que desea crear. La pantalla es similar a la de la figura
\ref{fig:seleccionar-tipo-check}. Como se explicó en la introducción de este
manual, es posible crear cuatro tipo de chequeos: chequeos de ping, chequeos de
puerto, chequeos de registros DNS y chequeos por HTTP.

Si pulsa en el botón del tipo que desea crear, aparecerá un formulario con
campos para rellenar.

\subsubsection{Creación de un chequeo por ping}

\begin{figure}[hbtp]
  \centering
  \includegraphics[width=0.8\textwidth]{apendice_manual_usuario/pantalla_crear_ping.png}
  \caption{Formulario de creación de chequeo por ping}
  \label{fig:crear-ping}
\end{figure}

El formulario para crear un chequeo de tipo ping aparece en la
figura~\ref{fig:crear-ping}. Los campos que deberá rellenar son los siguientes:

\begin{itemize}
\item \textbf{Target}: es el objetivo a chequear. Debe ser un \textit{hostname} o una IP.
\item \textbf{Título} del chequeo.
\item \textbf{Descripción} opcional del chequeo.
\item \textbf{Check interval}: intervalo entre chequeos, definido en minutos. El
  valor mínimo es de 1 minuto.
\item \textbf{Notify email}: indica si los cambios de estado en este chequeo
  deberán ser notificados por correo electrónico.
\item \textbf{Notify Android}: indica si los cambios de estado en este chequeo
  deberán ser notificados mediante la aplicación Android.
\item \textbf{Should check timeout}: indica si hay que verificar que el ping
  enviado vuelve dentro de un tiempo máximo determinado.
\item \textbf{Timeout value}: si la anterior opción está activa, este campo
  define el número de milisegundos que puede tardar en responder el servidor al
  ping, como máximo.
\end{itemize}

\subsubsection{Creación de un chequeo de puerto}

\begin{figure}[hbtp]
  \centering
  \includegraphics[width=0.8\textwidth]{apendice_manual_usuario/pantalla_crear_port.png}
  \caption{Formulario de creación de chequeo de puerto}
  \label{fig:crear-port}
\end{figure}

El formulario para crear un chequeo de puertos aparece en la
figura~\ref{fig:crear-port}. Los campos que deberá rellenar son los siguientes:

\begin{itemize}
\item \textbf{Título} del chequeo.
\item \textbf{Descripción} opcional del chequeo.
\item \textbf{Target}: es el objetivo a chequear. Debe ser un \textit{hostname} o una IP.
\item \textbf{Target port}: es el puerto a chequear. Debe ser un entero.
\item \textbf{Check for string}: opcionalmente, se puede indicar una cadena que
  debe estar presente en la respuesta al conectar con el puerto.
\item \textbf{Check interval}: intervalo entre chequeos, definido en minutos. El
  valor mínimo es de 1 minuto.
\item \textbf{Notify email}: indica si los cambios de estado en este chequeo
  deberán ser notificados por correo electrónico.
\item \textbf{Notify Android}: indica si los cambios de estado en este chequeo
  deberán ser notificados mediante la aplicación Android.
\end{itemize}

\subsubsection{Creación de un chequeo de DNS}

\begin{figure}[hbtp]
  \centering
  \includegraphics[width=0.8\textwidth]{apendice_manual_usuario/pantalla_crear_dns.png}
  \caption{Formulario de creación de chequeo de registro DNS}
  \label{fig:crear-dns}
\end{figure}

El formulario para crear un chequeo de registro DNS aparece en la figura
\ref{fig:crear-dns}. Los campos que deberá rellenar son los siguientes:

\begin{itemize}
\item \textbf{Título} del chequeo.
\item \textbf{Descripción} opcional del chequeo.
\item \textbf{Target}: es el objetivo a chequear. Debe ser un \textit{hostname}.
\item \textbf{Record type}: indica el tipo de registro a chequear. Puede ser A, AAAA, CNAME, TXT y MX.
\item \textbf{Resolved address}: indica el contenido que debe tener el registro.
\item \textbf{Check interval}: intervalo entre chequeos, definido en minutos. El
  valor mínimo es de 1 minuto.
\item \textbf{Notify email}: indica si los cambios de estado en este chequeo
  deberán ser notificados por correo electrónico.
\item \textbf{Notify Android}: indica si los cambios de estado en este chequeo
  deberán ser notificados mediante la aplicación Android.
\end{itemize}

\subsubsection{Creación de un chequeo por HTTP}

\begin{figure}[hbtp]
  \centering
  \includegraphics[width=0.8\textwidth]{apendice_manual_usuario/pantalla_crear_http.png}
  \caption{Formulario de creación de chequeo por HTTP}
  \label{fig:crear-http}
\end{figure}

El formulario para crear un chequeo por HTTP aparece en la figura
\ref{fig:crear-http}. Los campos que deberá rellenar son los siguientes:

\begin{itemize}
\item \textbf{Título} del chequeo.
\item \textbf{Descripción} opcional del chequeo.
\item \textbf{Target}: es el objetivo a chequear. Debe ser una URL.
\item \textbf{Status code}: es el código de estado que deberá devolver la petición..
\item \textbf{Check for string}: opcionalmente, se puede indicar una cadena que
  debe estar presente en la respuesta a la petición.
\item \textbf{Check interval}: intervalo entre chequeos, definido en minutos. El
  valor mínimo es de 1 minuto.
\item \textbf{Notify email}: indica si los cambios de estado en este chequeo
  deberán ser notificados por correo electrónico.
\item \textbf{Notify Android}: indica si los cambios de estado en este chequeo
  deberán ser notificados mediante la aplicación Android.
\end{itemize}


\subsection{Operaciones con chequeos}

\begin{figure}[hbtp]
  \centering
  \includegraphics[width=\textwidth]{apendice_manual_usuario/chequeo_valido.png}
  \caption{}
  \label{fig:chequeo}
\end{figure}

Una vez creado un chequeo, aparecerá en el panel de control con el aspecto de la
figura~\ref{fig:chequeo}. Como se aprecia en la figura, se muestra mediante un
icono el estado del chequeo, los datos (título, tipo, descripción) y una gráfica
que ilustra el estado del chequeo en las últimas 24 horas.

\begin{figure}[hbtp]
  \centering
  \includegraphics[width=\textwidth]{apendice_manual_usuario/chequeo_invalido.png}
  \caption{}
  \label{fig:chequeo-invalido}
\end{figure}

Si un chequeo detecta un error, cambia de apariencia, pasando a tener el aspecto
que ofrece la figura \ref{fig:chequeo-invalido}. Además, según se haya
configurado en las opciones del chequeo, se envía una notificación al usuario
informando del cambio de estado del chequeo.

Cuando se ve en la lista general de chequeos, cada chequeo tiene una serie de
botones que permiten llevar a cabo operaciones diversas, que se detallan a
continuación.

\subsubsection{Activación y desactivación de un chequeo}

Para desactivar un chequeo, tiene que pulsar el botón \textbf{Disable} de la
botonera, que solo aparecerá si el chequeo está activado. De igual modo, si el
chequeo está desactivado, el botón que aparecerá tendrá la etiqueta
\textbf{Enable}, que podrá pulsar para activarlo.

\subsubsection{Borrado de un chequeo}

En la botonera se sitúa el botón \textbf{Delete}, que podrá pulsar para borrar
el chequeo. Aparecerá un formulario pidiéndole confirmación. Si confirma la
acción, se procederá al borrado del chequeo y todos los datos dependientes de
éste.

\subsubsection{Edición de un chequeo}

Si desea editar un chequeo debe pulsar el botón \textbf{Edit}, situado en la
botonera del chequeo. Aparecerá un formulario similar al de creación de chequeo,
pero con los datos actuales ya introducidos. Puede modificar los datos que desee
y confirmar los cambios pulsando el botón \textbf{Update check} del formulario.

\subsubsection{Vista de detalle del chequeo}

\begin{figure}[hbtp]
  \centering
  \includegraphics[width=\textwidth]{apendice_manual_usuario/pantalla_detalle_chequeo.png}
  \caption{}
  \label{fig:detalle-chequeo}
\end{figure}

En la botonera del chequeo, en el panel de control, verá un botón
\textbf{Details}. Si lo pulsa será llevado a la pantalla de vista de detalle del
chequeo, que se ilustra en la figura \ref{fig:detalle-chequeo}. Esta pantalla
ofrece mucha información sobre el chequeo. En particular:

\begin{itemize}
\item En la parte superior se sitúa la información general:
  \begin{itemize}
  \item Icono con el estado del chequeo.
  \item Título.
  \item Tipo de chequeo.
  \item Descripción.
  \item Botonera de acciones.
  \end{itemize}

\item En la parte media de la pantalla se sitúa un bloque con pestañas. Mediante
  estas pestañas es posible indicar sobre qué periodo de tiempo queremos
  visualizar los datos: últimas 24 horas o última semana. La información
  disponible es la siguiente:

  \begin{itemize}
  \item Porcentaje de \textit{uptime} del chequeo.
  \item \textit{(Solo chequeos tipo ping)} Tiempo de respuesta mínimo, máximo y promedio.
  \item Gráfica con el estado del chequeo.
  \end{itemize}

\item En la parte inferior de la pantalla se sitúa el panel de eventos del
  chequeo. Éste permite ver los últimos cambios de estado del chequeo:

  \begin{itemize}
  \item Tipo de cambio de estado (\textit{Up}, \textit{Down}).
  \item Duración del estado.
  \item Hora de comienzo del estado.
  \item Hora de fin del estado. Si es el estado en el que el chequeo se
    encuentra actualmente, no aparecerá ninguna fecha.
  \item Extra: muestra información adicional para estados de error.
  \end{itemize}
\end{itemize}

\begin{figure}[hbtp]
  \centering
  \includegraphics[width=\textwidth]{apendice_manual_usuario/pantalla_grafica_semanal.png}
  \caption{Detalle de la gráfica para un intervalo de tiempo de una semana}
  \label{fig:grafica-semanal}
\end{figure} 

\subsection{Recepción de notificaciones por correo electrónico}

\subsubsection{Notificaciones de cambio de estado}

Cuando el sistema detecte un cambio de estado en un chequeo, y esté activada la
notificación por correo electrónico, se enviará un email como el que aparece en
la figura \ref{fig:email-estado}.

\begin{figure}[hbtp]
  \centering
  \includegraphics[width=0.7\textwidth]{apendice_manual_usuario/email_estado.png}
  \caption{Email de notificación de cambio de estado}
  \label{fig:email-estado}
\end{figure}

En ese email se le indicará el título del chequeo, el nuevo estado, la fecha y
hora en la que ha pasado al nuevo estado, y un enlace directo al panel de
detalle del chequeo en SiteUp, de forma que pueda obtener más información.

\subsubsection{Notificaciones de resumen diarias}

Si ha activado el envío de resúmenes diarios sobre el estado de sus chequeos, el
sistema enviará todos los días un correo electrónico con información. El aspecto
del correo que recibirá se puede ver en la figura \ref{fig:email-resumen}

\begin{figure}[hbtp]
  \centering
  \includegraphics[width=0.7\textwidth]{apendice_manual_usuario/email_resumen.png}
  \caption{Email diario de resumen}
  \label{fig:email-resumen}
\end{figure}


\section{Uso de la aplicación móvil}

La aplicación móvil \textbf{SiteUp Client} es una aplicación para para el
sistema operativo Android que permite a sus usuarios revisar los chequeos que
tiene dados de alta en la plataforma web SiteUp, así como recibir de manera
instantánea notificaciones \textit{push} sobre el estado de sus chequeos.

\subsection{Instalación}

La aplicación móvil se instala copiando y ejecutando el fichero APK en el
dispositivo. En un futuro se podrá encontrar en la \textit{Play Store} de
Google.

\subsection{Lanzamiento}

En la figura~\ref{fig:manual-android-launcher} se puede observar el aspecto de
la aplicación \textbf{SiteUp Client} en el lanzador de Android. Para ejecutarla,
simplemente pulse sobre el icono de la aplicación.

\begin{figure}[htbp]
  \centering
  \includegraphics[width=0.4\textwidth]{5_diseno/android-1}
  \caption{Icono de SiteUp en el lanzador}
  \label{fig:manual-android-launcher}
\end{figure}

\subsection{Carga de la aplicación}

Tras lanzar la aplicación, el sistema cargará los datos necesarios y se
comunicará con los servidores de Google para obtener los datos necesarios. La
apariencia de la aplicación durante este proceso se puede ver en la figura
\ref{fig:manual-android-login}.

\begin{figure}[htbp]
  \centering
  \includegraphics[width=0.4\textwidth]{5_diseno/android-2}
  \caption{Pantalla de carga de la aplicación móvil}
  \label{fig:manual-android-loading}
\end{figure}

\subsection{Inicio de sesión}

Tras la carga, la primera vez que lance la aplicación será necesario que inicie
sesión con las credenciales que utilizó en la plataforma web. Para ello, tendrá
que introducir su nombre de usuario y contraseña en el formulario que aparece en
la figura~\ref{fig:manual-android-loading}, y pulsar el botón \textbf{Sign in}.

\begin{figure}[htbp]
  \centering
  \includegraphics[width=0.4\textwidth]{5_diseno/android-3}
  \caption{Pantalla de inicio de sesión de la aplicación móvil}
  \label{fig:manual-android-login}
\end{figure}

La aplicación móvil se comunicará con los servidores centrales de SiteUp,
validará sus credenciales de usuario y, si son correctas, le dará acceso al
sistema. Si las credenciales de usuario no son correctas, el formulario mostrará
un error.

\subsection{Revisión de chequeos}

Si ha iniciado sesión correctamente, tendrá la opción de revisar la lista de
chequeos que tiene dados de alta en la plataforma web \textbf{SiteUp}. La
apariencia de la lista de chequeos es igual a la de la
figura~\ref{fig:manual-android-dashboard}.

\begin{figure}[hbtp]
  \centering
  \includegraphics[width=0.4\textwidth]{5_diseno/android-4}
  \caption{Pantalla de chequeos de la aplicación móvil}
  \label{fig:manual-android-dashboard}
\end{figure}

Si pulsa en cualquiera de los chequeos, será redirigido a la plataforma web, que
podrá navegar sin problema gracias a que se encuentra preparada para su uso en
dispositivos móviles.

\subsection{Recepción de notificaciones}

La utilidad principal de la aplicación es la posibilidad de recibir
notificaciones en cualquier momento provenientes del servidor -- conocidas como
notificaciones \textit{push}. 

El procedimiento es transparente al usuario. Simplemente lance la aplicación e
inice sesión. Tras hacer esto, la aplicación estará preparada para recibir
notificaciones desde el servidor.

Cuando se reciba una notificación, la aplicación mostará un mensaje en el área
de notificaciones del móvil, con información resumida. Puede ver el aspecto de
la notificación en la figura \ref{fig:manual-android-notification}. Si pulsa en
la notificación será llevado a los detalles del chequeo para obtener más
detalles sobre el origen del aviso.

\begin{figure}[hbtp]
  \centering
  \includegraphics[width=0.4\textwidth]{5_diseno/android-5}
  \caption{Aspecto de las notificaciones}
  \label{fig:manual-android-notification}
\end{figure}



%%% Local Variables:
%%% mode: latex
%%% TeX-master: "../memoria"
%%% End:


\chapter{Conclusiones}
\label{chap:conclusiones}
En este capítulo se explicarán en detalle las instrucciones de uso de la
plataforma web.

\section{Introducción a SiteUp}

\textbf{SiteUp} es un servicio web creado para la monitorización de servicios de
Internet, especialmente pensado para dueños de sitios web y servidores que
deseen estar informados en caso de cualquier eventualidad.

El uso de internet en todos los ámbitos de la vida, desde las interacciones
sociales hasta la búsqueda de empleo, han dado lugar a la creación de servicios
accesibles a través de Internet y a nuevos modelos de negocio basados en
aquellos. Resulta por tanto imprescindible contar con sistemas para monitorizar
la disponibilidad de estos servicios y, en caso necesario, actuar de manera que
puedan subsanarse las causas de los problemas detectados.

SiteUp cubre precisamente estas necesidades, permitiendo a los usuarios:

\begin{itemize}
\item Crear chequeos mediante envío de paquetes ICMP (Ping), para revisar que un
  servidor está activo y conectado a la red.
\item Crear chequeos de registros DNS, de forma que el contenido de estos
  registros sea el correcto en todo momento.
\item Crear chequeos de puertos remotos, para vigilar que los servicios
  concretos dentro de un servidor son accesibles.
\item Crear chequeos mediante peticiones HTTP, para controlar que un servicio
  web responde de forma correcta y con el contenido adecuado.
\item Establecer las opciones de periodicidad de estos chequeos, de acuerdo a
  las necesidades particulares de cada caso.
\item Recibir notificaciones tanto a través de correo electrónico como mediante
  una aplicación móvil para el sistema operativo Android.
\end{itemize}

\section{Requisitos previos}

Para el uso correcto de la plataforma web SiteUp, el principal requisito es
contar con un dispositivo con acceso a internet que disponga de un navegador
reciente, como por ejemplo Google Chrome. Cualquier tipo de tablet, teléfono
móvil u ordenador personal podrá acceder a la plataforma web SiteUp, pero para
aprovechar mejor el potencial del sistema se recomienda usar una computadora
personal con las siguientes características:

\begin{itemize}
\item Navegador reciente. Se recomienda el uso de Google Chrome.
\item Pantalla con resolución mínima de 1280 píxeles en su lado longitudinal.
\item Conexión a internet con un caudal de descarga de al menos 512kb/s.
\end{itemize}

Si desea hacer uso de la aplicación móvil, será necesario que el usuario cuente
con un dispositivo con las siguientes características mínimas:

\begin{itemize}
\item Sistema operativo Android, versión mínima 2.3, versión recomendada 4.3.
\item Pantalla de al menos 3 pulgadas de diagonal y resolución mínima de 320x480
  píxeles.
\end{itemize}

\section{Uso de la plataforma web}

\subsection{Acceso}

A fecha de abril de 2014, la dirección de acceso a la plataforma web de SiteUp
es \url{http://siteup.josetomastocino.com}. Simplemente introduzca esa dirección
en su navegador y automáticamente será redirigido a la plataforma web, donde
verá la pantalla de bienvenida.

\subsection{Registro de usuario}

\begin{figure}[hbtp]
  \centering
  \includegraphics[width=0.3\textwidth]{apendice_manual_usuario/general_botones.png}
  \caption{Botones principales para el usuario anónimo}
  \label{fig:botones-1}
\end{figure}

Si es la primera vez que accede al sitio web deberá registrarse. Para ello, debe
pulsar en el botón \textbf{Sign up}, que se ve en la figura \ref{fig:botones-1}.

Tras ello, aparecerá el formulario de registro de usuario, que se ve en la
figura \ref{fig:registro}

\begin{figure}[hbtp]
  \centering
  \includegraphics[width=0.5\textwidth]{apendice_manual_usuario/formulario_registro.png}
  \caption{Formulario de registro de usuario}
  \label{fig:registro}
\end{figure}

En el formulario deberá introducir los siguientes datos:

\begin{itemize}
\item \textbf{Nombre de usuario}. Debe tener 30 caracteres como máximo, y solo
  pueden usarse letras, números y guiones.
\item \textbf{Dirección de correo electrónico}. Asegúrese de que es una
  dirección de correo correcta.
\item \textbf{Contraseña}.
\end{itemize}

Tras introducir los datos, deberá pulsar el botón \textbf{Sign up} del
formulario. Si el registro es correcto, el sistema dará de alta su usuario y
mostrará la pantalla \ref{fig:registro-correcto}. Si no, volverá a aparecer el
formulario con información sobre los errores en los datos.

\begin{figure}[hbtp]
  \centering
  \includegraphics[width=\textwidth]{apendice_manual_usuario/pantalla_register_completo.png}
  \caption{Mensaje indicando el registro correcto}
  \label{fig:registro-correcto}
\end{figure}

\subsection{Inicio de sesión}

Una vez que haya registrado su usuario, deberá pulsar en el botón \textbf{Log
  in}, como aparece en la imagen \ref{fig:botones-1}. Esto le llevará al
formulario de inicio de sesión que aparece en la figura \ref{fig:login}

\begin{figure}[hbtp]
  \centering
  \includegraphics[width=0.5\textwidth]{apendice_manual_usuario/general_formulario_login.png}
  \caption{Formulario de inicio de sesión}
  \label{fig:login}
\end{figure}

Deberá introducir su nombre de usuario y la contraseña que escribió al
registrarse, y posteriormente pulsar en el botón \textbf{Log in} del
formulario. Si los datos son correctos, será llevado al panel principal de la
web, donde podrá gestionar sus chequeos. Si no son correctos, volverá a aparecer
el formulario, con información sobre los errores encontrados.

\subsection{Recuperación de contraseña}

Si se registró como usuario pero no recuerda su contraseña es posible
recuperarla.

Para ello, vaya al formulario de inicio de sesión como ser ha explicado en el
paso anterior. En ese formulario, en la zona inferior, verá un enlace titulado
\textbf{Forgot password?} que deberá pulsar. Después de hacerlo se mostrará el
formulario de recuperación de contraseña, ilustrado en la figura
\ref{fig:formulario-reset}.

\begin{figure}[hbtp]
  \centering
  \includegraphics[width=0.5\textwidth]{apendice_manual_usuario/general_formulario_reset.png}
  \caption{Formulario de recuperación de contraseña}
  \label{fig:formulario-reset}
\end{figure}

En este formulario deberá introducir la dirección de correo electrónico que
utilizó al registrarse, y pulsar el botón \textbf{Send reset code}. Si lo hace
correctamente, aparecerá la pantalla de confirmación que se ve en la figura
\ref{fig:pantalla-reset-enviado}.

\begin{figure}[hbtp]
  \centering
  \includegraphics[width=0.5\textwidth]{apendice_manual_usuario/pantalla_reset_enviado.png}
  \caption{Pantalla de confirmación de envío de código de reinicio}
  \label{fig:pantalla-reset-enviado}
\end{figure}

El sistema habrá enviado un enlace a su dirección de correo electrónico. Deberá
pulsar el enlace que aparece en ese correo. Asegúrese de que recibe
correctamente el correo. Si no lo recibe, revise la carpeta de correo no
deseado. El enlace del correo le llevará a un formulario de reseteo de
contraseña como el que aparece en la figura \ref{fig:formulario-reset-2}.

\begin{figure}[hbtp]
  \centering
  \includegraphics[width=0.5\textwidth]{apendice_manual_usuario/general_formulario_reset_2.png}
  \caption{}
  \label{fig:formulario-reset-2}
\end{figure}

En este formulario deberá escribir la nueva contraseña de su cuenta dos veces,
para confirmar que la escribe correctamente. Tras ello, deberá pulsar el botón
\textbf{Set new password}, lo cual confirmará los cambios. 

Tras este proceso, podrá iniciar sesión como se ha descrito previamente, esta
vez utilizando la nueva contraseña.


\subsection{Pantalla de listado de chequeos}

\begin{figure}[hbtp]
  \centering
  \includegraphics[width=0.8\textwidth]{apendice_manual_usuario/pantalla-dashboard.png}
  \caption{Panel de control}
  \label{fig:dashboard}
\end{figure}

\begin{figure}[hbtp]
  \centering
  \includegraphics[width=0.4\textwidth]{apendice_manual_usuario/general_botones_2.png}
  \caption{Botones principales para el usuario logueado}
  \label{fig:botones-2}
\end{figure}

En la figura \ref{fig:dashboard} se puede ver el aspecto del panel de control de
la plataforma web SiteUp. Desde este panel podrá hacer una gran cantidad de
acciones, desde crear grupos de chequeos hasta editar sus detalles de usuario.

Es posible acceder a esta pantalla desde cualquier sitio de la web si está
logueado, pulsando en el botón \textbf{Dashboard} de la cabecera de la web, tal
y como se ve en la figura \ref{fig:botones-2}.

\subsection{Cerrar sesión}

Si ha iniciado sesión y desea cerrarla, debe dirigirse a los botones
principales, situados en el lado derecho de la cabecera de la web, tal y como
aparecen en la figura \ref{fig:botones-2}. Deberá pulsar el botón \textbf{Log
  out}, tras lo cual se cerrara su sesión y será redirigido a la pantalla
inicial de SiteUp.

\subsection{Edición de datos de usuario}

Si desea editar alguno de los datos de usuario que ingresó cuando hizo el
registro, deberá primero pulsar el botón \textbf{Your profile}, situado en la
cabecera de la web. Es el segundo de los botones que aparecen, tal y como se ve
en la figura \ref{fig:botones-2}.

Una vez hecho esto, aparecerá el formulario de edición de usuario, ilustrado en
la figura \ref{fig:formulario-usuario}, en el que verá sus actuales datos de
usuario. 

\begin{figure}[hbtp]
  \centering
  \includegraphics[width=0.8\textwidth]{apendice_manual_usuario/pantalla_perfil_usuario.png}
  \caption{}
  \label{fig:formulario-usuario}
\end{figure}

En este formulario podrá editar:

\begin{itemize}
\item Su \textbf{nombre} de usuario.
\item Su \textbf{dirección} de correo electrónico.
\item Si desea o no recibir \textbf{resúmenes diarios} sobre el estado de sus chequeos.
\end{itemize}

Una vez editados los datos que desee deberá pulsar el botón \textbf{Update
  details}. También podrá editar su contraseña pulsando el botón \textbf{Change
  password}, situado al pie del formulario a la izquierda.


\subsection{Creación de grupo de chequeos}

En SiteUp, los chequeos se organizan en \textbf{grupos de chequeos}, lo que
permite trabajar con ellos más fácilmente a la hora de activarlos, desactivarlos
o borrarlos. 

Para crear un grupo de chequeos deberá pulsar el botón \textbf{Add group},
situado en la zona superior derecha del panel principal. El botón puede verse en
la figura \ref{fig:crear-grupo}.

\begin{figure}[hbtp]
  \centering
  \includegraphics[width=0.3\textwidth]{apendice_manual_usuario/grupo_boton_nuevo.png}
  \caption{Botón para crear nuevo grupo de chequeos}
  \label{fig:crear-grupo}
\end{figure}

Tras pulsar el botón aparecerá el formulario de creación de nuevo grupo. En él
deberá introducir los datos del nuevo grupo. En particular:

\begin{itemize}
\item El \textbf{título} que identifica al grupo.
\end{itemize}

Tras ello, deberá pulsar el botón \textbf{Create group}, y el grupo será creado

\begin{figure}[H]
  \centering
  \includegraphics[width=0.8\textwidth]{apendice_manual_usuario/grupo_formulario_crear.png}
  \caption{Formulario de creación de grupo de chequeos}
  \label{fig:}
\end{figure}


\subsection{Operaciones con grupos de chequeos}

En la figura~\ref{fig:operaciones-grupo} se puede ver el menú de operaciones
disponibles para un grupo de chequeos. Este menú aparece en el panel principal,
al lado derecho junto al título de cada grupo de chequeos.

\begin{figure}[hbtp]
  \centering
  \includegraphics[width=0.5\textwidth]{apendice_manual_usuario/grupo_submenu_acciones.png}
  \caption{}
  \label{fig:operaciones-grupo}
\end{figure}

\subsubsection{Edición de grupo de chequeos}

En el menú de operaciones del grupo de chequeos, pulse el botón
\textbf{Edit}. Aparecerá un formulario similar al de creación de grupo, en el
que podrá editar el título del grupo de chequeos.

\subsubsection{Borrado de grupo de chequeos}

Para borrar un grupo de chequeos, pulse el botón \textbf{Delete} en el menú de
operaciones del grupo. Aparecerá un formulario como el de la figura
\ref{fig:formulario-borrado-grupo}, pidiéndole confirmación del borrado. 

\begin{figure}[hbtp]
  \centering
  \includegraphics[width=0.5\textwidth]{apendice_manual_usuario/formulario_borrado_grupo}
  \caption{Formulario de confirmación de borrado}
  \label{fig:formulario-borrado-grupo}
\end{figure}

Si pulsa en el botón \textbf{Go back}, volverá al panel principal, cancelando el
borrado. Si pulsa el botón \textbf{Delete}, confirmará la operación y el sistema
borrará el grupo de chequeos, así como todos los chequeos que contenga.

\subsubsection{Activación y desactivación de grupo de chequeos}

Es posible activar y desactivar de forma grupal todos los chequeos de un
grupo. Para ello, deberá dirigirse al menú de operaciones del grupo y pulsar en
\textbf{Enable all checks} o en \textbf{Disable all checks}, según sea su decisión.

\subsection{Creación de chequeo}

\begin{figure}[hbtp]
  \centering
  \includegraphics[width=0.8\textwidth]{apendice_manual_usuario/pantalla_seleccion_tipo.png}
  \caption{Pantalla de selección del tipo de chequeo a crear}
  \label{fig:seleccionar-tipo-check}
\end{figure}

Como se ha comentado previamente, el elemento principal de SiteUp es el
\textbf{chequeo}. El usuario da de alta chequeos dentro de un grupo, y estos
chequeos se lanzan periódicamente, generando datos.

Para crear un chequeo dentro de un grupo deberá ir al menú de operaciones del
grupo, que aparece en la figura \ref{fig:operaciones-grupo} y se sitúa a la
derecha del título de un grupo de chequeos. Dentro del menú, deberá pulsar la
opción \textbf{Add check}.



Tras esto aparecerá una pantalla en la que deberá seleccionar el tipo de chequeo
que desea crear. La pantalla es similar a la de la figura
\ref{fig:seleccionar-tipo-check}. Como se explicó en la introducción de este
manual, es posible crear cuatro tipo de chequeos: chequeos de ping, chequeos de
puerto, chequeos de registros DNS y chequeos por HTTP.

Si pulsa en el botón del tipo que desea crear, aparecerá un formulario con
campos para rellenar.

\subsubsection{Creación de un chequeo por ping}

\begin{figure}[hbtp]
  \centering
  \includegraphics[width=0.8\textwidth]{apendice_manual_usuario/pantalla_crear_ping.png}
  \caption{Formulario de creación de chequeo por ping}
  \label{fig:crear-ping}
\end{figure}

El formulario para crear un chequeo de tipo ping aparece en la
figura~\ref{fig:crear-ping}. Los campos que deberá rellenar son los siguientes:

\begin{itemize}
\item \textbf{Target}: es el objetivo a chequear. Debe ser un \textit{hostname} o una IP.
\item \textbf{Título} del chequeo.
\item \textbf{Descripción} opcional del chequeo.
\item \textbf{Check interval}: intervalo entre chequeos, definido en minutos. El
  valor mínimo es de 1 minuto.
\item \textbf{Notify email}: indica si los cambios de estado en este chequeo
  deberán ser notificados por correo electrónico.
\item \textbf{Notify Android}: indica si los cambios de estado en este chequeo
  deberán ser notificados mediante la aplicación Android.
\item \textbf{Should check timeout}: indica si hay que verificar que el ping
  enviado vuelve dentro de un tiempo máximo determinado.
\item \textbf{Timeout value}: si la anterior opción está activa, este campo
  define el número de milisegundos que puede tardar en responder el servidor al
  ping, como máximo.
\end{itemize}

\subsubsection{Creación de un chequeo de puerto}

\begin{figure}[hbtp]
  \centering
  \includegraphics[width=0.8\textwidth]{apendice_manual_usuario/pantalla_crear_port.png}
  \caption{Formulario de creación de chequeo de puerto}
  \label{fig:crear-port}
\end{figure}

El formulario para crear un chequeo de puertos aparece en la
figura~\ref{fig:crear-port}. Los campos que deberá rellenar son los siguientes:

\begin{itemize}
\item \textbf{Título} del chequeo.
\item \textbf{Descripción} opcional del chequeo.
\item \textbf{Target}: es el objetivo a chequear. Debe ser un \textit{hostname} o una IP.
\item \textbf{Target port}: es el puerto a chequear. Debe ser un entero.
\item \textbf{Check for string}: opcionalmente, se puede indicar una cadena que
  debe estar presente en la respuesta al conectar con el puerto.
\item \textbf{Check interval}: intervalo entre chequeos, definido en minutos. El
  valor mínimo es de 1 minuto.
\item \textbf{Notify email}: indica si los cambios de estado en este chequeo
  deberán ser notificados por correo electrónico.
\item \textbf{Notify Android}: indica si los cambios de estado en este chequeo
  deberán ser notificados mediante la aplicación Android.
\end{itemize}

\subsubsection{Creación de un chequeo de DNS}

\begin{figure}[hbtp]
  \centering
  \includegraphics[width=0.8\textwidth]{apendice_manual_usuario/pantalla_crear_dns.png}
  \caption{Formulario de creación de chequeo de registro DNS}
  \label{fig:crear-dns}
\end{figure}

El formulario para crear un chequeo de registro DNS aparece en la figura
\ref{fig:crear-dns}. Los campos que deberá rellenar son los siguientes:

\begin{itemize}
\item \textbf{Título} del chequeo.
\item \textbf{Descripción} opcional del chequeo.
\item \textbf{Target}: es el objetivo a chequear. Debe ser un \textit{hostname}.
\item \textbf{Record type}: indica el tipo de registro a chequear. Puede ser A, AAAA, CNAME, TXT y MX.
\item \textbf{Resolved address}: indica el contenido que debe tener el registro.
\item \textbf{Check interval}: intervalo entre chequeos, definido en minutos. El
  valor mínimo es de 1 minuto.
\item \textbf{Notify email}: indica si los cambios de estado en este chequeo
  deberán ser notificados por correo electrónico.
\item \textbf{Notify Android}: indica si los cambios de estado en este chequeo
  deberán ser notificados mediante la aplicación Android.
\end{itemize}

\subsubsection{Creación de un chequeo por HTTP}

\begin{figure}[hbtp]
  \centering
  \includegraphics[width=0.8\textwidth]{apendice_manual_usuario/pantalla_crear_http.png}
  \caption{Formulario de creación de chequeo por HTTP}
  \label{fig:crear-http}
\end{figure}

El formulario para crear un chequeo por HTTP aparece en la figura
\ref{fig:crear-http}. Los campos que deberá rellenar son los siguientes:

\begin{itemize}
\item \textbf{Título} del chequeo.
\item \textbf{Descripción} opcional del chequeo.
\item \textbf{Target}: es el objetivo a chequear. Debe ser una URL.
\item \textbf{Status code}: es el código de estado que deberá devolver la petición..
\item \textbf{Check for string}: opcionalmente, se puede indicar una cadena que
  debe estar presente en la respuesta a la petición.
\item \textbf{Check interval}: intervalo entre chequeos, definido en minutos. El
  valor mínimo es de 1 minuto.
\item \textbf{Notify email}: indica si los cambios de estado en este chequeo
  deberán ser notificados por correo electrónico.
\item \textbf{Notify Android}: indica si los cambios de estado en este chequeo
  deberán ser notificados mediante la aplicación Android.
\end{itemize}


\subsection{Operaciones con chequeos}

\begin{figure}[hbtp]
  \centering
  \includegraphics[width=\textwidth]{apendice_manual_usuario/chequeo_valido.png}
  \caption{}
  \label{fig:chequeo}
\end{figure}

Una vez creado un chequeo, aparecerá en el panel de control con el aspecto de la
figura~\ref{fig:chequeo}. Como se aprecia en la figura, se muestra mediante un
icono el estado del chequeo, los datos (título, tipo, descripción) y una gráfica
que ilustra el estado del chequeo en las últimas 24 horas.

\begin{figure}[hbtp]
  \centering
  \includegraphics[width=\textwidth]{apendice_manual_usuario/chequeo_invalido.png}
  \caption{}
  \label{fig:chequeo-invalido}
\end{figure}

Si un chequeo detecta un error, cambia de apariencia, pasando a tener el aspecto
que ofrece la figura \ref{fig:chequeo-invalido}. Además, según se haya
configurado en las opciones del chequeo, se envía una notificación al usuario
informando del cambio de estado del chequeo.

Cuando se ve en la lista general de chequeos, cada chequeo tiene una serie de
botones que permiten llevar a cabo operaciones diversas, que se detallan a
continuación.

\subsubsection{Activación y desactivación de un chequeo}

Para desactivar un chequeo, tiene que pulsar el botón \textbf{Disable} de la
botonera, que solo aparecerá si el chequeo está activado. De igual modo, si el
chequeo está desactivado, el botón que aparecerá tendrá la etiqueta
\textbf{Enable}, que podrá pulsar para activarlo.

\subsubsection{Borrado de un chequeo}

En la botonera se sitúa el botón \textbf{Delete}, que podrá pulsar para borrar
el chequeo. Aparecerá un formulario pidiéndole confirmación. Si confirma la
acción, se procederá al borrado del chequeo y todos los datos dependientes de
éste.

\subsubsection{Edición de un chequeo}

Si desea editar un chequeo debe pulsar el botón \textbf{Edit}, situado en la
botonera del chequeo. Aparecerá un formulario similar al de creación de chequeo,
pero con los datos actuales ya introducidos. Puede modificar los datos que desee
y confirmar los cambios pulsando el botón \textbf{Update check} del formulario.

\subsubsection{Vista de detalle del chequeo}

\begin{figure}[hbtp]
  \centering
  \includegraphics[width=\textwidth]{apendice_manual_usuario/pantalla_detalle_chequeo.png}
  \caption{}
  \label{fig:detalle-chequeo}
\end{figure}

En la botonera del chequeo, en el panel de control, verá un botón
\textbf{Details}. Si lo pulsa será llevado a la pantalla de vista de detalle del
chequeo, que se ilustra en la figura \ref{fig:detalle-chequeo}. Esta pantalla
ofrece mucha información sobre el chequeo. En particular:

\begin{itemize}
\item En la parte superior se sitúa la información general:
  \begin{itemize}
  \item Icono con el estado del chequeo.
  \item Título.
  \item Tipo de chequeo.
  \item Descripción.
  \item Botonera de acciones.
  \end{itemize}

\item En la parte media de la pantalla se sitúa un bloque con pestañas. Mediante
  estas pestañas es posible indicar sobre qué periodo de tiempo queremos
  visualizar los datos: últimas 24 horas o última semana. La información
  disponible es la siguiente:

  \begin{itemize}
  \item Porcentaje de \textit{uptime} del chequeo.
  \item \textit{(Solo chequeos tipo ping)} Tiempo de respuesta mínimo, máximo y promedio.
  \item Gráfica con el estado del chequeo.
  \end{itemize}

\item En la parte inferior de la pantalla se sitúa el panel de eventos del
  chequeo. Éste permite ver los últimos cambios de estado del chequeo:

  \begin{itemize}
  \item Tipo de cambio de estado (\textit{Up}, \textit{Down}).
  \item Duración del estado.
  \item Hora de comienzo del estado.
  \item Hora de fin del estado. Si es el estado en el que el chequeo se
    encuentra actualmente, no aparecerá ninguna fecha.
  \item Extra: muestra información adicional para estados de error.
  \end{itemize}
\end{itemize}

\begin{figure}[hbtp]
  \centering
  \includegraphics[width=\textwidth]{apendice_manual_usuario/pantalla_grafica_semanal.png}
  \caption{Detalle de la gráfica para un intervalo de tiempo de una semana}
  \label{fig:grafica-semanal}
\end{figure} 

\subsection{Recepción de notificaciones por correo electrónico}

\subsubsection{Notificaciones de cambio de estado}

Cuando el sistema detecte un cambio de estado en un chequeo, y esté activada la
notificación por correo electrónico, se enviará un email como el que aparece en
la figura \ref{fig:email-estado}.

\begin{figure}[hbtp]
  \centering
  \includegraphics[width=0.7\textwidth]{apendice_manual_usuario/email_estado.png}
  \caption{Email de notificación de cambio de estado}
  \label{fig:email-estado}
\end{figure}

En ese email se le indicará el título del chequeo, el nuevo estado, la fecha y
hora en la que ha pasado al nuevo estado, y un enlace directo al panel de
detalle del chequeo en SiteUp, de forma que pueda obtener más información.

\subsubsection{Notificaciones de resumen diarias}

Si ha activado el envío de resúmenes diarios sobre el estado de sus chequeos, el
sistema enviará todos los días un correo electrónico con información. El aspecto
del correo que recibirá se puede ver en la figura \ref{fig:email-resumen}

\begin{figure}[hbtp]
  \centering
  \includegraphics[width=0.7\textwidth]{apendice_manual_usuario/email_resumen.png}
  \caption{Email diario de resumen}
  \label{fig:email-resumen}
\end{figure}


\section{Uso de la aplicación móvil}

La aplicación móvil \textbf{SiteUp Client} es una aplicación para para el
sistema operativo Android que permite a sus usuarios revisar los chequeos que
tiene dados de alta en la plataforma web SiteUp, así como recibir de manera
instantánea notificaciones \textit{push} sobre el estado de sus chequeos.

\subsection{Instalación}

La aplicación móvil se instala copiando y ejecutando el fichero APK en el
dispositivo. En un futuro se podrá encontrar en la \textit{Play Store} de
Google.

\subsection{Lanzamiento}

En la figura~\ref{fig:manual-android-launcher} se puede observar el aspecto de
la aplicación \textbf{SiteUp Client} en el lanzador de Android. Para ejecutarla,
simplemente pulse sobre el icono de la aplicación.

\begin{figure}[htbp]
  \centering
  \includegraphics[width=0.4\textwidth]{5_diseno/android-1}
  \caption{Icono de SiteUp en el lanzador}
  \label{fig:manual-android-launcher}
\end{figure}

\subsection{Carga de la aplicación}

Tras lanzar la aplicación, el sistema cargará los datos necesarios y se
comunicará con los servidores de Google para obtener los datos necesarios. La
apariencia de la aplicación durante este proceso se puede ver en la figura
\ref{fig:manual-android-login}.

\begin{figure}[htbp]
  \centering
  \includegraphics[width=0.4\textwidth]{5_diseno/android-2}
  \caption{Pantalla de carga de la aplicación móvil}
  \label{fig:manual-android-loading}
\end{figure}

\subsection{Inicio de sesión}

Tras la carga, la primera vez que lance la aplicación será necesario que inicie
sesión con las credenciales que utilizó en la plataforma web. Para ello, tendrá
que introducir su nombre de usuario y contraseña en el formulario que aparece en
la figura~\ref{fig:manual-android-loading}, y pulsar el botón \textbf{Sign in}.

\begin{figure}[htbp]
  \centering
  \includegraphics[width=0.4\textwidth]{5_diseno/android-3}
  \caption{Pantalla de inicio de sesión de la aplicación móvil}
  \label{fig:manual-android-login}
\end{figure}

La aplicación móvil se comunicará con los servidores centrales de SiteUp,
validará sus credenciales de usuario y, si son correctas, le dará acceso al
sistema. Si las credenciales de usuario no son correctas, el formulario mostrará
un error.

\subsection{Revisión de chequeos}

Si ha iniciado sesión correctamente, tendrá la opción de revisar la lista de
chequeos que tiene dados de alta en la plataforma web \textbf{SiteUp}. La
apariencia de la lista de chequeos es igual a la de la
figura~\ref{fig:manual-android-dashboard}.

\begin{figure}[hbtp]
  \centering
  \includegraphics[width=0.4\textwidth]{5_diseno/android-4}
  \caption{Pantalla de chequeos de la aplicación móvil}
  \label{fig:manual-android-dashboard}
\end{figure}

Si pulsa en cualquiera de los chequeos, será redirigido a la plataforma web, que
podrá navegar sin problema gracias a que se encuentra preparada para su uso en
dispositivos móviles.

\subsection{Recepción de notificaciones}

La utilidad principal de la aplicación es la posibilidad de recibir
notificaciones en cualquier momento provenientes del servidor -- conocidas como
notificaciones \textit{push}. 

El procedimiento es transparente al usuario. Simplemente lance la aplicación e
inice sesión. Tras hacer esto, la aplicación estará preparada para recibir
notificaciones desde el servidor.

Cuando se reciba una notificación, la aplicación mostará un mensaje en el área
de notificaciones del móvil, con información resumida. Puede ver el aspecto de
la notificación en la figura \ref{fig:manual-android-notification}. Si pulsa en
la notificación será llevado a los detalles del chequeo para obtener más
detalles sobre el origen del aviso.

\begin{figure}[hbtp]
  \centering
  \includegraphics[width=0.4\textwidth]{5_diseno/android-5}
  \caption{Aspecto de las notificaciones}
  \label{fig:manual-android-notification}
\end{figure}



%%% Local Variables:
%%% mode: latex
%%% TeX-master: "../memoria"
%%% End:


\appendix

\chapter{Manual de instalación de la plataforma web}
\label{chap:manual_instalacion}
En este capítulo se explicarán en detalle las instrucciones de uso de la
plataforma web.

\section{Introducción a SiteUp}

\textbf{SiteUp} es un servicio web creado para la monitorización de servicios de
Internet, especialmente pensado para dueños de sitios web y servidores que
deseen estar informados en caso de cualquier eventualidad.

El uso de internet en todos los ámbitos de la vida, desde las interacciones
sociales hasta la búsqueda de empleo, han dado lugar a la creación de servicios
accesibles a través de Internet y a nuevos modelos de negocio basados en
aquellos. Resulta por tanto imprescindible contar con sistemas para monitorizar
la disponibilidad de estos servicios y, en caso necesario, actuar de manera que
puedan subsanarse las causas de los problemas detectados.

SiteUp cubre precisamente estas necesidades, permitiendo a los usuarios:

\begin{itemize}
\item Crear chequeos mediante envío de paquetes ICMP (Ping), para revisar que un
  servidor está activo y conectado a la red.
\item Crear chequeos de registros DNS, de forma que el contenido de estos
  registros sea el correcto en todo momento.
\item Crear chequeos de puertos remotos, para vigilar que los servicios
  concretos dentro de un servidor son accesibles.
\item Crear chequeos mediante peticiones HTTP, para controlar que un servicio
  web responde de forma correcta y con el contenido adecuado.
\item Establecer las opciones de periodicidad de estos chequeos, de acuerdo a
  las necesidades particulares de cada caso.
\item Recibir notificaciones tanto a través de correo electrónico como mediante
  una aplicación móvil para el sistema operativo Android.
\end{itemize}

\section{Requisitos previos}

Para el uso correcto de la plataforma web SiteUp, el principal requisito es
contar con un dispositivo con acceso a internet que disponga de un navegador
reciente, como por ejemplo Google Chrome. Cualquier tipo de tablet, teléfono
móvil u ordenador personal podrá acceder a la plataforma web SiteUp, pero para
aprovechar mejor el potencial del sistema se recomienda usar una computadora
personal con las siguientes características:

\begin{itemize}
\item Navegador reciente. Se recomienda el uso de Google Chrome.
\item Pantalla con resolución mínima de 1280 píxeles en su lado longitudinal.
\item Conexión a internet con un caudal de descarga de al menos 512kb/s.
\end{itemize}

Si desea hacer uso de la aplicación móvil, será necesario que el usuario cuente
con un dispositivo con las siguientes características mínimas:

\begin{itemize}
\item Sistema operativo Android, versión mínima 2.3, versión recomendada 4.3.
\item Pantalla de al menos 3 pulgadas de diagonal y resolución mínima de 320x480
  píxeles.
\end{itemize}

\section{Uso de la plataforma web}

\subsection{Acceso}

A fecha de abril de 2014, la dirección de acceso a la plataforma web de SiteUp
es \url{http://siteup.josetomastocino.com}. Simplemente introduzca esa dirección
en su navegador y automáticamente será redirigido a la plataforma web, donde
verá la pantalla de bienvenida.

\subsection{Registro de usuario}

\begin{figure}[hbtp]
  \centering
  \includegraphics[width=0.3\textwidth]{apendice_manual_usuario/general_botones.png}
  \caption{Botones principales para el usuario anónimo}
  \label{fig:botones-1}
\end{figure}

Si es la primera vez que accede al sitio web deberá registrarse. Para ello, debe
pulsar en el botón \textbf{Sign up}, que se ve en la figura \ref{fig:botones-1}.

Tras ello, aparecerá el formulario de registro de usuario, que se ve en la
figura \ref{fig:registro}

\begin{figure}[hbtp]
  \centering
  \includegraphics[width=0.5\textwidth]{apendice_manual_usuario/formulario_registro.png}
  \caption{Formulario de registro de usuario}
  \label{fig:registro}
\end{figure}

En el formulario deberá introducir los siguientes datos:

\begin{itemize}
\item \textbf{Nombre de usuario}. Debe tener 30 caracteres como máximo, y solo
  pueden usarse letras, números y guiones.
\item \textbf{Dirección de correo electrónico}. Asegúrese de que es una
  dirección de correo correcta.
\item \textbf{Contraseña}.
\end{itemize}

Tras introducir los datos, deberá pulsar el botón \textbf{Sign up} del
formulario. Si el registro es correcto, el sistema dará de alta su usuario y
mostrará la pantalla \ref{fig:registro-correcto}. Si no, volverá a aparecer el
formulario con información sobre los errores en los datos.

\begin{figure}[hbtp]
  \centering
  \includegraphics[width=\textwidth]{apendice_manual_usuario/pantalla_register_completo.png}
  \caption{Mensaje indicando el registro correcto}
  \label{fig:registro-correcto}
\end{figure}

\subsection{Inicio de sesión}

Una vez que haya registrado su usuario, deberá pulsar en el botón \textbf{Log
  in}, como aparece en la imagen \ref{fig:botones-1}. Esto le llevará al
formulario de inicio de sesión que aparece en la figura \ref{fig:login}

\begin{figure}[hbtp]
  \centering
  \includegraphics[width=0.5\textwidth]{apendice_manual_usuario/general_formulario_login.png}
  \caption{Formulario de inicio de sesión}
  \label{fig:login}
\end{figure}

Deberá introducir su nombre de usuario y la contraseña que escribió al
registrarse, y posteriormente pulsar en el botón \textbf{Log in} del
formulario. Si los datos son correctos, será llevado al panel principal de la
web, donde podrá gestionar sus chequeos. Si no son correctos, volverá a aparecer
el formulario, con información sobre los errores encontrados.

\subsection{Recuperación de contraseña}

Si se registró como usuario pero no recuerda su contraseña es posible
recuperarla.

Para ello, vaya al formulario de inicio de sesión como ser ha explicado en el
paso anterior. En ese formulario, en la zona inferior, verá un enlace titulado
\textbf{Forgot password?} que deberá pulsar. Después de hacerlo se mostrará el
formulario de recuperación de contraseña, ilustrado en la figura
\ref{fig:formulario-reset}.

\begin{figure}[hbtp]
  \centering
  \includegraphics[width=0.5\textwidth]{apendice_manual_usuario/general_formulario_reset.png}
  \caption{Formulario de recuperación de contraseña}
  \label{fig:formulario-reset}
\end{figure}

En este formulario deberá introducir la dirección de correo electrónico que
utilizó al registrarse, y pulsar el botón \textbf{Send reset code}. Si lo hace
correctamente, aparecerá la pantalla de confirmación que se ve en la figura
\ref{fig:pantalla-reset-enviado}.

\begin{figure}[hbtp]
  \centering
  \includegraphics[width=0.5\textwidth]{apendice_manual_usuario/pantalla_reset_enviado.png}
  \caption{Pantalla de confirmación de envío de código de reinicio}
  \label{fig:pantalla-reset-enviado}
\end{figure}

El sistema habrá enviado un enlace a su dirección de correo electrónico. Deberá
pulsar el enlace que aparece en ese correo. Asegúrese de que recibe
correctamente el correo. Si no lo recibe, revise la carpeta de correo no
deseado. El enlace del correo le llevará a un formulario de reseteo de
contraseña como el que aparece en la figura \ref{fig:formulario-reset-2}.

\begin{figure}[hbtp]
  \centering
  \includegraphics[width=0.5\textwidth]{apendice_manual_usuario/general_formulario_reset_2.png}
  \caption{}
  \label{fig:formulario-reset-2}
\end{figure}

En este formulario deberá escribir la nueva contraseña de su cuenta dos veces,
para confirmar que la escribe correctamente. Tras ello, deberá pulsar el botón
\textbf{Set new password}, lo cual confirmará los cambios. 

Tras este proceso, podrá iniciar sesión como se ha descrito previamente, esta
vez utilizando la nueva contraseña.


\subsection{Pantalla de listado de chequeos}

\begin{figure}[hbtp]
  \centering
  \includegraphics[width=0.8\textwidth]{apendice_manual_usuario/pantalla-dashboard.png}
  \caption{Panel de control}
  \label{fig:dashboard}
\end{figure}

\begin{figure}[hbtp]
  \centering
  \includegraphics[width=0.4\textwidth]{apendice_manual_usuario/general_botones_2.png}
  \caption{Botones principales para el usuario logueado}
  \label{fig:botones-2}
\end{figure}

En la figura \ref{fig:dashboard} se puede ver el aspecto del panel de control de
la plataforma web SiteUp. Desde este panel podrá hacer una gran cantidad de
acciones, desde crear grupos de chequeos hasta editar sus detalles de usuario.

Es posible acceder a esta pantalla desde cualquier sitio de la web si está
logueado, pulsando en el botón \textbf{Dashboard} de la cabecera de la web, tal
y como se ve en la figura \ref{fig:botones-2}.

\subsection{Cerrar sesión}

Si ha iniciado sesión y desea cerrarla, debe dirigirse a los botones
principales, situados en el lado derecho de la cabecera de la web, tal y como
aparecen en la figura \ref{fig:botones-2}. Deberá pulsar el botón \textbf{Log
  out}, tras lo cual se cerrara su sesión y será redirigido a la pantalla
inicial de SiteUp.

\subsection{Edición de datos de usuario}

Si desea editar alguno de los datos de usuario que ingresó cuando hizo el
registro, deberá primero pulsar el botón \textbf{Your profile}, situado en la
cabecera de la web. Es el segundo de los botones que aparecen, tal y como se ve
en la figura \ref{fig:botones-2}.

Una vez hecho esto, aparecerá el formulario de edición de usuario, ilustrado en
la figura \ref{fig:formulario-usuario}, en el que verá sus actuales datos de
usuario. 

\begin{figure}[hbtp]
  \centering
  \includegraphics[width=0.8\textwidth]{apendice_manual_usuario/pantalla_perfil_usuario.png}
  \caption{}
  \label{fig:formulario-usuario}
\end{figure}

En este formulario podrá editar:

\begin{itemize}
\item Su \textbf{nombre} de usuario.
\item Su \textbf{dirección} de correo electrónico.
\item Si desea o no recibir \textbf{resúmenes diarios} sobre el estado de sus chequeos.
\end{itemize}

Una vez editados los datos que desee deberá pulsar el botón \textbf{Update
  details}. También podrá editar su contraseña pulsando el botón \textbf{Change
  password}, situado al pie del formulario a la izquierda.


\subsection{Creación de grupo de chequeos}

En SiteUp, los chequeos se organizan en \textbf{grupos de chequeos}, lo que
permite trabajar con ellos más fácilmente a la hora de activarlos, desactivarlos
o borrarlos. 

Para crear un grupo de chequeos deberá pulsar el botón \textbf{Add group},
situado en la zona superior derecha del panel principal. El botón puede verse en
la figura \ref{fig:crear-grupo}.

\begin{figure}[hbtp]
  \centering
  \includegraphics[width=0.3\textwidth]{apendice_manual_usuario/grupo_boton_nuevo.png}
  \caption{Botón para crear nuevo grupo de chequeos}
  \label{fig:crear-grupo}
\end{figure}

Tras pulsar el botón aparecerá el formulario de creación de nuevo grupo. En él
deberá introducir los datos del nuevo grupo. En particular:

\begin{itemize}
\item El \textbf{título} que identifica al grupo.
\end{itemize}

Tras ello, deberá pulsar el botón \textbf{Create group}, y el grupo será creado

\begin{figure}[H]
  \centering
  \includegraphics[width=0.8\textwidth]{apendice_manual_usuario/grupo_formulario_crear.png}
  \caption{Formulario de creación de grupo de chequeos}
  \label{fig:}
\end{figure}


\subsection{Operaciones con grupos de chequeos}

En la figura~\ref{fig:operaciones-grupo} se puede ver el menú de operaciones
disponibles para un grupo de chequeos. Este menú aparece en el panel principal,
al lado derecho junto al título de cada grupo de chequeos.

\begin{figure}[hbtp]
  \centering
  \includegraphics[width=0.5\textwidth]{apendice_manual_usuario/grupo_submenu_acciones.png}
  \caption{}
  \label{fig:operaciones-grupo}
\end{figure}

\subsubsection{Edición de grupo de chequeos}

En el menú de operaciones del grupo de chequeos, pulse el botón
\textbf{Edit}. Aparecerá un formulario similar al de creación de grupo, en el
que podrá editar el título del grupo de chequeos.

\subsubsection{Borrado de grupo de chequeos}

Para borrar un grupo de chequeos, pulse el botón \textbf{Delete} en el menú de
operaciones del grupo. Aparecerá un formulario como el de la figura
\ref{fig:formulario-borrado-grupo}, pidiéndole confirmación del borrado. 

\begin{figure}[hbtp]
  \centering
  \includegraphics[width=0.5\textwidth]{apendice_manual_usuario/formulario_borrado_grupo}
  \caption{Formulario de confirmación de borrado}
  \label{fig:formulario-borrado-grupo}
\end{figure}

Si pulsa en el botón \textbf{Go back}, volverá al panel principal, cancelando el
borrado. Si pulsa el botón \textbf{Delete}, confirmará la operación y el sistema
borrará el grupo de chequeos, así como todos los chequeos que contenga.

\subsubsection{Activación y desactivación de grupo de chequeos}

Es posible activar y desactivar de forma grupal todos los chequeos de un
grupo. Para ello, deberá dirigirse al menú de operaciones del grupo y pulsar en
\textbf{Enable all checks} o en \textbf{Disable all checks}, según sea su decisión.

\subsection{Creación de chequeo}

\begin{figure}[hbtp]
  \centering
  \includegraphics[width=0.8\textwidth]{apendice_manual_usuario/pantalla_seleccion_tipo.png}
  \caption{Pantalla de selección del tipo de chequeo a crear}
  \label{fig:seleccionar-tipo-check}
\end{figure}

Como se ha comentado previamente, el elemento principal de SiteUp es el
\textbf{chequeo}. El usuario da de alta chequeos dentro de un grupo, y estos
chequeos se lanzan periódicamente, generando datos.

Para crear un chequeo dentro de un grupo deberá ir al menú de operaciones del
grupo, que aparece en la figura \ref{fig:operaciones-grupo} y se sitúa a la
derecha del título de un grupo de chequeos. Dentro del menú, deberá pulsar la
opción \textbf{Add check}.



Tras esto aparecerá una pantalla en la que deberá seleccionar el tipo de chequeo
que desea crear. La pantalla es similar a la de la figura
\ref{fig:seleccionar-tipo-check}. Como se explicó en la introducción de este
manual, es posible crear cuatro tipo de chequeos: chequeos de ping, chequeos de
puerto, chequeos de registros DNS y chequeos por HTTP.

Si pulsa en el botón del tipo que desea crear, aparecerá un formulario con
campos para rellenar.

\subsubsection{Creación de un chequeo por ping}

\begin{figure}[hbtp]
  \centering
  \includegraphics[width=0.8\textwidth]{apendice_manual_usuario/pantalla_crear_ping.png}
  \caption{Formulario de creación de chequeo por ping}
  \label{fig:crear-ping}
\end{figure}

El formulario para crear un chequeo de tipo ping aparece en la
figura~\ref{fig:crear-ping}. Los campos que deberá rellenar son los siguientes:

\begin{itemize}
\item \textbf{Target}: es el objetivo a chequear. Debe ser un \textit{hostname} o una IP.
\item \textbf{Título} del chequeo.
\item \textbf{Descripción} opcional del chequeo.
\item \textbf{Check interval}: intervalo entre chequeos, definido en minutos. El
  valor mínimo es de 1 minuto.
\item \textbf{Notify email}: indica si los cambios de estado en este chequeo
  deberán ser notificados por correo electrónico.
\item \textbf{Notify Android}: indica si los cambios de estado en este chequeo
  deberán ser notificados mediante la aplicación Android.
\item \textbf{Should check timeout}: indica si hay que verificar que el ping
  enviado vuelve dentro de un tiempo máximo determinado.
\item \textbf{Timeout value}: si la anterior opción está activa, este campo
  define el número de milisegundos que puede tardar en responder el servidor al
  ping, como máximo.
\end{itemize}

\subsubsection{Creación de un chequeo de puerto}

\begin{figure}[hbtp]
  \centering
  \includegraphics[width=0.8\textwidth]{apendice_manual_usuario/pantalla_crear_port.png}
  \caption{Formulario de creación de chequeo de puerto}
  \label{fig:crear-port}
\end{figure}

El formulario para crear un chequeo de puertos aparece en la
figura~\ref{fig:crear-port}. Los campos que deberá rellenar son los siguientes:

\begin{itemize}
\item \textbf{Título} del chequeo.
\item \textbf{Descripción} opcional del chequeo.
\item \textbf{Target}: es el objetivo a chequear. Debe ser un \textit{hostname} o una IP.
\item \textbf{Target port}: es el puerto a chequear. Debe ser un entero.
\item \textbf{Check for string}: opcionalmente, se puede indicar una cadena que
  debe estar presente en la respuesta al conectar con el puerto.
\item \textbf{Check interval}: intervalo entre chequeos, definido en minutos. El
  valor mínimo es de 1 minuto.
\item \textbf{Notify email}: indica si los cambios de estado en este chequeo
  deberán ser notificados por correo electrónico.
\item \textbf{Notify Android}: indica si los cambios de estado en este chequeo
  deberán ser notificados mediante la aplicación Android.
\end{itemize}

\subsubsection{Creación de un chequeo de DNS}

\begin{figure}[hbtp]
  \centering
  \includegraphics[width=0.8\textwidth]{apendice_manual_usuario/pantalla_crear_dns.png}
  \caption{Formulario de creación de chequeo de registro DNS}
  \label{fig:crear-dns}
\end{figure}

El formulario para crear un chequeo de registro DNS aparece en la figura
\ref{fig:crear-dns}. Los campos que deberá rellenar son los siguientes:

\begin{itemize}
\item \textbf{Título} del chequeo.
\item \textbf{Descripción} opcional del chequeo.
\item \textbf{Target}: es el objetivo a chequear. Debe ser un \textit{hostname}.
\item \textbf{Record type}: indica el tipo de registro a chequear. Puede ser A, AAAA, CNAME, TXT y MX.
\item \textbf{Resolved address}: indica el contenido que debe tener el registro.
\item \textbf{Check interval}: intervalo entre chequeos, definido en minutos. El
  valor mínimo es de 1 minuto.
\item \textbf{Notify email}: indica si los cambios de estado en este chequeo
  deberán ser notificados por correo electrónico.
\item \textbf{Notify Android}: indica si los cambios de estado en este chequeo
  deberán ser notificados mediante la aplicación Android.
\end{itemize}

\subsubsection{Creación de un chequeo por HTTP}

\begin{figure}[hbtp]
  \centering
  \includegraphics[width=0.8\textwidth]{apendice_manual_usuario/pantalla_crear_http.png}
  \caption{Formulario de creación de chequeo por HTTP}
  \label{fig:crear-http}
\end{figure}

El formulario para crear un chequeo por HTTP aparece en la figura
\ref{fig:crear-http}. Los campos que deberá rellenar son los siguientes:

\begin{itemize}
\item \textbf{Título} del chequeo.
\item \textbf{Descripción} opcional del chequeo.
\item \textbf{Target}: es el objetivo a chequear. Debe ser una URL.
\item \textbf{Status code}: es el código de estado que deberá devolver la petición..
\item \textbf{Check for string}: opcionalmente, se puede indicar una cadena que
  debe estar presente en la respuesta a la petición.
\item \textbf{Check interval}: intervalo entre chequeos, definido en minutos. El
  valor mínimo es de 1 minuto.
\item \textbf{Notify email}: indica si los cambios de estado en este chequeo
  deberán ser notificados por correo electrónico.
\item \textbf{Notify Android}: indica si los cambios de estado en este chequeo
  deberán ser notificados mediante la aplicación Android.
\end{itemize}


\subsection{Operaciones con chequeos}

\begin{figure}[hbtp]
  \centering
  \includegraphics[width=\textwidth]{apendice_manual_usuario/chequeo_valido.png}
  \caption{}
  \label{fig:chequeo}
\end{figure}

Una vez creado un chequeo, aparecerá en el panel de control con el aspecto de la
figura~\ref{fig:chequeo}. Como se aprecia en la figura, se muestra mediante un
icono el estado del chequeo, los datos (título, tipo, descripción) y una gráfica
que ilustra el estado del chequeo en las últimas 24 horas.

\begin{figure}[hbtp]
  \centering
  \includegraphics[width=\textwidth]{apendice_manual_usuario/chequeo_invalido.png}
  \caption{}
  \label{fig:chequeo-invalido}
\end{figure}

Si un chequeo detecta un error, cambia de apariencia, pasando a tener el aspecto
que ofrece la figura \ref{fig:chequeo-invalido}. Además, según se haya
configurado en las opciones del chequeo, se envía una notificación al usuario
informando del cambio de estado del chequeo.

Cuando se ve en la lista general de chequeos, cada chequeo tiene una serie de
botones que permiten llevar a cabo operaciones diversas, que se detallan a
continuación.

\subsubsection{Activación y desactivación de un chequeo}

Para desactivar un chequeo, tiene que pulsar el botón \textbf{Disable} de la
botonera, que solo aparecerá si el chequeo está activado. De igual modo, si el
chequeo está desactivado, el botón que aparecerá tendrá la etiqueta
\textbf{Enable}, que podrá pulsar para activarlo.

\subsubsection{Borrado de un chequeo}

En la botonera se sitúa el botón \textbf{Delete}, que podrá pulsar para borrar
el chequeo. Aparecerá un formulario pidiéndole confirmación. Si confirma la
acción, se procederá al borrado del chequeo y todos los datos dependientes de
éste.

\subsubsection{Edición de un chequeo}

Si desea editar un chequeo debe pulsar el botón \textbf{Edit}, situado en la
botonera del chequeo. Aparecerá un formulario similar al de creación de chequeo,
pero con los datos actuales ya introducidos. Puede modificar los datos que desee
y confirmar los cambios pulsando el botón \textbf{Update check} del formulario.

\subsubsection{Vista de detalle del chequeo}

\begin{figure}[hbtp]
  \centering
  \includegraphics[width=\textwidth]{apendice_manual_usuario/pantalla_detalle_chequeo.png}
  \caption{}
  \label{fig:detalle-chequeo}
\end{figure}

En la botonera del chequeo, en el panel de control, verá un botón
\textbf{Details}. Si lo pulsa será llevado a la pantalla de vista de detalle del
chequeo, que se ilustra en la figura \ref{fig:detalle-chequeo}. Esta pantalla
ofrece mucha información sobre el chequeo. En particular:

\begin{itemize}
\item En la parte superior se sitúa la información general:
  \begin{itemize}
  \item Icono con el estado del chequeo.
  \item Título.
  \item Tipo de chequeo.
  \item Descripción.
  \item Botonera de acciones.
  \end{itemize}

\item En la parte media de la pantalla se sitúa un bloque con pestañas. Mediante
  estas pestañas es posible indicar sobre qué periodo de tiempo queremos
  visualizar los datos: últimas 24 horas o última semana. La información
  disponible es la siguiente:

  \begin{itemize}
  \item Porcentaje de \textit{uptime} del chequeo.
  \item \textit{(Solo chequeos tipo ping)} Tiempo de respuesta mínimo, máximo y promedio.
  \item Gráfica con el estado del chequeo.
  \end{itemize}

\item En la parte inferior de la pantalla se sitúa el panel de eventos del
  chequeo. Éste permite ver los últimos cambios de estado del chequeo:

  \begin{itemize}
  \item Tipo de cambio de estado (\textit{Up}, \textit{Down}).
  \item Duración del estado.
  \item Hora de comienzo del estado.
  \item Hora de fin del estado. Si es el estado en el que el chequeo se
    encuentra actualmente, no aparecerá ninguna fecha.
  \item Extra: muestra información adicional para estados de error.
  \end{itemize}
\end{itemize}

\begin{figure}[hbtp]
  \centering
  \includegraphics[width=\textwidth]{apendice_manual_usuario/pantalla_grafica_semanal.png}
  \caption{Detalle de la gráfica para un intervalo de tiempo de una semana}
  \label{fig:grafica-semanal}
\end{figure} 

\subsection{Recepción de notificaciones por correo electrónico}

\subsubsection{Notificaciones de cambio de estado}

Cuando el sistema detecte un cambio de estado en un chequeo, y esté activada la
notificación por correo electrónico, se enviará un email como el que aparece en
la figura \ref{fig:email-estado}.

\begin{figure}[hbtp]
  \centering
  \includegraphics[width=0.7\textwidth]{apendice_manual_usuario/email_estado.png}
  \caption{Email de notificación de cambio de estado}
  \label{fig:email-estado}
\end{figure}

En ese email se le indicará el título del chequeo, el nuevo estado, la fecha y
hora en la que ha pasado al nuevo estado, y un enlace directo al panel de
detalle del chequeo en SiteUp, de forma que pueda obtener más información.

\subsubsection{Notificaciones de resumen diarias}

Si ha activado el envío de resúmenes diarios sobre el estado de sus chequeos, el
sistema enviará todos los días un correo electrónico con información. El aspecto
del correo que recibirá se puede ver en la figura \ref{fig:email-resumen}

\begin{figure}[hbtp]
  \centering
  \includegraphics[width=0.7\textwidth]{apendice_manual_usuario/email_resumen.png}
  \caption{Email diario de resumen}
  \label{fig:email-resumen}
\end{figure}


\section{Uso de la aplicación móvil}

La aplicación móvil \textbf{SiteUp Client} es una aplicación para para el
sistema operativo Android que permite a sus usuarios revisar los chequeos que
tiene dados de alta en la plataforma web SiteUp, así como recibir de manera
instantánea notificaciones \textit{push} sobre el estado de sus chequeos.

\subsection{Instalación}

La aplicación móvil se instala copiando y ejecutando el fichero APK en el
dispositivo. En un futuro se podrá encontrar en la \textit{Play Store} de
Google.

\subsection{Lanzamiento}

En la figura~\ref{fig:manual-android-launcher} se puede observar el aspecto de
la aplicación \textbf{SiteUp Client} en el lanzador de Android. Para ejecutarla,
simplemente pulse sobre el icono de la aplicación.

\begin{figure}[htbp]
  \centering
  \includegraphics[width=0.4\textwidth]{5_diseno/android-1}
  \caption{Icono de SiteUp en el lanzador}
  \label{fig:manual-android-launcher}
\end{figure}

\subsection{Carga de la aplicación}

Tras lanzar la aplicación, el sistema cargará los datos necesarios y se
comunicará con los servidores de Google para obtener los datos necesarios. La
apariencia de la aplicación durante este proceso se puede ver en la figura
\ref{fig:manual-android-login}.

\begin{figure}[htbp]
  \centering
  \includegraphics[width=0.4\textwidth]{5_diseno/android-2}
  \caption{Pantalla de carga de la aplicación móvil}
  \label{fig:manual-android-loading}
\end{figure}

\subsection{Inicio de sesión}

Tras la carga, la primera vez que lance la aplicación será necesario que inicie
sesión con las credenciales que utilizó en la plataforma web. Para ello, tendrá
que introducir su nombre de usuario y contraseña en el formulario que aparece en
la figura~\ref{fig:manual-android-loading}, y pulsar el botón \textbf{Sign in}.

\begin{figure}[htbp]
  \centering
  \includegraphics[width=0.4\textwidth]{5_diseno/android-3}
  \caption{Pantalla de inicio de sesión de la aplicación móvil}
  \label{fig:manual-android-login}
\end{figure}

La aplicación móvil se comunicará con los servidores centrales de SiteUp,
validará sus credenciales de usuario y, si son correctas, le dará acceso al
sistema. Si las credenciales de usuario no son correctas, el formulario mostrará
un error.

\subsection{Revisión de chequeos}

Si ha iniciado sesión correctamente, tendrá la opción de revisar la lista de
chequeos que tiene dados de alta en la plataforma web \textbf{SiteUp}. La
apariencia de la lista de chequeos es igual a la de la
figura~\ref{fig:manual-android-dashboard}.

\begin{figure}[hbtp]
  \centering
  \includegraphics[width=0.4\textwidth]{5_diseno/android-4}
  \caption{Pantalla de chequeos de la aplicación móvil}
  \label{fig:manual-android-dashboard}
\end{figure}

Si pulsa en cualquiera de los chequeos, será redirigido a la plataforma web, que
podrá navegar sin problema gracias a que se encuentra preparada para su uso en
dispositivos móviles.

\subsection{Recepción de notificaciones}

La utilidad principal de la aplicación es la posibilidad de recibir
notificaciones en cualquier momento provenientes del servidor -- conocidas como
notificaciones \textit{push}. 

El procedimiento es transparente al usuario. Simplemente lance la aplicación e
inice sesión. Tras hacer esto, la aplicación estará preparada para recibir
notificaciones desde el servidor.

Cuando se reciba una notificación, la aplicación mostará un mensaje en el área
de notificaciones del móvil, con información resumida. Puede ver el aspecto de
la notificación en la figura \ref{fig:manual-android-notification}. Si pulsa en
la notificación será llevado a los detalles del chequeo para obtener más
detalles sobre el origen del aviso.

\begin{figure}[hbtp]
  \centering
  \includegraphics[width=0.4\textwidth]{5_diseno/android-5}
  \caption{Aspecto de las notificaciones}
  \label{fig:manual-android-notification}
\end{figure}



%%% Local Variables:
%%% mode: latex
%%% TeX-master: "../memoria"
%%% End:


\chapter{Manual de usuario}
\label{chap:manual_usuario}
En este capítulo se explicarán en detalle las instrucciones de uso de la
plataforma web.

\section{Introducción a SiteUp}

\textbf{SiteUp} es un servicio web creado para la monitorización de servicios de
Internet, especialmente pensado para dueños de sitios web y servidores que
deseen estar informados en caso de cualquier eventualidad.

El uso de internet en todos los ámbitos de la vida, desde las interacciones
sociales hasta la búsqueda de empleo, han dado lugar a la creación de servicios
accesibles a través de Internet y a nuevos modelos de negocio basados en
aquellos. Resulta por tanto imprescindible contar con sistemas para monitorizar
la disponibilidad de estos servicios y, en caso necesario, actuar de manera que
puedan subsanarse las causas de los problemas detectados.

SiteUp cubre precisamente estas necesidades, permitiendo a los usuarios:

\begin{itemize}
\item Crear chequeos mediante envío de paquetes ICMP (Ping), para revisar que un
  servidor está activo y conectado a la red.
\item Crear chequeos de registros DNS, de forma que el contenido de estos
  registros sea el correcto en todo momento.
\item Crear chequeos de puertos remotos, para vigilar que los servicios
  concretos dentro de un servidor son accesibles.
\item Crear chequeos mediante peticiones HTTP, para controlar que un servicio
  web responde de forma correcta y con el contenido adecuado.
\item Establecer las opciones de periodicidad de estos chequeos, de acuerdo a
  las necesidades particulares de cada caso.
\item Recibir notificaciones tanto a través de correo electrónico como mediante
  una aplicación móvil para el sistema operativo Android.
\end{itemize}

\section{Requisitos previos}

Para el uso correcto de la plataforma web SiteUp, el principal requisito es
contar con un dispositivo con acceso a internet que disponga de un navegador
reciente, como por ejemplo Google Chrome. Cualquier tipo de tablet, teléfono
móvil u ordenador personal podrá acceder a la plataforma web SiteUp, pero para
aprovechar mejor el potencial del sistema se recomienda usar una computadora
personal con las siguientes características:

\begin{itemize}
\item Navegador reciente. Se recomienda el uso de Google Chrome.
\item Pantalla con resolución mínima de 1280 píxeles en su lado longitudinal.
\item Conexión a internet con un caudal de descarga de al menos 512kb/s.
\end{itemize}

Si desea hacer uso de la aplicación móvil, será necesario que el usuario cuente
con un dispositivo con las siguientes características mínimas:

\begin{itemize}
\item Sistema operativo Android, versión mínima 2.3, versión recomendada 4.3.
\item Pantalla de al menos 3 pulgadas de diagonal y resolución mínima de 320x480
  píxeles.
\end{itemize}

\section{Uso de la plataforma web}

\subsection{Acceso}

A fecha de abril de 2014, la dirección de acceso a la plataforma web de SiteUp
es \url{http://siteup.josetomastocino.com}. Simplemente introduzca esa dirección
en su navegador y automáticamente será redirigido a la plataforma web, donde
verá la pantalla de bienvenida.

\subsection{Registro de usuario}

\begin{figure}[hbtp]
  \centering
  \includegraphics[width=0.3\textwidth]{apendice_manual_usuario/general_botones.png}
  \caption{Botones principales para el usuario anónimo}
  \label{fig:botones-1}
\end{figure}

Si es la primera vez que accede al sitio web deberá registrarse. Para ello, debe
pulsar en el botón \textbf{Sign up}, que se ve en la figura \ref{fig:botones-1}.

Tras ello, aparecerá el formulario de registro de usuario, que se ve en la
figura \ref{fig:registro}

\begin{figure}[hbtp]
  \centering
  \includegraphics[width=0.5\textwidth]{apendice_manual_usuario/formulario_registro.png}
  \caption{Formulario de registro de usuario}
  \label{fig:registro}
\end{figure}

En el formulario deberá introducir los siguientes datos:

\begin{itemize}
\item \textbf{Nombre de usuario}. Debe tener 30 caracteres como máximo, y solo
  pueden usarse letras, números y guiones.
\item \textbf{Dirección de correo electrónico}. Asegúrese de que es una
  dirección de correo correcta.
\item \textbf{Contraseña}.
\end{itemize}

Tras introducir los datos, deberá pulsar el botón \textbf{Sign up} del
formulario. Si el registro es correcto, el sistema dará de alta su usuario y
mostrará la pantalla \ref{fig:registro-correcto}. Si no, volverá a aparecer el
formulario con información sobre los errores en los datos.

\begin{figure}[hbtp]
  \centering
  \includegraphics[width=\textwidth]{apendice_manual_usuario/pantalla_register_completo.png}
  \caption{Mensaje indicando el registro correcto}
  \label{fig:registro-correcto}
\end{figure}

\subsection{Inicio de sesión}

Una vez que haya registrado su usuario, deberá pulsar en el botón \textbf{Log
  in}, como aparece en la imagen \ref{fig:botones-1}. Esto le llevará al
formulario de inicio de sesión que aparece en la figura \ref{fig:login}

\begin{figure}[hbtp]
  \centering
  \includegraphics[width=0.5\textwidth]{apendice_manual_usuario/general_formulario_login.png}
  \caption{Formulario de inicio de sesión}
  \label{fig:login}
\end{figure}

Deberá introducir su nombre de usuario y la contraseña que escribió al
registrarse, y posteriormente pulsar en el botón \textbf{Log in} del
formulario. Si los datos son correctos, será llevado al panel principal de la
web, donde podrá gestionar sus chequeos. Si no son correctos, volverá a aparecer
el formulario, con información sobre los errores encontrados.

\subsection{Recuperación de contraseña}

Si se registró como usuario pero no recuerda su contraseña es posible
recuperarla.

Para ello, vaya al formulario de inicio de sesión como ser ha explicado en el
paso anterior. En ese formulario, en la zona inferior, verá un enlace titulado
\textbf{Forgot password?} que deberá pulsar. Después de hacerlo se mostrará el
formulario de recuperación de contraseña, ilustrado en la figura
\ref{fig:formulario-reset}.

\begin{figure}[hbtp]
  \centering
  \includegraphics[width=0.5\textwidth]{apendice_manual_usuario/general_formulario_reset.png}
  \caption{Formulario de recuperación de contraseña}
  \label{fig:formulario-reset}
\end{figure}

En este formulario deberá introducir la dirección de correo electrónico que
utilizó al registrarse, y pulsar el botón \textbf{Send reset code}. Si lo hace
correctamente, aparecerá la pantalla de confirmación que se ve en la figura
\ref{fig:pantalla-reset-enviado}.

\begin{figure}[hbtp]
  \centering
  \includegraphics[width=0.5\textwidth]{apendice_manual_usuario/pantalla_reset_enviado.png}
  \caption{Pantalla de confirmación de envío de código de reinicio}
  \label{fig:pantalla-reset-enviado}
\end{figure}

El sistema habrá enviado un enlace a su dirección de correo electrónico. Deberá
pulsar el enlace que aparece en ese correo. Asegúrese de que recibe
correctamente el correo. Si no lo recibe, revise la carpeta de correo no
deseado. El enlace del correo le llevará a un formulario de reseteo de
contraseña como el que aparece en la figura \ref{fig:formulario-reset-2}.

\begin{figure}[hbtp]
  \centering
  \includegraphics[width=0.5\textwidth]{apendice_manual_usuario/general_formulario_reset_2.png}
  \caption{}
  \label{fig:formulario-reset-2}
\end{figure}

En este formulario deberá escribir la nueva contraseña de su cuenta dos veces,
para confirmar que la escribe correctamente. Tras ello, deberá pulsar el botón
\textbf{Set new password}, lo cual confirmará los cambios. 

Tras este proceso, podrá iniciar sesión como se ha descrito previamente, esta
vez utilizando la nueva contraseña.


\subsection{Pantalla de listado de chequeos}

\begin{figure}[hbtp]
  \centering
  \includegraphics[width=0.8\textwidth]{apendice_manual_usuario/pantalla-dashboard.png}
  \caption{Panel de control}
  \label{fig:dashboard}
\end{figure}

\begin{figure}[hbtp]
  \centering
  \includegraphics[width=0.4\textwidth]{apendice_manual_usuario/general_botones_2.png}
  \caption{Botones principales para el usuario logueado}
  \label{fig:botones-2}
\end{figure}

En la figura \ref{fig:dashboard} se puede ver el aspecto del panel de control de
la plataforma web SiteUp. Desde este panel podrá hacer una gran cantidad de
acciones, desde crear grupos de chequeos hasta editar sus detalles de usuario.

Es posible acceder a esta pantalla desde cualquier sitio de la web si está
logueado, pulsando en el botón \textbf{Dashboard} de la cabecera de la web, tal
y como se ve en la figura \ref{fig:botones-2}.

\subsection{Cerrar sesión}

Si ha iniciado sesión y desea cerrarla, debe dirigirse a los botones
principales, situados en el lado derecho de la cabecera de la web, tal y como
aparecen en la figura \ref{fig:botones-2}. Deberá pulsar el botón \textbf{Log
  out}, tras lo cual se cerrara su sesión y será redirigido a la pantalla
inicial de SiteUp.

\subsection{Edición de datos de usuario}

Si desea editar alguno de los datos de usuario que ingresó cuando hizo el
registro, deberá primero pulsar el botón \textbf{Your profile}, situado en la
cabecera de la web. Es el segundo de los botones que aparecen, tal y como se ve
en la figura \ref{fig:botones-2}.

Una vez hecho esto, aparecerá el formulario de edición de usuario, ilustrado en
la figura \ref{fig:formulario-usuario}, en el que verá sus actuales datos de
usuario. 

\begin{figure}[hbtp]
  \centering
  \includegraphics[width=0.8\textwidth]{apendice_manual_usuario/pantalla_perfil_usuario.png}
  \caption{}
  \label{fig:formulario-usuario}
\end{figure}

En este formulario podrá editar:

\begin{itemize}
\item Su \textbf{nombre} de usuario.
\item Su \textbf{dirección} de correo electrónico.
\item Si desea o no recibir \textbf{resúmenes diarios} sobre el estado de sus chequeos.
\end{itemize}

Una vez editados los datos que desee deberá pulsar el botón \textbf{Update
  details}. También podrá editar su contraseña pulsando el botón \textbf{Change
  password}, situado al pie del formulario a la izquierda.


\subsection{Creación de grupo de chequeos}

En SiteUp, los chequeos se organizan en \textbf{grupos de chequeos}, lo que
permite trabajar con ellos más fácilmente a la hora de activarlos, desactivarlos
o borrarlos. 

Para crear un grupo de chequeos deberá pulsar el botón \textbf{Add group},
situado en la zona superior derecha del panel principal. El botón puede verse en
la figura \ref{fig:crear-grupo}.

\begin{figure}[hbtp]
  \centering
  \includegraphics[width=0.3\textwidth]{apendice_manual_usuario/grupo_boton_nuevo.png}
  \caption{Botón para crear nuevo grupo de chequeos}
  \label{fig:crear-grupo}
\end{figure}

Tras pulsar el botón aparecerá el formulario de creación de nuevo grupo. En él
deberá introducir los datos del nuevo grupo. En particular:

\begin{itemize}
\item El \textbf{título} que identifica al grupo.
\end{itemize}

Tras ello, deberá pulsar el botón \textbf{Create group}, y el grupo será creado

\begin{figure}[H]
  \centering
  \includegraphics[width=0.8\textwidth]{apendice_manual_usuario/grupo_formulario_crear.png}
  \caption{Formulario de creación de grupo de chequeos}
  \label{fig:}
\end{figure}


\subsection{Operaciones con grupos de chequeos}

En la figura~\ref{fig:operaciones-grupo} se puede ver el menú de operaciones
disponibles para un grupo de chequeos. Este menú aparece en el panel principal,
al lado derecho junto al título de cada grupo de chequeos.

\begin{figure}[hbtp]
  \centering
  \includegraphics[width=0.5\textwidth]{apendice_manual_usuario/grupo_submenu_acciones.png}
  \caption{}
  \label{fig:operaciones-grupo}
\end{figure}

\subsubsection{Edición de grupo de chequeos}

En el menú de operaciones del grupo de chequeos, pulse el botón
\textbf{Edit}. Aparecerá un formulario similar al de creación de grupo, en el
que podrá editar el título del grupo de chequeos.

\subsubsection{Borrado de grupo de chequeos}

Para borrar un grupo de chequeos, pulse el botón \textbf{Delete} en el menú de
operaciones del grupo. Aparecerá un formulario como el de la figura
\ref{fig:formulario-borrado-grupo}, pidiéndole confirmación del borrado. 

\begin{figure}[hbtp]
  \centering
  \includegraphics[width=0.5\textwidth]{apendice_manual_usuario/formulario_borrado_grupo}
  \caption{Formulario de confirmación de borrado}
  \label{fig:formulario-borrado-grupo}
\end{figure}

Si pulsa en el botón \textbf{Go back}, volverá al panel principal, cancelando el
borrado. Si pulsa el botón \textbf{Delete}, confirmará la operación y el sistema
borrará el grupo de chequeos, así como todos los chequeos que contenga.

\subsubsection{Activación y desactivación de grupo de chequeos}

Es posible activar y desactivar de forma grupal todos los chequeos de un
grupo. Para ello, deberá dirigirse al menú de operaciones del grupo y pulsar en
\textbf{Enable all checks} o en \textbf{Disable all checks}, según sea su decisión.

\subsection{Creación de chequeo}

\begin{figure}[hbtp]
  \centering
  \includegraphics[width=0.8\textwidth]{apendice_manual_usuario/pantalla_seleccion_tipo.png}
  \caption{Pantalla de selección del tipo de chequeo a crear}
  \label{fig:seleccionar-tipo-check}
\end{figure}

Como se ha comentado previamente, el elemento principal de SiteUp es el
\textbf{chequeo}. El usuario da de alta chequeos dentro de un grupo, y estos
chequeos se lanzan periódicamente, generando datos.

Para crear un chequeo dentro de un grupo deberá ir al menú de operaciones del
grupo, que aparece en la figura \ref{fig:operaciones-grupo} y se sitúa a la
derecha del título de un grupo de chequeos. Dentro del menú, deberá pulsar la
opción \textbf{Add check}.



Tras esto aparecerá una pantalla en la que deberá seleccionar el tipo de chequeo
que desea crear. La pantalla es similar a la de la figura
\ref{fig:seleccionar-tipo-check}. Como se explicó en la introducción de este
manual, es posible crear cuatro tipo de chequeos: chequeos de ping, chequeos de
puerto, chequeos de registros DNS y chequeos por HTTP.

Si pulsa en el botón del tipo que desea crear, aparecerá un formulario con
campos para rellenar.

\subsubsection{Creación de un chequeo por ping}

\begin{figure}[hbtp]
  \centering
  \includegraphics[width=0.8\textwidth]{apendice_manual_usuario/pantalla_crear_ping.png}
  \caption{Formulario de creación de chequeo por ping}
  \label{fig:crear-ping}
\end{figure}

El formulario para crear un chequeo de tipo ping aparece en la
figura~\ref{fig:crear-ping}. Los campos que deberá rellenar son los siguientes:

\begin{itemize}
\item \textbf{Target}: es el objetivo a chequear. Debe ser un \textit{hostname} o una IP.
\item \textbf{Título} del chequeo.
\item \textbf{Descripción} opcional del chequeo.
\item \textbf{Check interval}: intervalo entre chequeos, definido en minutos. El
  valor mínimo es de 1 minuto.
\item \textbf{Notify email}: indica si los cambios de estado en este chequeo
  deberán ser notificados por correo electrónico.
\item \textbf{Notify Android}: indica si los cambios de estado en este chequeo
  deberán ser notificados mediante la aplicación Android.
\item \textbf{Should check timeout}: indica si hay que verificar que el ping
  enviado vuelve dentro de un tiempo máximo determinado.
\item \textbf{Timeout value}: si la anterior opción está activa, este campo
  define el número de milisegundos que puede tardar en responder el servidor al
  ping, como máximo.
\end{itemize}

\subsubsection{Creación de un chequeo de puerto}

\begin{figure}[hbtp]
  \centering
  \includegraphics[width=0.8\textwidth]{apendice_manual_usuario/pantalla_crear_port.png}
  \caption{Formulario de creación de chequeo de puerto}
  \label{fig:crear-port}
\end{figure}

El formulario para crear un chequeo de puertos aparece en la
figura~\ref{fig:crear-port}. Los campos que deberá rellenar son los siguientes:

\begin{itemize}
\item \textbf{Título} del chequeo.
\item \textbf{Descripción} opcional del chequeo.
\item \textbf{Target}: es el objetivo a chequear. Debe ser un \textit{hostname} o una IP.
\item \textbf{Target port}: es el puerto a chequear. Debe ser un entero.
\item \textbf{Check for string}: opcionalmente, se puede indicar una cadena que
  debe estar presente en la respuesta al conectar con el puerto.
\item \textbf{Check interval}: intervalo entre chequeos, definido en minutos. El
  valor mínimo es de 1 minuto.
\item \textbf{Notify email}: indica si los cambios de estado en este chequeo
  deberán ser notificados por correo electrónico.
\item \textbf{Notify Android}: indica si los cambios de estado en este chequeo
  deberán ser notificados mediante la aplicación Android.
\end{itemize}

\subsubsection{Creación de un chequeo de DNS}

\begin{figure}[hbtp]
  \centering
  \includegraphics[width=0.8\textwidth]{apendice_manual_usuario/pantalla_crear_dns.png}
  \caption{Formulario de creación de chequeo de registro DNS}
  \label{fig:crear-dns}
\end{figure}

El formulario para crear un chequeo de registro DNS aparece en la figura
\ref{fig:crear-dns}. Los campos que deberá rellenar son los siguientes:

\begin{itemize}
\item \textbf{Título} del chequeo.
\item \textbf{Descripción} opcional del chequeo.
\item \textbf{Target}: es el objetivo a chequear. Debe ser un \textit{hostname}.
\item \textbf{Record type}: indica el tipo de registro a chequear. Puede ser A, AAAA, CNAME, TXT y MX.
\item \textbf{Resolved address}: indica el contenido que debe tener el registro.
\item \textbf{Check interval}: intervalo entre chequeos, definido en minutos. El
  valor mínimo es de 1 minuto.
\item \textbf{Notify email}: indica si los cambios de estado en este chequeo
  deberán ser notificados por correo electrónico.
\item \textbf{Notify Android}: indica si los cambios de estado en este chequeo
  deberán ser notificados mediante la aplicación Android.
\end{itemize}

\subsubsection{Creación de un chequeo por HTTP}

\begin{figure}[hbtp]
  \centering
  \includegraphics[width=0.8\textwidth]{apendice_manual_usuario/pantalla_crear_http.png}
  \caption{Formulario de creación de chequeo por HTTP}
  \label{fig:crear-http}
\end{figure}

El formulario para crear un chequeo por HTTP aparece en la figura
\ref{fig:crear-http}. Los campos que deberá rellenar son los siguientes:

\begin{itemize}
\item \textbf{Título} del chequeo.
\item \textbf{Descripción} opcional del chequeo.
\item \textbf{Target}: es el objetivo a chequear. Debe ser una URL.
\item \textbf{Status code}: es el código de estado que deberá devolver la petición..
\item \textbf{Check for string}: opcionalmente, se puede indicar una cadena que
  debe estar presente en la respuesta a la petición.
\item \textbf{Check interval}: intervalo entre chequeos, definido en minutos. El
  valor mínimo es de 1 minuto.
\item \textbf{Notify email}: indica si los cambios de estado en este chequeo
  deberán ser notificados por correo electrónico.
\item \textbf{Notify Android}: indica si los cambios de estado en este chequeo
  deberán ser notificados mediante la aplicación Android.
\end{itemize}


\subsection{Operaciones con chequeos}

\begin{figure}[hbtp]
  \centering
  \includegraphics[width=\textwidth]{apendice_manual_usuario/chequeo_valido.png}
  \caption{}
  \label{fig:chequeo}
\end{figure}

Una vez creado un chequeo, aparecerá en el panel de control con el aspecto de la
figura~\ref{fig:chequeo}. Como se aprecia en la figura, se muestra mediante un
icono el estado del chequeo, los datos (título, tipo, descripción) y una gráfica
que ilustra el estado del chequeo en las últimas 24 horas.

\begin{figure}[hbtp]
  \centering
  \includegraphics[width=\textwidth]{apendice_manual_usuario/chequeo_invalido.png}
  \caption{}
  \label{fig:chequeo-invalido}
\end{figure}

Si un chequeo detecta un error, cambia de apariencia, pasando a tener el aspecto
que ofrece la figura \ref{fig:chequeo-invalido}. Además, según se haya
configurado en las opciones del chequeo, se envía una notificación al usuario
informando del cambio de estado del chequeo.

Cuando se ve en la lista general de chequeos, cada chequeo tiene una serie de
botones que permiten llevar a cabo operaciones diversas, que se detallan a
continuación.

\subsubsection{Activación y desactivación de un chequeo}

Para desactivar un chequeo, tiene que pulsar el botón \textbf{Disable} de la
botonera, que solo aparecerá si el chequeo está activado. De igual modo, si el
chequeo está desactivado, el botón que aparecerá tendrá la etiqueta
\textbf{Enable}, que podrá pulsar para activarlo.

\subsubsection{Borrado de un chequeo}

En la botonera se sitúa el botón \textbf{Delete}, que podrá pulsar para borrar
el chequeo. Aparecerá un formulario pidiéndole confirmación. Si confirma la
acción, se procederá al borrado del chequeo y todos los datos dependientes de
éste.

\subsubsection{Edición de un chequeo}

Si desea editar un chequeo debe pulsar el botón \textbf{Edit}, situado en la
botonera del chequeo. Aparecerá un formulario similar al de creación de chequeo,
pero con los datos actuales ya introducidos. Puede modificar los datos que desee
y confirmar los cambios pulsando el botón \textbf{Update check} del formulario.

\subsubsection{Vista de detalle del chequeo}

\begin{figure}[hbtp]
  \centering
  \includegraphics[width=\textwidth]{apendice_manual_usuario/pantalla_detalle_chequeo.png}
  \caption{}
  \label{fig:detalle-chequeo}
\end{figure}

En la botonera del chequeo, en el panel de control, verá un botón
\textbf{Details}. Si lo pulsa será llevado a la pantalla de vista de detalle del
chequeo, que se ilustra en la figura \ref{fig:detalle-chequeo}. Esta pantalla
ofrece mucha información sobre el chequeo. En particular:

\begin{itemize}
\item En la parte superior se sitúa la información general:
  \begin{itemize}
  \item Icono con el estado del chequeo.
  \item Título.
  \item Tipo de chequeo.
  \item Descripción.
  \item Botonera de acciones.
  \end{itemize}

\item En la parte media de la pantalla se sitúa un bloque con pestañas. Mediante
  estas pestañas es posible indicar sobre qué periodo de tiempo queremos
  visualizar los datos: últimas 24 horas o última semana. La información
  disponible es la siguiente:

  \begin{itemize}
  \item Porcentaje de \textit{uptime} del chequeo.
  \item \textit{(Solo chequeos tipo ping)} Tiempo de respuesta mínimo, máximo y promedio.
  \item Gráfica con el estado del chequeo.
  \end{itemize}

\item En la parte inferior de la pantalla se sitúa el panel de eventos del
  chequeo. Éste permite ver los últimos cambios de estado del chequeo:

  \begin{itemize}
  \item Tipo de cambio de estado (\textit{Up}, \textit{Down}).
  \item Duración del estado.
  \item Hora de comienzo del estado.
  \item Hora de fin del estado. Si es el estado en el que el chequeo se
    encuentra actualmente, no aparecerá ninguna fecha.
  \item Extra: muestra información adicional para estados de error.
  \end{itemize}
\end{itemize}

\begin{figure}[hbtp]
  \centering
  \includegraphics[width=\textwidth]{apendice_manual_usuario/pantalla_grafica_semanal.png}
  \caption{Detalle de la gráfica para un intervalo de tiempo de una semana}
  \label{fig:grafica-semanal}
\end{figure} 

\subsection{Recepción de notificaciones por correo electrónico}

\subsubsection{Notificaciones de cambio de estado}

Cuando el sistema detecte un cambio de estado en un chequeo, y esté activada la
notificación por correo electrónico, se enviará un email como el que aparece en
la figura \ref{fig:email-estado}.

\begin{figure}[hbtp]
  \centering
  \includegraphics[width=0.7\textwidth]{apendice_manual_usuario/email_estado.png}
  \caption{Email de notificación de cambio de estado}
  \label{fig:email-estado}
\end{figure}

En ese email se le indicará el título del chequeo, el nuevo estado, la fecha y
hora en la que ha pasado al nuevo estado, y un enlace directo al panel de
detalle del chequeo en SiteUp, de forma que pueda obtener más información.

\subsubsection{Notificaciones de resumen diarias}

Si ha activado el envío de resúmenes diarios sobre el estado de sus chequeos, el
sistema enviará todos los días un correo electrónico con información. El aspecto
del correo que recibirá se puede ver en la figura \ref{fig:email-resumen}

\begin{figure}[hbtp]
  \centering
  \includegraphics[width=0.7\textwidth]{apendice_manual_usuario/email_resumen.png}
  \caption{Email diario de resumen}
  \label{fig:email-resumen}
\end{figure}


\section{Uso de la aplicación móvil}

La aplicación móvil \textbf{SiteUp Client} es una aplicación para para el
sistema operativo Android que permite a sus usuarios revisar los chequeos que
tiene dados de alta en la plataforma web SiteUp, así como recibir de manera
instantánea notificaciones \textit{push} sobre el estado de sus chequeos.

\subsection{Instalación}

La aplicación móvil se instala copiando y ejecutando el fichero APK en el
dispositivo. En un futuro se podrá encontrar en la \textit{Play Store} de
Google.

\subsection{Lanzamiento}

En la figura~\ref{fig:manual-android-launcher} se puede observar el aspecto de
la aplicación \textbf{SiteUp Client} en el lanzador de Android. Para ejecutarla,
simplemente pulse sobre el icono de la aplicación.

\begin{figure}[htbp]
  \centering
  \includegraphics[width=0.4\textwidth]{5_diseno/android-1}
  \caption{Icono de SiteUp en el lanzador}
  \label{fig:manual-android-launcher}
\end{figure}

\subsection{Carga de la aplicación}

Tras lanzar la aplicación, el sistema cargará los datos necesarios y se
comunicará con los servidores de Google para obtener los datos necesarios. La
apariencia de la aplicación durante este proceso se puede ver en la figura
\ref{fig:manual-android-login}.

\begin{figure}[htbp]
  \centering
  \includegraphics[width=0.4\textwidth]{5_diseno/android-2}
  \caption{Pantalla de carga de la aplicación móvil}
  \label{fig:manual-android-loading}
\end{figure}

\subsection{Inicio de sesión}

Tras la carga, la primera vez que lance la aplicación será necesario que inicie
sesión con las credenciales que utilizó en la plataforma web. Para ello, tendrá
que introducir su nombre de usuario y contraseña en el formulario que aparece en
la figura~\ref{fig:manual-android-loading}, y pulsar el botón \textbf{Sign in}.

\begin{figure}[htbp]
  \centering
  \includegraphics[width=0.4\textwidth]{5_diseno/android-3}
  \caption{Pantalla de inicio de sesión de la aplicación móvil}
  \label{fig:manual-android-login}
\end{figure}

La aplicación móvil se comunicará con los servidores centrales de SiteUp,
validará sus credenciales de usuario y, si son correctas, le dará acceso al
sistema. Si las credenciales de usuario no son correctas, el formulario mostrará
un error.

\subsection{Revisión de chequeos}

Si ha iniciado sesión correctamente, tendrá la opción de revisar la lista de
chequeos que tiene dados de alta en la plataforma web \textbf{SiteUp}. La
apariencia de la lista de chequeos es igual a la de la
figura~\ref{fig:manual-android-dashboard}.

\begin{figure}[hbtp]
  \centering
  \includegraphics[width=0.4\textwidth]{5_diseno/android-4}
  \caption{Pantalla de chequeos de la aplicación móvil}
  \label{fig:manual-android-dashboard}
\end{figure}

Si pulsa en cualquiera de los chequeos, será redirigido a la plataforma web, que
podrá navegar sin problema gracias a que se encuentra preparada para su uso en
dispositivos móviles.

\subsection{Recepción de notificaciones}

La utilidad principal de la aplicación es la posibilidad de recibir
notificaciones en cualquier momento provenientes del servidor -- conocidas como
notificaciones \textit{push}. 

El procedimiento es transparente al usuario. Simplemente lance la aplicación e
inice sesión. Tras hacer esto, la aplicación estará preparada para recibir
notificaciones desde el servidor.

Cuando se reciba una notificación, la aplicación mostará un mensaje en el área
de notificaciones del móvil, con información resumida. Puede ver el aspecto de
la notificación en la figura \ref{fig:manual-android-notification}. Si pulsa en
la notificación será llevado a los detalles del chequeo para obtener más
detalles sobre el origen del aviso.

\begin{figure}[hbtp]
  \centering
  \includegraphics[width=0.4\textwidth]{5_diseno/android-5}
  \caption{Aspecto de las notificaciones}
  \label{fig:manual-android-notification}
\end{figure}



%%% Local Variables:
%%% mode: latex
%%% TeX-master: "../memoria"
%%% End:


\chapter{GNU Free Documentation License}
\label{sec:fdl}
\input{0_misc/fdl.tex}

\bibliographystyle{paquetes/hispa-annote}
\bibliography{0_misc/bibliografia}

\end{document}

%%% Local Variables:
%%% mode: latex
%%% TeX-master: t
%%% End:
